\documentclass[a4paper,11pt,openany]{ltjsbook}
\usepackage{luatexja}
\usepackage{luatexja-fontspec}
\usepackage{amsmath,amssymb,amsfonts,amsthm}
\usepackage{geometry}
\geometry{top=24mm,bottom=24mm,left=24mm,right=24mm,footskip=10mm}
\usepackage{tcolorbox}
\tcbuselibrary{skins,breakable}
\usepackage{hyperref}
\usepackage{enumitem}
\usepackage{booktabs}
\usepackage{fancyhdr}

\definecolor{laplaceblue}{RGB}{20, 50, 90}
\definecolor{goldaccent}{RGB}{180, 130, 20}
\definecolor{softblue}{RGB}{240, 245, 252}
\definecolor{darkslate}{RGB}{30, 45, 60}
\definecolor{boxframe}{RGB}{40, 80, 130}

\hypersetup{
    colorlinks=true,
    linkcolor=laplaceblue,
    citecolor=goldaccent,
    urlcolor=laplaceblue
}

\newtcolorbox{mathreviewbox}[2][]{
    enhanced,
    breakable,
    colback=softblue,
    colframe=boxframe,
    fonttitle=\bfseries\gtfamily,
    title={#2},
    sharp corners,
    boxrule=0.8pt,
    left=10pt, right=10pt, top=8pt, bottom=8pt,
    #1
}

\pagestyle{fancy}
\fancyhf{}
\fancyhead[LE,RO]{\thepage}
\fancyhead[LO]{\nouppercase{\rightmark}}
\fancyhead[RE]{\nouppercase{\leftmark}}
\renewcommand{\headrulewidth}{0.4pt}

\title{\Huge\textbf{確率の哲学的試論}\\[4mm]
\huge\textbf{決定論的宇宙観と確率計算の全体系}\\[6mm]
\Large\textbf{—— 原著全18章 完全翻訳・数学校閲版 ——}\\[5mm]
\large 原著：ピエール＝シモン・ラプラス（Pierre-Simon Laplace, 1814）\\[2mm]
\normalsize Project Gutenberg eBook \#58881: \textit{A Philosophical Essay on Probabilities}}

\author{\textbf{ピエール＝シモン・ラプラス} 著\\[2mm]
\textbf{教育情報分析研究会} 訳・数理校閲}
\date{2026年10月 校訂第1版}

\begin{document}

\maketitle

\frontmatter

\chapter*{解題：『確率の哲学的試論』の現代的意義}
\addcontentsline{toc}{chapter}{解題：『確率の哲学的試論』の現代的意義}

ピエール＝シモン・ラプラス（Pierre-Simon, marquis de Laplace, 1749--1827）が1814年に発表した『確率の哲学的試論』（\textit{Essai philosophique sur les probabilités}）は、彼が著した金字塔的巨著『確率の解析的理論』（\textit{Théorie analytique des probabilités}, 1812）の序論として講義された内容を発展させた、数学史上屈指の古典的著作である。

本書は高度な微積分や差分方程式の展開を伴う『解析的理論』のエッセンスを、あらゆる学問的探求者、哲学者、市民に向けて言葉で説き明かした傑作であり、同時に近代科学における「決定論（いわゆるラプラスの悪魔）」の宣言書としても世界史にその名を刻んでいる。

本書の特徴は、コイン投げやサイコロといった偶然の遊戯から出発しながら、条件付き確率、原因の確率（ベイズの定理）、大数の法則、中心極限定理、最小二乗法、そして天文学、測地学、気象学、生命表、保険制度、裁判の陪審員判決、議会の投票理論、さらには人間の認知的錯覚（バイアス）に至るまで、確率論が人間社会と自然認識の全領域を照らす普遍的道具であることを実証した点にある。

本版は、Project Gutenberg（eBook \#58881）として公開されているF.~W.~トラスコットおよびF.~L.~エモリーによる英訳決定版を底本とし、現代的な数理モデル（\texttt{qwen2.5-7b-instruct} による忠実な翻訳および \texttt{qwen2.5-math-7b-instruct} による数理・LaTeX検証）の厳密な照合を経て、全18章の完全翻訳として組版・製本されたものである。

\tableofcontents

\mainmatter


\chapter{序論（Introduction）}
\label{ch:i}

\begin{center}
\large\textit{確率論の哲学的動機と人間精神の認識論的基礎}
\end{center}
\vspace{4mm}

この哲学的な論文は、1795年に国民会議の決議により数学者としてラグランジュと共に教授職に就き、通常学校へ招かれた際に講義した確率に関する講義を発展させたものです。最近では、同じ主題について『解析的確率理論』というタイトルで著作を出版しました。ここでは分析を使わずにこの理論の原理と一般的な結果を提示し、これらを人生に最も重要な問題に適用します。これらの問題は大抵は確率的な問題です。厳密には、ほとんどすべての知識が問題的であると言えるかもしれません。そして、数学その他の科学で我々が確実に知ることができる少数のものさえも、真理を確認する主要な手段である帰納法と類推は確率に基づいており、したがって人間の知識全体はこの論文で述べる理論と密接に関連しています。

ここでは興味を持って観察されるべき点として、永遠の原理である理性、正義、人道についても考慮する際には、常にこれらの原則に沿う有利な確率が付随していることに注意してください。これらを遵守することは大きな利点があり、それらから逸脱することは重大な不便をもたらします。その確率は宝くじの有利な確率のように、偶然の揺れの中で最終的に優勢を占めます。

この論文で述べた反省が哲学者たちに注目してもらい、彼らの心にふさわしい主題へと導かれることを願っています。

\newpage

\chapter{確率について（Concerning Probability）}
\label{ch:ii}

\begin{center}
\large\textit{決定論的宇宙観・ラプラスの悪魔と主観的確率の定義}
\end{center}
\vspace{4mm}

すべての出来事、その出来事がいかに微細なものであっても大きな自然法則から独立しているように見えていても、それは必然的にその法則の結果である。太陽の回転のように、そのような出来事と宇宙全体との間のつながりを知らない限り、これらの出来事は最終的な原因や偶然によって説明されるべきものと考えられてきた。それらが一定の規則性を持って繰り返されるか、あるいは秩序に無関心なように見えるかに関係なく、このような仮定された原因は知識の範囲が広がるにつれて徐々に消失し、正確な哲学によって完全に否定されることになる。正確な哲学には、それらは実際の原因に対する私たちの知識不足の表現であるとしか見えない。

現在の出来事は、あることが起こるには必ずその原因があるという明白な原理に基づいて、前段階の出来事が結びついている。この命題は「充足理由律」として知られ、中立的な行動に対しても適用される。自由意志は決定的な動機なしにこれらの行動を生み出すことができない。もし我々が二つの状況が完全に類似しているとし、その一つでは意志が活動的であり、もう一つでは無関心である場合、私たちはその選択が原因のない結果であると言える。レービンツはこれをエピクロス派の盲目な運命と呼ぶ。

この反対の見解は、意識の幻覚である。意志の選択が中立的な事柄において逃れられない理由から生じることを失う視点から、意志の選択は動機なしに決定されるという錯覚を持つ。

我々は、宇宙の現在の状態を過去の状態による結果とし、また将来の状態による原因と考えるべきである。一瞬でもそのような知識を持つ知性があれば、自然に宿るすべての力と、その構成体がそれぞれ占める位置について与えられたデータを分析的に処理できるのであれば、それは一つの公式で宇宙最大の天体の動きから最も軽量な原子の動きまでを把握することができるだろう。そして、そのような知性には、過去も未来も現在の目で見ることができるほど確定的なものが何も存在しないことになる。

人類の心は、天文学者が達成した完成度において、この知識を持つ知性の微かな概念を提供している。機械学と幾何学の発見、そして万有引力の理解により、過去と未来の世界系の状態を同一の解析的表現で把握することができた。そのような方法を他の知識の対象にも適用することで、現象を一般的な法則に結びつけ、特定の条件がもたらすべき結果を見越すことに成功した。

これらの真実探求における努力は、我々が先ほど述べたような無限に遠く離れた大きな知性へと常に導かれる傾向がある。しかし、それは人類にとって特有の傾向であり、動物から永遠に隔てられている。この傾向は、動物を超越させ、その進歩は国家や時代における真の栄光を構成する。

次章：確率について

かつて、遠くない時代においても、異常な雨や極端な乾燥、彗星の非常に長い尾、日食、極光、一般にすべての異常現象は天の怒りの兆候とされ、これらの影響を避けるために天を求めた祈祷が行われていた。誰一人として惑星や太陽の動きを止める祈りを捧げることはなかった：観測はそのような祈りの無意味さをすぐに明らかにした。

しかし、これらの現象が長期的な間隔で発生し消失するため、自然の秩序に反しているように思われたことから、天が地上の罪によって激怒したとされ、それらは「天の報復を告げるために作られた」と考えられた。1456年の彗星の長い尾は、すでにトルコ軍による東羅馬帝国の崩壊という急速な成功で混乱に陥ったヨーロッパ全体を恐怖に包んだ。しかし、この星はその後四回の軌道を巡り、我々には非常に異なる興味を惹きつけた。

間隔が開いた間に得られた宇宙系の法則に関する知識により、人間と宇宙との真の関係に対する無知から生じた恐怖は払拭された。そしてハレーは1531年、1607年、1682年に見られた彗星とこの彗星が同一であることを認定し、その次の出現を1758年末または1759年初頭に予告した。

学術界は、この出現が科学の偉大な発見の一つを確認するものであり、セネカが言及した那些星体から巨大な高さを落下する星の回転について述べたときの予言を実現させるために、これに対する期待を抱いていた。彼は「数世紀にわたる研究によって明らかになる日が来るでしょう。現在隠されていることが明確に示される時が来るでしょう。そして後世は、私たちが逃げてしまったほど明白な真実があったことに驚くことでしょう」と述べていた。

クラロワはその後、この彗星が木星と土星の作用によって経験した干渉を分析するために取り組み、膨大な計算を行った結果、その次の近日点通過は4月の初め頃になると予測し、これは観測によって実際的に確認された。

天文学者が彗星の運動において見せる規則性は、疑いなくすべての現象にも存在するものである。

単純な気体分子や蒸気の描く曲線は、惑星の軌道のように一定の方法で制御される。それらの違いは、私たちの知識不足によるものである。

確率は、部分的にはこの知識不足、部分的には私たちの知識に関連している。三つ以上の事象から一つが発生することを知っているが、どの一つが発生するかについては何の理由も与えられていない。このような不確定な状態では、これらの事象の発生を確実に予告することは不可能である。しかし、任意に選んだ一つの事象が起こらない確率は高い。なぜなら、いくつかの同様に可能な場合があり、その一つが排除される一方で、他の一つだけがそれを支持するからである。

確率論は、同じ種類のすべての事象を一定数の等可能性のある状況に帰着させ、つまり、私たちがその存在について同等に不確定であるような状況に帰着させ、そして求める確率に対応する有利な状況の数を決定することから成る。この数の比は、

7

この確率を測定し、それは従って単なる分数であり、その分子は有利な状況の数で、分母はすべての可能な状況の数である。

次の確率の概念は、有利な状況の数とすべての可能性のある状況の数が同じ比率で増加するとき、確率が一定であることを仮定しています。これを確認するために、二つの壺AとBを考えましょう。壺Aには4個の白球と2個の黒球があり、壺Bには2個の白球と1個の黒球があります。

まず、壺Aの2つの黒球を糸で結びつけ、その糸が一つの黒球が取り出される瞬間に切れることを想像します。これにより4つの白球は2つの類似したシステムになります。有利な状況すべてが壺から黒球を取り出すことに有利になると仮定すると、これらの機会は全て黒球を引き出すことにつながります。

次に、これらの糸が全く切れないと想像します。この場合、可能性のある状況の数と黒球を引くための有利な状況の数は変わりませんが、同時に2つの球が引き出されます。すると壺Aから黒球を引く確率は元のものと同じになります。

しかし、これにより壺Bの場合を考えることができます。この壺には3つの球があり、それらを常に結ばれた2つの球のシステムに置き換えることができます。

すべての状況が一事件に対して有利であるとき、確率は確定性へと変化し、その表現は1になります。この条件において、確定性と確率は比較可能ですが、真実が厳密に証明されると、まだわずかな誤りの源を感じている場合との心の状態には根本的な違いがあるかもしれません。

ある事象がただの確率的なものである場合、各人がそれについて持っているデータの違いは、同じ対象に対する異なる意見が生じる主要な原因の一つです。たとえば、三つの甕（うん）A, B, Cを考えましょう。そのうち一つには黒い球しか入っておらず、他の二つには白い球しか入っていないとします。Cから一粒の球を引く確率が求められます。もしどの甕が黒い球を入れているか全く分からず、A, B, Cのいずれであるかについて理由もなくCよりBやAの方がより可能性が高いとは思わない場合、これらの三つの仮説は等しく可能なように見えます。そして、黒い球を引くのは最初の仮説だけなので、その確率は三分の一になります。

もしAが白い球しか入れていないことが分かっているとします。このとき、不確定性はBとCのみに拡がりますので、Cから引いた球が黒い確率は二分の一になります。最後には、AとBが白い球しか入れていないことを確認された場合、その確率は確実性に変わります。

したがって、多数の聴衆に関連する出来事は、聴衆の知識の程度によって異なる信頼度を持つことになります。報告者が完全にその事象を信じている場合、またその立場や人格から大きな信用を与えられる場合、彼の声明は情報がない聴衆にとって、驚くべきものであっても通常の声明と同じ確率を持ちます。

同様の人物によって証明され、その論理は彼らに完全な信頼をもたらすだろう。しかし、もし他の同等に信頼できる人々が同じ事象を否認していることが知られれば、疑念が生じ、聡明な聴衆はその事象を信用しないし、事実の確証があるか自然界の不変の法則に関わらず、それを否認するだろう。

誤りの蔓延は、多数が最も重要な生活上の問題について最善に情報を持っていると判断し、その信頼を与えている人々の意見の影響によって引き起こされる。魔術と占星術は、これらの誤りの二つの大きな例である。幼少期から灌漁され、検討せずに受け入れられ、普遍的な信用に基づくこれらのはずみが長期間にわたって存続した後も、科学の進歩により、啓蒙された人々の意見によって、最終的にそれらは一般の人々の中からも消え去った。この力は、道徳的世界における最も豊かな資源であり、一国全体において全く異なる概念を確立し、維持するものである。私たちが他の場面で同じ権威を持って支持している概念とは異なり。私たちは自分たちの意見とそれらの相対的な可能性を批判的に検討し、公平に weighs すべきではないか。我々が不十分に教育されていると判断した人々を啓蒙するべきであるが、まず私たち自身の意見を慎重に検討し、偏見なくその確率を考慮すべきである。

意見の相違は、既知のデータによる影響をどのように決定するかに依存します。確率論は非常に繊細な考慮事項を扱うため、同じデータであっても二人の人が異なる結果を得るのも不思議ではありません、特に非常に複雑な問題ではなおさらです。現在、この理論の一般原則を見ていきましょう。

\newpage

\chapter{確率計算の一般原理（General Principles of the Calculus）}
\label{ch:iii}

\begin{center}
\large\textit{同等確からしさ・独立試行・条件付き確率・ベイズの逆確率（10の根本原理）}
\end{center}
\vspace{4mm}

【初等原理】これらの原理の第一は、確率そのものに対する定義であり、これは以前見てきたように、有利な状況の数とすべての可能性のある状況の数との比を表します。

【第2原理】しかし、それには各々の状況が等しく可能なことを仮定しています。もしもそうではない場合、まずそれぞれの状況の可能性を決定することから始めます。その正確な評価は確率論における最も精密な点の一つです。その後、確率は各有利な状況の可能性の合計となります。この原理を例によって説明しましょう。

大きな非常に薄いコインを投げることを考えましょう。このコインには、対向する二つの大きな面があり、これらを表と裏と呼びます。これらの面は完全に相似しています。二回投げて少なくとも一度表が出る確率を求めましょう。明らかに四つの等しく可能な状況が生じることが分かります：最初の投げで表、そして次の投げでも表；最初の投げで表、次の投げでは裏；最初の投げでは裏、次の投げでは表；最後に両方とも裏です。これらの最初の三つの状況は求められるイベントを有利としますから、この確率は | となります。したがって、二回投げて少なくとも一度表が出るという賭けは三対一です。

このゲームでは、私たちが数えることができる異なる場合は三つだけです。つまり、最初の投げで表が出た場合、二度目の投げは不要となります；最初の投げで裏が出て、二度目の投げで表が出る場合；最後に、最初と二度目の投げで裏が出る場合です。ダランベルトがこれら三つの場合を等可能性として考慮するなら、その確率は | になります。

しかし明らかに、最初の投げで表が出る確率は f であり、他の二つの場合の確率は J であり、最初のケースは単純な事象であり、二つの事象が組み合わさったものに対応します：最初と二度目の投げで表が出る場合、そして最初の投げで表、二度目の投げで裏が出る場合です。したがって、第二の原理に従って、最初の投げで表が出る確率 f を最初の投げで裏が出て二度目の投げで表が出る確率 J に加えると、求められる確率は f となり、これは二回投げる場合での見出されるものと一致します。

この「仮定」は、そのイベントに賭ける人にとっての機会を全く変更しません；それは単に様々なケースを等可能性のケースに簡略化する役割を果たすだけです。

【第3原理】—確率論の理論において最も重要な点であり、最も誤解を生む要素は、これらの確率が互いに組み合わさることによってどのように増加または減少するかの方法です。もし事象が互いに関連しないなら、その結合した存在の確率はそれぞれの確率の積となります。つまり、

確率論の計算。13

一サイコロを投げて1の目が出る確率は$\frac{1}{6}$であり、二つのサイコロを同時に投げて両方とも1の目が出る確率は$\frac{1}{36}$です。一方のサイコロの各面が他方のサイコロの六つの面と組み合わされるため、実際には三十六通りの等しい可能性があるうちの一つだけが二つの1の目を出す結果となります。一般に、同じ状況で繰り返し発生する単純な事象の確率はその単純な事象の確率を与えられた回数のべき乗として表されます。したがって、ある分数が1未満である連続したべき乗は無限に減少します。従って、非常に多くの小さな確率に依存する事象は極めて低い確率になることがあります。

たとえば、証人二十人が順次その証言を伝える場合を考えます。最初の証人が次の証人に、その証人が次の証人にと連続して証言を伝え、それぞれの証言の確率が$\frac{1}{9}$であるとします。すると、これらの証言から推測される事象の確率は$\frac{1}{36}$以下になります。

この確率の減少は、物体の光が複数枚のガラス板によって遮られるときの暗転に似ています。少量のガラス板でも、単一のガラス板で明瞭に見える物体を全く見えなくするのに十分な量となります。

歴史家たちは、多くの世代を通じて伝わる事象の確率が著しく低下することについて十分に注意していないように思われます。多くの歴史的事象が確実とされてきたとしても、このテストによって少なくとも疑問視されるべきであることが明らかになります。

純粋な数学的な科学においては最も遠くまで及ぶ推論もその確実性を共有します。

彼らから導かれる円周率。解析を物理学に適用した場合、その結果は事実や経験と同等の確固たる certainty を持つ。しかし、倫理学などの人文科学では、各推論が前段階から推測的に導き出されるだけであるため、これらの推論がいかに確かなものであっても、その数が増えるにつれて誤りの可能性は増大し、結果として非常に遠く離れた推論においては真実の確率よりも大きな誤りの確率を上回る。

【第4原理】—二つの事象が互いに関連している場合、複合事象の確率は最初の事象の確率と、その事象が発生したとき次に起こる第二の事象の確率の積である。前記の三つのurn A, B, Cの場合を考えてみよう。二つは白い球のみを含んでおり、一つは黒い球のみを含んでいる。 urn Cから白い球を引く確率は f であり、この理由は三つのurn中で二つがその色の球を含むためである。しかし、 urn Cから白い球が引かれた後では、唯一の黒い球を含むurnの不確定性はurn Aと Bにまで及ぶ。 urn Bから白い球を引く確率は $p$ であり、 f と $p$ の積、すなわち $p$ が同時にurn Bと Cから白い球を引く確率である。

この例から、過去の事象が将来の事象の確率に与える影響を見ることができる。 urn Bから白い球を引く確率は最初 f であったが、 urn Cから白い球が引かれた後では $p$ となる。同じurnから黒い球が引かれた場合、その確率は確定的になる。

この影響を定量的に定めるために、以下の原理を使用する。これは前記の原理の推論である。

【第5原理】—事前に生じた事件の確率と、その一つと別の予想される二つ目のイベントを含むイベントの確率を計算した場合、二つ目の確率を最初の確率で割ったものは、観察されたイベントから予想されるイベントの確率となります。

ここでは、一部の哲学者が過去が将来の確率に与える影響について提起した問題が提示されます。硬貨投げにおいて、表が裏よりも頻繁に現れたとします。これだけで、硬貨の構造の中に有利な秘密の原因があると考えるようになります。生活における恒久的な幸せは、適切性を示す証拠となり、それに基づいて幸運な人々を優先的に選ぶべきだと考えさせます。しかし、状況の不確実性によって絶対的な判断に至れない場合、例えば硬貨投げにおいて表が出た後毎回異なる硬貨を使用する場合、過去は未来に対して何らの光を投げかけませんし、そのような過去を考慮することは荒唐無稽となります。

【第6原理】—観察された事件が原因とされる各々の要因は、その要因が存在すると仮定したときにそのイベントが発生する確率と同じくらいの可能性で示されます。これらの各々の要因の存在確率は、その要因から生じるイベントの確率を分子とし、すべての要因に対する類似の確率の合計を分母とする分数となります；これらの様々な要因が事前に不等確率である場合、それは必要です。

各原因から生じる事象の確率ではなく、その確率をそれぞれの原因自体の可能性で乗じることに代えて考える。これはこの分野の確率解析の基本原理であり、事象から原因へと移行する方法である。

この原理は、私たちは規則的な事象を特定の原因に帰属させる理由を説明する。一部の哲学者は、これらの事象が他の事象よりも可能性が低いと考えており、例えばコイン投げで表が出る二十回連続という組み合わせは、表と裏が不規則な形で混ざった組み合わせよりも性質上容易ではないと考えたことがある。しかし、過去の事象が未来の事象の可能性に影響を与えるという考え方は全く容認されない。規則的な組み合わせは稀に起こるだけではなく、それが少ない数であるからである。対称的な事象をどこかで見ると、その原因を探そうとするのは、対称的な事象が他の事象よりも可能性が低いとは考えていないからであり、この事象が規則的な原因の結果であるべきか、あるいは偶然の結果であるべきかという仮定の中で、前者の方が後者よりより確実であると考えるためである。テーブルで文字列コンスタンティノープルを見つけると、これは偶然の結果ではないと判断するが、それはその組み合わせが他のものよりも可能性が低いからではなく、この単語が我々の言語に使用されていない場合でも特定の原因によるものと考えないからである。しかし、この単語が我々の言語で使われているなら、ある人がこれらの文字をこのような順序に並べたというよりは、その組み合わせが偶然によるものであるよりも遥かに確実であると判断する。

ここでは「Extraordinary」という言葉を定義する場所である。私たちは思い描く中で全ての可能な事象を様々な形で並べ替える。

クラスについて；そして、我々は非常に少ない数の要素を含むクラスを特別なものと見なす。例えば、コイン投げにおいて、100回連続して表が出るという出来事は、その組み合わせのほとんど無限にあることから我々に特別に思える。これらの組み合わせを規則的な系列と不規則な系列とに分類し、後者の数は圧倒的に多い。

また、百万個の中から一つだけ白い球が入っている urn から白い球を引くという出来事も、その他の球は黒であるという状況下では我々に特別な印象を与える。これは、二つの色に対して二つのクラスしか形成しないからである。しかし、百万個の数字の中から例えば 475813 を引くという出来事は通常の出来事と見做される。なぜなら、個々の数字を比較し、分類せずに並べ替えると、一つの数字が他の数字よりも早く現れる理由がないからである。

以上のことを考慮すると、より特別な出来事ほどその実現性を裏付ける強い証拠が必要であることが一般に結論付けられる。証言する者は欺く可能性もしくは自己欺瞞の可能性があるため、これらの二つの要因が出来事が実在しない場合よりももっと確率が高いということである。この点については、証言の確率について議論する際に特に明らかになる。

【第7原理】—将来の出来事の確率は、観察された出来事を基にした各原因の確率と、その原因が存在する場合に将来の出来事が起こる確率の積の合計である。以下の例を用いてこの原理を説明する。

二つの球、つまり白い球または黒い球が入っている壷を考えましょう。これらの球の一つを引き取り、それを元の壷に戻し、新しい引き取りを行います。最初の二回の引き取りで白い球が引き出されたと仮定します。第三次の引き取りで白い球が引き出される確率を求めます。

ここでは、以下の二つの仮説だけを考えることができます：一つは白い球で、もう一つは黒い球である、または両方とも白い球である。最初の仮説では、観察されたイベントの確率は 1/2 であり、それは第二の仮説では確定性（確実性）となります。したがって、これらの仮説を原因として考えると、第六原理においてそれぞれの仮説に対する確率は % と | となります。しかし、最初の仮説が起こる場合、第三次に白い球が引き出される確率は 1/2 であり、それは第二の仮説では確定性と等しいです。したがって、最後の確率をそれぞれの対応する仮説の確率で掛けてから足し合わせると、これらの積の合計、すなわち T9 が第三次に白い球が引き出される確率となります。

単一のイベントの確率が未知の場合、その確率はゼロから一つまでの任意の値と仮定することができます。観察されたイベントから導かれた各仮説の確率は、第六原理によれば、その仮説におけるイベントの確率の分子であり、すべての仮説に対する類似する確率の合計が分母となります。したがって、イベントの可能性が特定の限界内に含まれる確率は、これらの限界内の分数の合計となります。もし

19

各分数を、対応する仮説で決定された未来の事象の確率に乘じます。すべての仮説に対する積の合計は、第七原理により、観察された事象から導き出される未来の事象の確率となります。したがって、過去に何回も連続して起こった事象が再び起こる確率は、その回数を1増やした値を同じ回数に2足した値で割ったものと等しいことが分かります。

歴史の最も古い時代を5000年前、または182623日とする場合、太陽が24時間ごとに恒常的に昇るという間隔において常に昇っていると仮定すると、明日も昇る確率は1826214対1となります。しかし、現時点での何らかの理由でその進行を止めることがないことを認識する人にとっては、この数字は著しく大きいものとなります。

ブフォンは『政治的算術』において、前記の確率を異なる方法で計算しています。彼は、その確率が1から分数だけ異なり、その分数の分子は1であり、分母は経過した日数に2を累乗した値であると仮定しています。しかし、過去の事象と原因および未来の事象の確率との関連性を真実の方法で結びつけることは、この有名な著者には未知でした。

\newpage

\chapter{希望値について（Concerning Hope）}
\label{ch:iv}

\begin{center}
\large\textit{数学的期待値と道徳的期待値（モラル・ホープ / 効用理論）}
\end{center}
\vspace{4mm}

確率は、その存在に関心を持つ人々の希望や恐怖を決定します。希望という言葉には様々な意味があり、それは一般的に、ある特定の利益を得ることに対する期待を表現し、それがしかり確かなものであるとは限らない仮定においてです。この理論における利益は、望む金額とその獲得確率との積から得られます；これは、リスクを避けるために事件が起こる可能性に比例して分配されるべき部分の合計です。このような分配は、すべての不確定要素を取り除いた場合に唯一公正なものです；なぜなら、等しい確率は等しい利益を与えます。これを数学的希望と呼びます。

【第8原理】——複数の事件によって得られる利益は、それぞれの事件の確率をその発生に付随する利益で掛けてから合計したものです。

この原理をいくつかの例で適用しましょう。まず、

希望について。21

パウルが表と裏のゲームで最初の投げで表が出ると2フラン、二回目だけ出ると5フランを受け取る場合を考えましょう。最初の場合の確率$\frac{1}{2}$に2フランを掛けて、二回目の場合の確率$\frac{1}{2}$に5フランを掛けると、その積の合計、つまり2フラン半分（2.5フラン）がパウルの利益となります。これは彼がこの優位性を得るために先払いすべき金額であり、ゲームの公平性を維持するためには投げはその利益と同じくらい行われるべきです。

パウルが最初の投げで表が出ると2フラン、二つ目の投げで表が出ると5フランを獲得する場合、最初に投げたかどうかに関わらず、二つ目の投げで表が出る確率が$p$であるとき、2フランと5フランを$p$で掛けて得られる積の合計はパウルにとって3分半フランの利益を与え、これこそ彼の賭け金となる。

【第9原理】—一連の確率的な出来事において、その中の有利な出来事が利益をもたらし、不利な出来事が損失をもたらす場合、その結果得られる利益は、各有利な出来事の確率とその利益の積の合計から、各不利な出来事の確率とその損失の積の合計を引いたもので求めることができる。もし後者の合計が前者より大きいとき、利益は損失となり、希望は恐怖に変わる。

したがって、生活における行動においては、期待される利益の確率で掛けて得られる積が、類似の積（損失に対する）よりも少なくならないようにしなければならない。しかし、これを達成するためには、正確な利益、損失、そしてそれぞれの確率を評価することが必要である。これを行うために必要なものは、大きな心の正確さ、鋭い判断力、そして豊富な経験であり、自己愛に食いつく人々が抱く運や幸せという誤った概念から自分を守る能力も求められる。

前節で述べた原理を以下の問題に適用することは、幾何学者たちにとって大きな課題となりました。パウルは最初の投げで表が出れば2フラン、二回目の投げのみで表が出れば4フラン、三回目の投げのみで表が出れば8フラン、そしてそのようなパターンが続く場合、彼の賭け金は投げwhelmsの数に等しくなります。つまり、ゲームが無限に続けば、賭け金も無限になります。しかし、誰もがこのゲームでさえ少量でも、例えば5フランを賭けることを望まないでしょう。計算結果と一般的な感覚との間のこのような差異はどこから来ているのでしょうか？すぐにわかるように、それは以下の点に帰着します：利益が私たちに与える道徳的利益はその利益に対して比例しないこと、そしてそれが千種類もの状況に依存すること、これらはしばしば定義するのが難しいものですが、最も一般的で重要なものは運です。

実際には、百フランしか持っていない人にとって1フランは、億万長者にとっての価値よりも遥かに大きいことが明白です。したがって、期待される利益において絶対的な価値と相対的な価値を区別する必要があります。

相対的な価値は、その利益が望ましい理由となる動機によって規制されますが、絶対的な価値はそれらに依存しません。この相対的価値を評価するための一般的な原理は与えられませんが、ダニエル・ベルヌーイによって提案された次の原理は多くの場合に適用されます。

「望みに関する論考」23

【第10原理】無限小総額の相対的な価値は、その総額が関心を持つ人の総利益で割ったものと等しい。この原理は、誰しも一定の利益があることを仮定しており、その価値は決してゼロと評価されることはないという。実際、何も所有していない人でも、自分の労働の産物や希望に対して、絶対的に必要な生活を維持するための最低限の価値を与える。

この原理を分析に適用すると、以下の規則を得る：個々人の財産の一部（その期待とは無関係）を単位と定義しよう。この財産が期待とその確率によって異なる価値を持つ場合、それぞれの価値をその確率で示すべき指数に上げた積は、単位として取り扱った財産と期待から得られる道徳的な利益と同じ程度の物理的財産を与えるものとなる。この積から単位を引いた差が個々人の物理的財産の増加であり、これを道徳的希望という。容易に理解できるように、単位として取り扱った財産が無限大になる場合、この増加は数学的な期待と一致する。しかし、これらの変動が単位の一部である場合、二つの希望は著しく異なる可能性がある。

このルールは、普通の感覚に合致した結果を導き出し、その正確さを一定の程度まで把握できるようにする。例えば、パウルの運が二百フランである場合、彼は九フランを超えて賭けるべきではないことが前問から明らかになる。同じルールはまた、期待される利益を危険に分散することを我々に導き、その全体を危険にさらすよりも好ましいと示唆する。

同様の結果として、最も公正なゲームにおいても損失は常に利益より大きいことがわかる。例えば、百フランの財産を持つプレイヤーがコイン投げで五十フランを賭ける場合を考える。彼の賭け後の財産は八十七フランに減少し、つまりこの残額がプレイヤーに与える道徳的な利益は、賭け前の彼の財産状態と同じである。このゲームは、期待される金額がその確率と等しい賭け金でも不利であり続ける。

このような判断により、期待される金額がその確率と等しくないゲームの非道徳性を評価できる。これらのゲームは誤った論理と、それらによって刺激される貪欲さによってのみ存続する。人々は自己の必需品を幻滅的な希望のために犠牲にし、その不確実性を十分に評価できないまま、無限の悪事が生じる源となる。

賭博ゲームの欠点、期待される利益全体を同じ危険から守ることの利点、そして普通の感覚によって示唆されたすべて類似した結果は、個々人の物理的財産がその道徳的財産を表す関数の性質に関わらず存在する。重要なのは、この関数の増加に対する物理的財産の増加の割合が、後者が増大するにつれて減少することである。

数学的論理と語調に

\newpage

\chapter{確率計算の解析的方法（Analytical Methods）}
\label{ch:v}

\begin{center}
\large\textit{母関数理論・有限差分方程式・ラプラス変換と漸近積分展開}
\end{center}
\vspace{4mm}

我々が最近説明した原理を、確率に関する様々な問題に適用するには、その調査のために生じた分析法のいくつかが必要となり、特に組合せ論と有限差分の計算が生まれました。

各文字を1つずつ、2つずつ、3つずつなど、それぞれの組み合わせに対して1を係数とする形で取り出すために、最初の文字からn番目の文字までの各文字を含む二項式の積を形成し、その展開した積から1を引くと、結果はこれらの文字を1つずつ、2つずつ、3つずつなどそれぞれの組み合わせで表されます。この原理を利用して、これらnつの文字をs回ずつ取り出す組み合わせの数を得るためには、まずこれらの文字が等しいものとして考えると、その前の積は最初の文字を1とする二項式のn乗になりますから、この二項式のst次方の係数がnつの文字をs回ずつ取り出す組み合わせの数となります。そしてこの数は既知の二項定理によって得られます。

確率論における計算 27

各組合せにおける文字の位置に注意を払うことが重要である。最初の文字と二番目の文字が結合される場合、二番目の文字は最初のまたは第二位の位置に配置され得るため、2つの組合せが生じる。これらの組合せに第三の文字を追加するとき、各組合せにおいて第一、第二、第三の順位を与えられるため、それぞれの他の2つの組合せに対して3つの組合せが形成され、全体として6つの組合せが得られる。これから容易にわかるように、s個の文字の並べ替え可能な数は、1からsまでの数の積であることが分かる。文字のそれぞれの位置を考慮するためには、n文字からs個を選ぶ組合せの数に対して、この積でs回乗算を行う必要がある。これは二項係数が表すこの数の係数の分母を取り除くことに等しい。

$n$ 個の数字から $r$ 個ずつ引くロトゲームを考えましょう。ある一回の抽選で与えられた $s$ 個の特定の数字が当たる確率を求めます。これを求めるために、まず全ての可能な場合の数（つまり $n$ 個から $r$ 個ずつ取り出す組み合わせの数）を分母とする分数を考えます。そして、分子は与えられた $s$ 個の数字が含まれるすべての組み合わせの数です。この最後の数は明らかに残りの数字（つまり $n-s$ 個）から $r$ 個ずつ取り出す組み合わせの数と等しいことがわかります。

この分数は求める確率であり、容易に求められる形で分子が $s$ 個の数字を $r$ 个取る組み合わせの数、分母が $n$ 個の数字を $s$ 個ずつ取り出す組み合わせの数となることが分かります。フランスのロトでは、90個の数字から5個ずつ引くものがあり、特定の組み合わせが当たる確率は $\frac{\binom{5}{s}}{\binom{90}{s}}$ となります。平等な賭けのために、このロトは賭け金の18倍を支払うべきです。

90個の数字から2つずつ取り出す組み合わせの総数は4005であり、5個ずつ取り出す組み合わせの総数は10となります。したがって、特定のペアが当たる確率は $\frac{10}{4005} = \frac{2}{801}$ となり、ロトは賭け金の401.5倍を支払うべきです。特定のトリオ（3つ）が当たる確率は約11748倍、特定のクォータニー（4つ）が当たる確率は約511038倍、そして特定のクインティ（5つ）が当たる確率は約43949268倍となります。このロトはプレイヤーにこれらの利点をほとんど提供していません。

urn 内に白球 a 個、黒球 b 個が存在し、引いた球を元の urn に戻すという条件下で、m 回の引き出しで白球が a 個、n - m 回の引き出しで黒球が b 個となる確率を求めよ。各回の引き出しが起こる可能性は a + b 個あることが明らかである。二回目の各結果は第一回のすべての結果と組み合わせることができるので、二回の引き出しにおける可能な結果の数は二項式 (a + b) の平方となる。この平方展開において、a の平方は白球が 2 回引かれた場合の可能性を表し、a と b の倍積は白球と黒球がそれぞれ一度ずつ引かれた場合の可能性を表す。最後に、b の平方は二回とも黒球が引かれた場合の可能性を表す。同様にして考えると、二項式 (a + b) の m 乗は、m 回の白球と n - m 回の黒球がそれぞれ一度ずつ引かれた場合の可能性を含むことが一般にわかる。

確率論における計算. 29

n回の引く場合に可能なすべての組み合わせ数を表す式は、そのべき展開においてaのm乗項がb倍された項は、m個の白球とn-m個の黒球を引き得る組み合わせ数を表します。この項を二項定理の全べき乗で割ると、m個の白球とn-m個の黒球を引く確率が得られます。aとa+bの比が一つの引くときに白球を引く確率を表し、bとa+bの比が一つの引くときに黒球を引く確率を表す場合、これらの確率をfとgと呼びます。n回の引く中でm個の白球を引く確率は、二項定理P+qのm乗項が展開された式においてfのm乗項が表す確率となります。このように、関数が変数をそれぞれの差分増加させたときの連続する条件を比較することで、提案された問題はしばしば非常に単純な比を提供します。これが通常微分方程式または部分微分方程式と呼ばれるものです。変数が一つである場合が通常微分方程式で、複数の場合が部分微分方程式です。この方法のいくつかの例を見てみましょう。

三名の能力が等しいプレイヤーが以下の条件で一緒にゲームをプレイします：最初の二人のプレイヤーのうち、勝者が出るとそのプレイヤーは第三のプレイヤーと対戦し、もし彼が勝つならゲームは終了します。しかし、彼が敗れた場合、勝者は第二のプレイヤーと対戦し続け、いずれかのプレイヤーが連続して他の二人を制するまでゲームは進行します。このとき、ゲームが特定のn回目のプレイで終了する確率を求めます。そして、ゲームがちょうど nth 回目で終わる確率を見つけ出しましょう。そのためには、勝者が nth-1 回目のプレイでゲームに参加し、その後のプレイで勝つべきです。しかし、もし nth-1 回目のプレイで彼が敗れ、直前に他のプレイヤーを制した相手に敗れた場合、ゲームはこのプレイで終了します。したがって、プレイヤーが nth-1 回目のプレイでゲームに参加し、それを勝つ確率は、ゲームがこのプレイで終わる確率と等しいです。そして、このプレイヤーが次のプレイで勝つためにはゲームが nth 回目で終了するための確率は、前回の確率の半分になります。この確率は明らかに n の関数であり、この関数は n が1減ったときの同じ関数の半分と等しいです。この等式は普通の有限差分方程式の一種です。

私たちはその使用によって、ゲームが特定の一手で終了する確率を簡単に決定することができます。明らかに、ゲームは第二手以前には終了しないでしょう；そしてこれは最初の二人のプレイヤーのうちの一人が相手を破った後、第三のプレイヤーに対して第二手で勝つことが必要であるためです。この一手でゲームが終了する確率は$p$です。従って、前の式により、ゲームが続く各々の確率は、第三手では$p$、第四手では $\frac{1}{8}$、そしてその以降も同様に推移します。

確率論における微積分. 31

一般には、n手目では$p$のn-1乗となります。これらの$p$のべき乗の総和は最後のべき乗より1小さいことから、n手以内でゲームが終了する確率です。

次の問題について再考しましょう。これは確率を用いて解くことができる最初の難問であり、PascalはFermatにこの問題を解決するように提案しました。二人のプレイヤーAとBが等しい技術を持っている場合、彼らは一定の条件のもとで一緒にゲームをプレイします。その条件は、最初に相手を指定された回数以上勝つプレイヤーがゲームを制し、賭け金の合計を獲得することです。いくつかの投げ合いを行った後、プレイヤーたちはゲームを終了せずに引分することに同意します：彼らの間でその合計がどのように分割されるべきかを求めます。明らかに、それぞれの部分はゲームを勝つ確率に比例すべきです。問題はそれらの確率の決定に帰着されます。これらの確率は显然的に各プレイヤーが指定された回数に達するまでの欠けているポイントの数によって依存します。したがって、Aの確率は二つの数字（私たちはこれらを指数と呼びます）の関数となります。もし両プレイヤーがさらに一投をプレイすることに合意した場合（この合意は彼らの状況を変えない条件のもとでは常に新しい勝つ確率に比例して分割される）、Aがその投げ合いを勝ち、その結果彼が欠けているポイント数が1減る一方で、Bが勝利し、最後のプレイヤーが欠けているポイント数が1減少する場合があります。それぞれのこれらの状況の確率は\ ;したがって求められる関数は、最初の指数を1減らす関数の半分に加えて、同じ関数の2番目の変数を1減らす関数の半分で与えられます。この等式は部分微分方程式と呼ばれる方程式の一つです。

私たちは、その使用により最小の数値で確率を決定することができます。プレイヤーAが一つも点を欠いていない場合や最初の指標がゼロであるとき、その確率またはそれを表現する関数は1になります。また、この関数は第二の指標とともに0になります。したがって、プレイヤーAが一つだけ点を欠いていると仮定すると、彼の確率はf, f, |, などとなります。Bが一つ、二つ、三つなどの点を欠いている場合に応じてです。一般には、これはBが欠いている点の数分だけ1から$p$のべき乗が引かれたものになります。したがって、プレイヤーAが二つの点を欠いていると仮定すると、彼の確率はJ, $p$, $\frac{1}{4}$, などとなります。Bが一つ、二つ、三つなどの点を欠いている場合に応じてです。さらに、プレイヤーAが三つの点を欠いていると仮定します。

このように、差分方程式を使って連続する量の値を得る方法は長い作業で労働集約的です。幾何学者たちは、この方程式を満足する一般の指標関数の方法を探求し、特定のケースに対して必要なのはその関数に適応した指標の値を代入することだけです。この主題について一貫して考えましょう。そのためには、水平線上に並んだ項の系列を考えます。各項は前の項から与えられた法則に従って導き出されます。これらの項とその指数（順位を示す数）との間に関係式があると仮定します。

確率論における微分方程式の計算. 33

連続する項において、ある数列の項が占める位置を表す方程式は、単一の添字による有限差分の方程式と呼ぶ。この方程式の階数や次数とは、その方程式の両端の項のランクの差である。この方程式を使用することで、数列の各項を順次決定し、無限に延長することが可能だが、そのためには数列のランクが方程式の階数と等しい数の項を知っていることが必要である。これらの項は、数列の一般項の表記式や差分方程式の積分における任意定数となる。

次に、先ほどの数列の項の下に水平方向に並んだ二番目の数列を想像し、さらにその二番目の数列の下に三番目の水平方向の数列を想像する。そして無限までこの操作を繰り返すとしよう。また、これらのすべての数列の項が、横と縦の連続した項間で一般の方程式によって結びつけられていることを想像する。その方程式は、各項のランクを二つの指標で表すものである。この方程式は二重添字による有限差分の方程式と呼ばれる。

さらに同じように、先ほどの数列の平面の下に類似した二番目の平面を想像し、その二番目の平面上の項が最初の平面上の各項に対応する位置にあるように配置する。また、この二番目の平面上の下に三番目の平面を想像し、無限までこの操作を繰り返すとしよう。そしてこれらのすべての項が長さ、幅、深さという三次元での連続した各項間で方程式によって結びつけられることを想像する。その方程式は三次添字による有限差分の方程式と呼ばれる。

最後に、空間の次元に関係なく抽象的に考え、一般的な観点から、ある数の指標の関数である一定の量の系を想像し、これらの量とその指標に対する相対的な差異および指標自体について、量の数と同じ数の方程式があるとする。これらの方程式は特定の数の指標による部分的有限差分となる。

これらの方程式を使用することで、これらの量を順次決定することができる。一方で、一つの指標による方程式が一定数の項の列を知らなければならないように、二つの指標による方程式は一または複数の系列の行を知らなければならない。その一般的な項が各々一つの指標の任意の関数として表されるべきである。同様に、三つの指標による方程式は一または複数の系列の平面を知らなければならない。その一般的な項が各々二つの指標の任意の関数として表されるべきである。そして、これらの場合全てで、順次消去によって決定できる一定の列項がある。

これらの方程式のうちから消去するすべての方程式が同じ方程式系に含まれているため、これらの順次項の表現は一つの一般的な表現に収められるべきであり、それが指標のランクを定める指標の関数である。この表現は提案された差分方程式の積分であり、その探求は積分解析の対象となる。

テイラーは『増分法則』という著作の中で有限差分の線形方程式を最初に考慮した人物である。彼はそれらの方程式を積分する方法を示している。

確率論における微積分学. 35

一次方程式の係数と最終項を持つ関数、級数の指数。実際、算術および幾何級数の項との関係は常に考慮されてきたもので最も単純な差分方程式の例であるが、それらはこの観点から考慮されることなく存在していた。それは一般的理論に付随し、その理論へと導くものであり、従って真の発見と言える。

同じ頃、モイブールは定数係数を持つ一定次数の有限差分方程式を「再帰系列」として考えていた。彼はそれらを非常に巧妙な方法で積分することができた。発明者の進歩を追うことは常に興味深いので、私はその方法を説明するために、三項間関係が与えられた再帰系列に適用することにする。まず、彼は連続項の関係を考える。これは幾何級数の場合で、二項間方程式として表される。この式を1未満の項に参照し、定数因子で乗じて最初の方程式からその積を引くことで、幾何級数の連続項間の関係を得ることができる。モイブールは次に、同じ因子を持つ別の級数を考える。彼はこの新しい級数の項の指数を1減らし、それらの方程式を最初の級数の項の比で乗じて、その積から二番目の級数の方程式を引くことで、この級数の連続項間の関係を得ることができる。

$$
\begin{aligned}
& \text{まず、幾何級数の連続項の関係を考えます。} \\
& a_{n+1} = ra_n, \quad (r \text{は定数}) \\
& \text{これを 1 未満の項に参照し、定数因子 } k \text{ で乗じて最初の方程式からその積を引きます。} \\
& ka_{n+1} = rka_n, \quad (r \text{は定数}) \\
& ka_{n+1} - ra_n = 0. \\
& \text{次に、同じ因子を持つ別の級数を考えます。} \\
& b_{n+1} = rb_n, \quad (r \text{は定数}) \\
& \text{この新しい級数の項の指数を 1 减らし、それらの方程式を最初の級数の項の比で乗じて、その積から二番目の級数の方程式を引きます。} \\
& r(b_{n+1} - b_n) = rb_n - rb_{n-1}, \quad (r \text{は定数}) \\
& r(b_{n+1} - b_n) - rb_n + rb_{n-1} = 0. \\
\end{aligned}
$$

この条件で得られる関係は、この級数の連続項間の完全な関係である。

彼は、最初の級数に対する結果に似たものを発見した。次に、彼は、これらの連続する項のうちのどの三つでも同じ比が存在することを観察し、二つの級数の各項を逐次的に加えると、その比は依然として存在すると述べる。彼はこの比の係数を提案された再帰系列の項の関係の係数に比較し、二つの幾何学的級数の比率を決定する方程式を見つけた。その方程式の根がこれらの比率であることが判明した。

モイブールは、最初の二項級数の項を使用して任意定数を決定することで、再帰系列を二つの幾何学的級数に分解する巧妙な過程で、この問題を解決した。この方法は実際にはダランベールが無限小差分方程式の積分のために使用し、ラグランジュが有限差分の類似方程式に変換したものである。

最後に、私は部分的な有限差分の方程式について線形方程式を考えた。最初は再帰系列という名前で、その後その本来の名前に戻した。これらのすべての方程式を統一的に解く最も一般的かつ簡単な方法は、判別関数の考慮に基づいたものであり、そのアイデアはここに示される。

もし私たちは変数 / に関する関数 V をこの変数の冪級数展開と仮定すると、これらの冪のどれか一つの係数は、この冪の指数または指標（これを x と呼ぶ）に関連する関数となる。V は、この係数または指標の関数を指すものとして私自身が呼んでいる判別関数である。

もし V の展開列を同じ変数に関する別の関数で乗じると、例えば、

$$
V \cdot f(x)
$$

のようにすると、これが部分的な有限差分方程式の解法の一例となる。

単位数にこの変数の2倍を加えたとき、その積は新しい判別関数となり、この変数tのx乗項の係数はVの同じ指数項の係数と等しく、さらにVの係数の2倍となる。したがって、積における変数tの指数xの関数は、Vにおける同じ指数xの関数に2倍するものとなる。この指数xの関数は、Vの展開における同じ指数の関数の導関数であると結論付けることができる。私はこの原始関数を指すために、この導関数を表す文字dを用いることにしよう。この文字dによって示される微分は、Vの乗数Tに依存し、これをTと呼ぶことにする。そして、変数tのべき乗に似た比率でVを開発することも可能であると仮定する。

次にVをTのn乗で再び掛けると、積における変数tのx乗項の係数は、前の積における同じ指数項の係数と同じ形の導関数となる。これを表現するために、前の導関数の同一の文字8を使用し、その後ろに原始関数のxを一度書くことにより表現できる。しかし、それを一度だけ書く代わりに、その指数を2とする。

このように続けば、一般にVをTのn乗で掛けると、VとTのn乗の積における変数tのx乗項の係数は、原始関数の前に文字8とnを指数として用いることで表現できる。

たとえば、T を / の逆数とする場合を考えましょう。すると、F と T の積において / の r 番目のべきの係数は V における / のべきが 1 大きくなるときの係数となります。この係数は、V と nth べきの Z の積において / の r 番目のべきを持つ項の係数を表します。

次に、新たな関数 Z を / について V と同じようにべき級数で展開し、原始関数の前に置かれた文字 A を / の r 番目のべきを持つ項の係数と定義しましょう。この係数は V と nth べきの Z の積において / の r 番目のべきを持つ項の係数を表し、その表現は文字 A に n の指数を付けて原始関数の前に置くことで与えられます。

たとえば、Z が t - 1 の逆数である場合、V と Z の積において / の x 番目のべきの係数は V における / の (x + 1) 番目のべきの係数から / の x 番目のべきの係数を引いたものとなります。これは x の原始関数の有限差分を表します。したがって、文字 A は指数が 1 大きくなるときの原始関数の有限差分を示し、この文字の n 乗を原始関数の前に置くと、その原始関数の有限 n 番目の差分を表します。

また、T を / の逆数とする場合、7 は Z + 1 に等しいことがわかります。V と nth べきの T の積は V と nth べきの Z + 1 の積に等しくなります。このべきを Z のべきについて展開すると、V とその展開項の各項との積は、原始関数の有限差分を代入した同様の discriminant 関数となります。

THE CALCULUS OF PROBABILITIES. 39

ここで、Z のべきを代入する代わりに、原始関数の指数に対する有限差分を代入します。

現在、VをTのn乗で掛けると、指数xがn単位増加した原始関数を得ます。判別関数からその係数へと戻ると、この原始関数がn単位増加したものが、Z + \ のn乗展開に等しいことがわかります。ただし、この展開においてZのべきは原始関数の差分に置き換えられ、これらのべきの独立項は原始関数で掛け算されます。これにより、任意の数n単位増加した指数を持つ原始関数を得ることができます。

TとZが常に先述の値を保有するという仮定のもとでは、Zはbinomial T - 1に等しいことになります。VをZのn乗で掛けることは、Vをbinomial T - 1のn乗展開で掛けることに等しいです。判別関数からその係数へと戻す作業を行い、先ほどと同じようにすると、原始関数のn階差分は、binomial T - 1のn乗展開によって表され、Tのべきは指数がべきの指数に増加した同じ関数で置き換えられ、Tの独立項（単位）は原始関数で置き換えられます。これにより、この関数の連続する項から差分を得ることができます。

原始関数を表す変数 \( S \) を用いて、この関数の導関数が \( V \) の製品において \( x \) のべき乗 \( / \) を \( T \) で乗算するものとし、同様に \( V \) の製品において \( Z \) で乗算するものを表す変数 \( A \) を用いる。それにより先の議論から、任意の変数 \( x \) の関数を \( T \) および \( Z \) によって表した場合、これらの関数間のすべての恒等式の展開において \( T \) と \( Z \) を \( d \) および \( // \) に置き換えることができることが導かれる。ただし、原始関数を指標と共に級数形式で書き、この関数をこれらの変数の独立項にも乗算する必要がある。

我々はこの一般の結果を利用して、原始関数の指数xにおける任意の一定のべき乗を変数xが1ずつ変わるときの差分に変換し、逆もまた真であることを示すことができる。つまり、原始関数の指数xにおける任意の一定のべき乗を変数xが特定の数だけ変わるときの差分に変換することができる。

まず、T を 1 の i 乗で / - i で割った値、Z を常に 1 で / - 1 で割った値と仮定する。すると、V と T の積において x のべき乗の係数は、V 内の x + i のべき乗の係数から x のべき乗の係数を引いたものとなる。したがって、これは指数xが数iだけ変わるときの原始関数の有限差分である。

容易に理解できるように、T は二項式 Z + 1 の i 乗と 1 の差に等しい。また、T の w 乗はこの差の w 乗に等しい。この等式において T と Z をそれぞれ 6 と J に置き換え、展開した後各項の末尾に指数xの原始関数を付加すると、指数xが * 単位ずつ変わるときの w 乗差は同じ関数の差分の系列で表される。この系列は、

確率論における微積分学. 41

その差分を表現するものであり、かつそれが同一であるが、このような変換の力は類似した変換に宿る。

解析の一般性により、この表現において n を負数と仮定することができます。すると、tF および A の負の冪は積分を示します。実際、原始関数の n 番目の差分が判別関数として V と t の n 番目の二項式の商（1 - t）で割った積に与えられるとき、その原始関数である n 番目の積分は同じ差分の判別関数を n 番目の冪で減じたものを持つことがわかります。この冪は文字 A の同じ冪に対応し、それ自身も同様に同じ順序の積分を示します。指標 x は 1 ずつ変化するとき、負の冪の 6 も等しく x が i 単位ずつ変化する積分を表します。したがって、正の冪と差分、そして負の冪と積分における分析の合理性を見ることは最も明瞭かつ簡単な方法です。

もし、$T$ に示される関数が原始関数の前に置かれたとき零であれば、我々は有限差分方程式を得ます。その積分の discriminant function である F を得るためには、V と T の積においてすべての t のべきが消失し、差分方程式の階数未満のべきだけが残ることを観察します。そうすると V は T の分数であり、その分子は差分方程式の階数より一つ低い最高べきを持つ多項式となります。

この多項式的な分子における t の各べき（包括的に零のべきも含む）の任意係数は、原始関数の指数が順にゼロ、一、二と等しい x の値において与えられる原始関数の値によって決定されます。差分方程式が与えられたとき、T を求めることで最初の項を第一項のメンバーに、第二項のメンバーには零を入れます。最大の指数を持つ関数の代わりに第一べきの t と置き換え、次に指数が一つ減少した原始関数の代わりに第二べきの t と置き換えるなどして最初の項のメンバーを操作します。

前記 V の表現の展開における t の x 空白べきの係数は、原始関数または有限差分方程式の積分である原始関数 x を得ます。解析学はこの展開のために様々な方法を提供し、その中から提案された問題に最も適したものを選択することができます。これはこの積分法の利点です。

今、V を変数 / と /' の二つの変数の関数として考えましょう。V はこれらの変数のべき乗や積に従って展開されます。任意のべき乗 x と x' の積の係数は、これらのべき乗の指数 x と x' の関数となります。この関数を原始関数の discriminant function と呼びます。

次に、V を変数 t と /' の二つの変数の関数 T で乘じてみましょう。T も V と同じようにこれらの変数のべき乗や積に従って展開されます。その積は原始関数の導関数の discriminant function となります。例えば、T が変数 / と t' の和から2を引いたものである場合、この導関数は指数 x を1減じた原始関数と、同じ原始関数の指数 x' を1減じたものを引いて2倍した原始関数の差となります。任意の T を d という文字で表すことにします。この導関数は V と T の w乗の積であり、その積は原始関数の ;/乗の文字 8 が置かれた原始関数の導関数の discriminant function です。したがって、単変数の関数に対する定理に類似する定理が得られます。

次のように仮定すると、部分差分方程式を得ます。例えば、変数$\phi$を$t'$−2に等しいとするTを再度定義する場合、0は原始関数の指数xを1減じたものと、同様な原始関数の指数x'を1減じたものの差に加えて、その原始関数の2倍からなる方程式となります。

この原始関数またはこの方程式の積分の判別関数Vは、Tとの積がt'によって生成されるすべての項を含まないべきです。しかしFはtとt'のそれぞれの冪を別々に含めることができ、つまり任意のtの関数と任意の$t'$の関数となります。その結果、Vは分子がこれらの2つの任意の関数の和であり、分母がTである分数になります。

この分数の展開において、tのr乗とt'のx'乗の積の係数は前述の部分差分方程式の積分となります。このような方程式を積分する方法としては、様々な解析的手法を用いた有理分数の展開過程を通じて最も簡単で理解しやすいと考えます。

この事項についてより詳しい説明は、微分積分の助けなしではほとんど理解できないでしょう。

無限小部分微分の方程式を有限小部分微分の方程式として考え、何も省略せずに扱うことで、その計算における曖昧な点に光を当てることができる。幾何学者たちの間で大きな議論が行われたこれらの計算について、このような方法で私は連続しない関数を積分に入れる可能性を証明した。ただし、その不連続性がこれらの方程式の微分またはそれ以上の階数の微分に限られている場合である。

微分積算学における超越的な結果は、理解の抽象化と同様に一般的な符号であり、その真の意味は、これらを導いた原始的概念に戻って形而上学的分析を行うことでしか確認することができない。これにはしばしば大きな困難が伴うため、人間の心は未来へ向かうよりも自身の中へ退くことをまだあまり好まない。

無限小部分微分と有限小部分微分との比較も同様に、微分積算学の形而上学的な原理に大きな光を当てる能力がある。

変数の増加がEである関数の有限n次の差分をEのn乗で割ることで証明できる。この商は比級数によってEの増加の冪から形成される最初の項で、Eが減少するにつれてこの最初の項に近づき、その違いは任意に小さい量より小さい量だけである。

この最初の項は次のように定義される列の極限であり、微分積算学において関数の無限小n次の差分をその無限小増加のn乗で割ったものと一致する。

確率論の計算 45

この項は列の極限であり、微分積算学において関数の無限小n次の差分をその無限小増加のn乗で割ったものの表現である。

この観点から無限小の差異を考慮すると、微分 calculus の様々な操作は、同一の式の展開において有限項または変数の増分に依存しない項を別々に比較することに他なりません。これは厳密に正確であり、これらの増分は未定義です。したがって、微分 calculus は他の代数的操作と同様に完全な正確さを持っています。

微分 calculus の幾何学や力学への適用においても同じ程度の正確さが見られます。例えば、曲線が二つの隣接する点で切れる正接線を考えると、これらの二点の縦座標間の距離を E とします。E は最初の縦座標から次の縦座標までの横座標の増分です。この E によって対応する縦座标的増分は、最初の縦座標とその正接線との比に E を掛けることで得られます。その後、曲線の方程式において最初の縦座標をこの増分で増加させることにより、次の縦座標に関する方程式を得ます。これらの二つの方程式の差は、E のべき乗の比率で展開し E で割った第三の方程式となり、その最初の項が E に依存しない限界となります。この項がゼロであることは、正接線の限界が明らかに接線（ subtangent）となることを意味します。

この特異な方法はフェルマーによって開発され、彼はそれを超越曲線にも拡張しました。この偉大な数学者は、 абсцисса の増分を E という文字で表し、この増分の最初のべき乗のみを考えることにより、微分積分学によって我々が行っているように、曲線の subtangent, 割点、縦座標の最大値と最小値、そして一般に有理関数のそれらを正確に決定します。彼はまた、デカルトの手紙集に挿入された光の屈折問題の美しい解法から、彼は無理関数を無理性を取り除くために根をべき乗することでその方法をどのように拡張するかを知っていることもわかります。したがってフェルマーは微分積分学の真正な発見者とみなされるべきです。

ニュートンはその後、流動法という彼の方法でこの計算をより分析的にし、美しい二項定理によって過程を簡素化し、一般化しました。最後に、ほぼ同じ時期にライプニッツは微分積分学に独自の表記法を加え、有限から無限小へと移行することにより、一般的な結果を表現する利点に加えて、差分や量の和の最初の近似値を与えるという利点も追加しました。この表記法は部分微分の計算にも適しています。

しばしば、項や因子が非常に多くなり、数値的な代入が実行不可能な表現に遭遇します。これは特に多くの事象を考慮する確率論の問題で起こります。しかし、結果が何らかの確率で示されるためには、公式の数値的な値が必要です。これらの事象は乗算によって発展します。特に必要なのは、

確率論の計算 47

この確率が無限に近づくにつれて、それが最終的に確定性に達するという法則は、事象の数が無限である場合に成立します。この法則を得るためには、私は微分を因子の高次方程式と乗じた定積分が積分によって多くの項と因子を持つ公式を与えることを考慮しました。この観察結果から、解析の複雑な表現や差分方程式の積分を類似した積分に変換するアイデアを得ました。私は、被積分関数と積分の限界を同時に得る方法でこの条件を満たしました。これは、被積分関数が表現と提案された方程式の判別関数であり、そのような方法が判別関数理論に付加される補完であるという特徴を持ちます。さらに、それは定積分を収束する級数に帰着させる問題に過ぎません。私は、表現される式ほど複雑なほど収束速度がより速くなるプロセスでこれを達成しましたので、必要性が増すにつれてより正確になります。しばしば、その因子として円周率と直径の比の平方根が含まれていますし、時には無限に存在する他の超越関数にも依存することがあります。

重要な注記があり、これは分析の広範な一般性に関連しており、この方法を差分方程式や確率理論が頻繁に提示する公式にも拡張できるものである。それは、実数で正である限界を仮定した定積分の級数が、これらの限界を決定する方程式が負または虚数の根だけを持つ場合でも同様に行われるという点についてである。私は最初に利用したこれらの正から負への移行と実数から虚数への移行は、多くの単純な定積分の値へとさらに進み、それらを直接証明することにより明らかにした。

これらからの移行は、幾何学者が初期には極端な慎重さを持って、その後完全な信頼性を得て長年にわたって使用してきた帰納法や類推法と並行する発見の手段として考慮できる。多くの例がその有用性を証明したためである。一方で、これらの様々な方法によって得られた結果は、常に直接的な証明により確認することが必要である。

私は先述の方法の集まりを差別関数の計算と名付けた。この計算は、確率論的解析理論というタイトルで出版された私の著作の基礎となっている。それは、ある量自身またはその完全な正のべき乗を繰り返し乗算する際、その表現する文字の上部にこれらのべき乗の度数を表す数字を記述することで単純な概念と結びついている。

この表記は、デカルトが幾何学において用いたもので、その後重要な著作が出版されてから広く採用されたものである。特に、この偉大な数学者が現代の微分積分学の基礎を築くために確立した曲線論と変数関数の理論と比べると、小さなことではある。しかし、分析の言語は、最も完璧なものであり、その自身が発見の強力な道具であるため、特に必要で幸運にも考えられた表記は、新たな計算法の芽となるのである。このことは次の例によって明らかになる。

解析の言語、最高のものであるが故に、その表記は必要でうまく考案されたときには、新しい計算法の種々を孕むものであることが認識される。これは以下の例により明らかになる。

ウォリスは、『無限の算術』という著作において、分析の進展に大きく貢献した作品の一つで、帰納法と類推を追跡することに特に興味を持っていた。彼は、ある文字の指数を2で割り、3で割ったときの商が、その結果はカルテルの表記法であり、割り算が可能であれば、その指数は量の平方根、立方根など、分子に示されるべき指数によって表現される量の根となると述べた。この結果を類推により拡張し、割り算が不可能な場合を考えたとき、彼は分数の指数を持つ量を、分母が示す度数の根として考えた。つまり、分子が示すべき指数によって表現される量の根とした。

ウォリスはさらに、カルテルの表記法によれば、同じ文字の冪の乗算はその指数を加算し、除算は被除数の指数から除数の指数を引くことであると観察した。特に、二番目の指数が最初の指数よりも大きい場合に限ってこの結果を拡張した。最初の指数が第二の指数以上または等しい場合、差はゼロや負の値となる。彼はそれらの指数が負の場合を仮定し、その冪は同じ指数を正として取り、それが示す量で1を割ることを述べた。

これらの観察から、彼は一般的に単項微分を統合し、特定の種類の二項微分の確定積分を推測した。そして、これらの積分を表す数列の法則性と、補間と幸運な帰納法を見出し、それらは定積分の計算の芽を含んでいることを悟った。これは私の新しい確率論の基礎の一端となっている。

ウォリスは無限積によって円の面積に対するその直径の平方の比を表現し、それを打ち切ることでこの比がより近づく限界に制限されることを見出した。これは分析において最も興味深い結果の一つである。しかし、ウォリスが分数指数を持つ根号の冪についてよく考えていたにもかかわらず、それらの冪を以前の方法で記述し続けたことに驚きを感じる。

私は誤りがないとすれば、ニュートンはオールドルムブルグへの手紙の中でこれらの冪を分数指数で表す最初の人物だった。ウォリスが帰納法を非常に美しく用いたように、二項の冪の指数とその展開における項の係数を比較し、整数かつ正の指数の場合に限ってこの指数の法則を見つけ出し、類推により分数や負の指数にも適用した。これらの様々な結果はデカルトの表記法に基づいており、彼が分析の進歩に与えた影響を示している。

さらに、これは対数の最も単純で公正な概念を与える利点がある。対数は実質的にある量の指数であり、その連続した冪は無限小の段階で増加し、すべての数を表現できるからである。

しかし、この表記が受けた最も重要な拡張は変数の指数にあり、これは指数計算というものを成し立たせ、その成果は非常に豊かである。

確率論の微積分学。5章

現代分析の分野において、リーブニッツは変数の指数を使用して超越関数を示す最初の人物であり、これにより有限な関数が構成される要素の体系を完成させた。つまり、任意の有限で明示的な変数関数は最終的には加法、減法、乗法、除法および定数または変数のべき乗によって組み合わせられた単純な量に帰着できる。これらの要素から形成された方程式の根は変数の暗黙的関数である。

リーブニッツは自己の微分符号にも同じ指数を用いたが、これは同じ量の重複した積を示すものではなく、同じ関数の重複した微分を示すものであった。このカルテシアン記法への新しい拡張により、リーブニッツは正のべき乗と微分、負のべき乗と積分との類似性を見出した。

ラグランジュはこの不思議な類似性をすべての発展に追随し、差分と積分の変換における多項式の有限増分や無限小の増分を持つ変数についての一般公式を得た。しかし彼はその証明を与えていない。

判別関数の理論はカルテシアン記法を一部の特性にまで拡張し、これらの特性によって示されるべき乗と操作の類似性を証明するものであるから、これは依然として符号の指数計算と考えることができる。差分方程式およびその積分の級数展開や解は、これから極めて容易に導き出される。

以上の翻訳では、原著の論理構造と語調を尊重しつつ、自然な日本語表現に仕上げています。

第二部。

確率計算の応用。

\newpage

\chapter{偶然の遊戯（Games of Chance）}
\label{ch:vi}

\begin{center}
\large\textit{公平な賭け・破産の確率・サイコロとトランプの組合せ論}
\end{center}
\vspace{4mm}

ゲームが提示する組合せのパターンは、確率に関する最初の研究の対象となりました。これらの無数の組合せの中には、多くのものは容易に計算可能ですが、他のものはより複雑な計算を必要とします。そして、組合せがより複雑になるとその難易度も増大し、それらを克服する欲望と好奇心は幾何学者たちをこの種の解析をさらに発展させることに駆り立てました。

既に述べたように、当選確率を決定するために Lottery の利益は組合せ理論によって容易に計算できます。しかし、例えばすべての番号が引かれるまでにどれだけの抽選回数が必要かを知るためには、より困難な問題となります。ここで n を番号の総数、r を各抽選で引かれる番号の数、i を未知の抽選回数とします。

すべての 53 番号が引かれる確率は、「i 次の連続する r 個の数の積の i 乗の有限差分」によって表されます。番号 n が相当大きい場合、この確率を J に等しくする i の値を求めることは少なくとも、その差分を非常に収束する級数に変換しない限り不可能となります。これは、非常に大きな数の関数の近似法により容易に行うことができます。

したがって、1万個の番号から各抽選で 1 個ずつ引かれる Lottery の場合、すべての番号が 95767 回目の抽選までに引かれることの確率は不利であり、95768 回目の抽選まで待つ方が有利となります。フランスの Lottery の場合、この賭けは 85 回目までの抽選では不利ですが、86 回目から有利となります。

二つのプレイヤー、AとBが表と裏のゲームをプレイする状況を考えましょう。各投げで表が出ればBに1カウンターを渡し、裏が出ればAに1カウンターを渡します。Bのカウンター数は制限がありますが、Aのカウンター数は無限であり、ゲームはBがもうカウンターを持たないときだけ終了します。ゲームが終了するまでの投げ回数について、一対一で賭けられるべき投げ回数を問います。

ゲームがi回目の投げで終了する確率の表現は、Bのカウンター数が相当大きい場合、多くの項と因子を持つ級数によって与えられます。この級数の値を求める未知数iを探すことは、その数が相当大きい場合、不可能となります。しかし、収束性が高い級数に簡略化することで可能になります。

先ほど述べた方法を適用すると、未知数について非常に単純な表現を見つけることができます。例えば、Bが100カウンターを持っている場合、ゲームが23780回目の投げで終了する確率は少々1より小さいとし、一方で23781回目の投げで終了する確率は少々1より大きいとします。

これらの二つの例を既に与えた例と合わせて、ゲームの問題が分析の完成度向上にどのように貢献したかが示されます。

\newpage

\chapter{等しいと仮定された偶然に潜在する未知の不平等}
\label{ch:vii}

\begin{center}
\large\textit{前提条件のゆらぎと事後検証によるベイズ的修正}
\end{center}
\vspace{4mm}

確率が等しい場合における不均一性は、確率計算の結果に感覚的な影響を与え、特に注意すべきです。硬貨投げのゲームを考えましょう。硬貨のどちらの面も投げるのに同様の容易さがあると仮定しましょう。すると最初の一回で表が出る確率は$$\frac{1}{2}$$であり、連続して表が出る確率は$$\left(\frac{1}{2}\right)^2 = \frac{1}{4}$$です。しかし、硬貨の中にどちらかの面がもう少し出やすいという不均一性がある場合（どの面が有利であるかは知りません）、最初の一回で表が出る確率は依然として$$\frac{1}{2}$$となります。なぜなら、不均一性が有利であれば単純な事象の確率は増加し、不利であれば減少するからです。しかし、この不確定性の中で、連続して表が出る確率は増加します。実際、その確率は最初の一回で表が出る確率を$$\frac{1}{2}$$とすると、それが有利である確率に掛け算することで得られます。

未知の不均一性が存在する場合における確率について56節
未知の不均一性が存在する場合における確率について57節

最初の投げで表が出ると、次回も表が出る確率が高まる。しかし、最初に表が出ること自体は、コインの不均一性が表を有利に働かせていることを信じる理由となる。この未知の不均一性は、次の投げでも表が出る確率を増加させ、その結果、両者の確率の積も増加する。

これを計算にかけるため、この不均一性が単純な有利な事象の確率を20分の1だけ増加させるという仮定を立てましょう。もしその事象が表である場合、その確率は$p$ + -fo すなわち \pm となり、二回続けて表が出る確率は（\pm）$^2$ すなわち |f $p$ となります。もし有利な事象が裏である場合、表の確率は | - > すなわち *V となり、二回続けて表が出る確率は TVo となります。

最初からどちらの不均一性が有利に働くかを信じる理由がないため、両方の事象が等しく可能であると判断される場合、複合事象「表 表」の確率は前記の二つの確率を足し合わせた半分を求めることで得られます。これにより $p$$p$| という確率を得られ、これは $p$ よりも J-g- すなわち有利性 -fa の平方だけ大きいことが分かります。

同様に、「裏 裏」の確率は $p$f となります。しかし、「表 裏」と「裏 表」の確率はそれぞれ jVT であり、これら四つの確率の合計は確実性または1と等しくなります。

したがって、常に単純な可能性が等しい事象を有利に働く未知の定数要因は、同じ単純な事象の再発生の確率を増加させることが分かります。

偶数回の投げ入れにおいて、表と裏はどちらも偶数回または奇数回起こるはずです。これらのそれぞれの場合の確率は、両面が等しい場合 $p$ です；しかし、それらの間にある未知の不均衡がある場合は、その不均衡は最初のケースに常に有利となります。

技能が等しい2人のプレイヤーが、各投げ入れで負けた方が相手に1つのコインを渡す条件でゲームを行い、あるプレイヤーが最後にコインを持たなくなるまで続けます。確率の計算によれば、ゲームの均衡状態では、プレイヤーの投げ入れは彼らのコインの数に対して逆比例するべきです。しかし、プレイヤー間にある小さな未知の不均衡がある場合、それはコインの数が少ないプレイヤーに有利となります。彼の勝利確率は、プレイヤーがコインの数を倍数または三倍にする合意があれば増加しますし、その確率は無限大になる場合の他のプレイヤーの確率と等しくなります。

これらの未知の不均衡の影響を修正するためには、それら自身を偶然の機会に委ねることができます。例えば、表と裏のゲームにおいて、最初のコインと一緒に二つ目のコインも投げ入れ、常に二つ目のコインで上向きになる面を表とする合意があれば、最初のコインで連続して表が出る確率は単独のコインの場合よりも著しく $p$ に近づきます。最後のケースでは、この差異は未知の不均衡が有利とする最初のコインの面に対する小さな可能性の増加の平方となります；一方、他の場合ではこれは

未知の不均衡における機会.         59

この平方を、二番目のコインに対する対応する平方に四倍して積みましょう。

次のようにします。1から100までの百個の数を順序どおりにurn（玉入れ）に入れ、それをよく振って混ぜます。その後、一つの数を引くとします。明らかに、混ぜ方が適切であれば、各数が引かれる確率は同じになります。しかし、各数がurnに入られた順序によって微小な差異がある懸念があれば、最初のurnから引かれた順序で二番目のurnに入れ直し、それをよく振って再び混ぜます。これにより、これらの差異は著しく小さくなります。さらに、三番目のurn、四番目のurnなどに入れて同じ操作を繰り返すことで、差異はさらに減少し、二番目のurnではすでにほとんど感じられなくなるまでになります。

\newpage

\chapter{事象の無限回の反復から生じる極限確率法則}
\label{ch:viii}

\begin{center}
\large\textit{大数の法則・中心極限定理・ド・モアブル＝ラプラスの定理と正規分布}
\end{center}
\vspace{4mm}

無数の事象から生じる結果。

確率という名の不明瞭な原因によって、事象の進行が不確定で規則性がないと理解される中で、これらの事象が増えるにつれて、驚くべき規則性が現れ、設計されたものであるかのように見受けられます。これにより、神の存在が証明されるように思われます。しかし、これを深く考えると、この規則性は単純な事象のそれぞれの可能性の開発に過ぎず、より高い確率を持つ事象はより頻繁に現れるべきであることが認識されます。

たとえば、白い球と黒い球が入った urn を想像してみましょう。そして、各回の抽選後、ボールを元の urn に戻すという条件で考えます。最初の数回の抽選では、白いボールと黒いボールの比は非常に不規則であることが多くなりますが、この不規則性の変動要因が全体として互いに打ち消し合うことで、多くの抽選を行うことで urn 内の白いボールと黒いボールの比率をより正確に把握できるようになります。つまり、それぞれの抽選における白いボールや黒いボールを引く確率も明らかになるのです。

この結果から以下の定理が導かれます。

白いボールを引いた回数に対する総引く回数の比が urn 内の白いボールと黒いボールの比率から与えられた区間を超えない確率は、無数の事象によって増加し、この区間がいかに狭くてもその確率は必ずしも確定的に近づきます。

この定理は直感的に理解できるものではあったが、解析によって証明するのは困難であった。その結果、最初にこの問題に関わった有名な幾何学者ジャン・ベルヌーイは、自身が与えた証明を非常に重視している。差別関数の微分 calculus をこの問題に適用すると、この定理だけでなく、観察された事象の比が真の比から特定の限界内だけずれる確率も容易に示すことができる。

前記の定理から、自然の行為の比はこれらの行為を大量に考えるときには非常にほぼ一定であるという結論を引き出すことができ、これは一般的な法則として考慮されるべきである。つまり、多くの年数を考えるとき、生産総額は大体同じであり、そのため人間は有用な予測によって季節の不規則性から身を守ることができる。この法則は、道徳的な原因による結果にも適用される。年間出生率に対する人口比と、結婚数に対する出生数比は僅かな変動しか見られない。パリでは年間出生数はほぼ同じであり、普通の季節における郵便局での不良な住所による投棄された手紙の数も年々大きく変わらないという話も聞く。ロンドンでも同様の観察が行われている。

この定理からまたもや、無限に延長された一連の事件において、長期的には規則的かつ恒常的な原因の作用が、不規則な原因の作用を上回るはずです。これは、当選 lottery の利益が農業の生産と同様に確実であることを可能にするのです；それらが準備している機会は、多くの投げ出しの中で利益を保証します。

有利で頻繁に付随する機会は、理、正義、人道という永遠の原理の観察に常に結びついており、これらの原理によって社会が確立され維持されるため、これらに従うことは大きな利点であり、それらから離れるのは重大な不便です。歴史を参照し、自分の経験を問いただすと、この計算の結果を支持するすべての事実を見ることができます。

理と人間の自然権に基づく制度が知恵ある人々によって確立され維持された国々における幸運な効果を考えてみてください。また、誠実性がその行動の基礎とした政府が得た利益を考え、慎重な正確さでそれを保つために受け入れられた犠牲を補填したことを考えましょう。

無限に増加する事件。63

彼らの侵略行為は、彼らに莫大な信用を国内で生み出し、また国外においても優越性をもたらした。しかし、対照的に、国家はしばしば、その指導者の野心と裏切りによって、いかに深い不幸の淵に落とされるかを思い起こすべきである。

偉大な国が征服の愛に陶然とし、世界支配を目論むたびに、独立心という感情は脅威を感じる諸民族間で連合を形成し、その結果としてしばしばその国自身が犠牲となる。また、多様な国家が拡大するか縮小するかを左右する変動要因の中においても、自然の限界が恒常的な要因として作用することにより最終的に優勢を占めるべきである。

帝国の安定と幸福にとって重要なのは、その境界を超えて国土を拡張しないことである。これらの要因の作用によって再び無休無怨に導かれる境界を超える国土を拡大することは、海が激しい嵐によって引き上げられた水滴が重力の力により再び海盆に戻るように、自然の流れに逆らうものである。

これは確率論的な計算の結果であり、多くの悲劇的な経験によって確認されている。歴史を恒常的要因による影響という視点から扱えば、好奇心に対する興味と同時に人間に対して最も有用な教訓を提供する利益も得られるだろう。

しばしば我々はこれらの原因が生み出す不可避免の結果を、その作用に至るまでの偶然的な状況に帰属させる。例えば、一つの民族が別の民族によって支配されるのは、広大な海や大きな距離が隔てている限り自然の流れから逸脱することである。長期的には、この恒常的要因は時間と共に作用する変動要因と連携し、その結果として支配された民族に自然の独立性を回復させたり、あるいは近隣の強大な国家との統合を促す力を持つことになるだろう。

多くの場合、特にリスクの分析において重要な場合においては、単純な事象の可能性が未知であり、私たちは観察された事象から指標を探し求めることで、その依存する原因についての推測を行う余地に迫られます。上記で述べられた原理を基にして、観察された事象から原因の確率を分析するために discriminant functions の解析を適用すると、以下の定理が導き出されます。

単純な事象または複数の単純な事象から成る事象（例えばゲームにおいて）が多数回繰り返される場合、観察された事象が最も確実に生じる可能性が高い単純な事象は、その観察が最大の確率で示唆するものであり、観察された事象が繰り返されるほどこの確率は増加し、無限回の繰り返しとなった場合、確実性に近づいていくことになります。

近似には二種類があります：一つは可能性のすべての側面における極限に対する相対的な近似であり、もう一つはこれらの可能性がこれらの極限内にある確率に関連しています。複合事象の繰り返しは、これらの極限が固定されたままである一方で、この確率をより多く増加させます；また、これらの確率が固定されたままである一方で、これらの極限間の区間をより狭めます。無限に近づくにつれて、この区間はゼロになり、確率は確定性へと変化します。

この定理を生後比について適用すると、

無限の事象の乗算. 65

ヨーロッパの諸国で観察された男児と女児の比率を比較すると、この比が22対21というところにほぼ等しいことから、男児の出産がより容易である可能性は極めて高いことが示されます。さらにナポレオンとサンクトペテルブルクでも同じ比率であることから、気候による影響はないと考えられます。従って、一般的な信念とは異なる可能性として、この男性優位の出生比は東洋においても存在するかもしれません。そこで私たちはフランスの学者たちにエジプトへ派遣し、この興味深い問題に関心を持たせるよう依頼しましたが、正確な出産に関する情報を得る困難さからその解決には至りませんでした。幸いにも、 Humboldt 氏はアメリカで彼の鋭さ、恒心性、勇敢さによって観察し収集した無数の新しい事項の中でこの問題に注意を払い、熱帯地域でもパリで観察される出生比と同じであることを発見しました。これにより私たちは男性優位の出産が人類全体に対する一般的な法則であると考えるべきだと考えられます。異なる種類の動物がこの点で従う法則は自然学者たちにとって注意を向けさせる価値があると私は思います。

ある場所で観察された多くの出産の数において、男児と女児の出生比がほとんど1に近い値を示すという事実は、一般法則から逸脱した結果であり、その存在しないことを結論付けるには十分ではない。この結果を得るためには、大規模なデータを使用し、それが高い確率で示唆されていることを確認する必要がある。

ブフォンは、例えば『政治算術』においてブルゴーニュのいくつかのコミュニティを引用し、そこで女児の出生が男児よりも多い例を挙げている。その中でもカルセル＝レ＝グリニヨンのコミュニティでは、5年間に2039人の出産があり、その内1026人が女児で983人が男児だった。

これらの数字は確かに大きなものであるが、それらはただ女児の出生の方がやや多い可能性を示唆し、その確率は -fa であり、これは四回続けてコイン投げで表が出ない確率よりも小さい。この確率は十分に高いものではなく、この異常性の原因を探るためには不十分である。そしておそらく、この異常性は百数年の間このコミュニティでの出産を追跡すれば消失するだろう。

出生簿は、市民の状態を確認するために丁寧に保管されており、大帝国の人口を決定するのに使用できるかもしれません。これには、住民の数を直接数えることなく、その inhabitant の数え上げを行う必要がありません ——これは労働集約的で正確性を得る上で難しい作業です。しかし、この目的を達成するためには、人口と年間出生数の比率を知る必要があります。最も精密な方法は、まず帝国全体のほぼ等しい面積に分布した地域を選んで、一般結果が地方状況に依存しないようにすることです；次に、これらの地域内のいくつかのコミュニティの住民を特定の時期に丁寧に数えることです；そして、その時期以前とその後の数年間に報告された出生の記述から、年間出生数に対応する平均数を求めることです。この数は、住民の数で割ることにより、人口に対する年間出生数の比率を求めることが可能になり、数え上げがより大きくなるにつれてその正確性が高まります。

無限の事件の乗算. 67

政府は類似の数え上げの有用性に気づき、私の要請に基づいてその実行を命令することに決定しました。フランス全体に均等に広がる30の地域で、最も正確な情報を提供できるコミュニティを選んでいます。これらのコミュニティの数え上げにより、1802年9月23日に合計2,037,615人の住民が得られました。これらのコミュニティでの1800年、1801年、1802年の出生報告は以下の通りです：

出生. 婚姻. 死亡.

男の子 11031 46037
女の子 105287 99443

人口と年間出産数の比は、これまで推定されていたよりも28T3I75oWb7r> が大きい。フランスの年間出産数にこの比率を乗じることで、この王国の人口を得ることができる。しかし、そのように決定された人口が真正な人口から与えられた限界を超えて乖離しない確率は何か？この問題を解決し、前節で得たデータを適用することで、私は以下の結論を得た：仏ランドの年間出産数を100万と仮定すると、その結果は28352845人の人口となり、この結果の誤差が半百万を超えない確率は私に約3万対1の賭けである。

前節で提示された声明によれば、男児の出産数に対する女児の出産数の比は22対21であり、結婚数に対する年間出産数の比は4対3である。

パリにおいて、男女の洗礼児数は約22対21という比率からわずかに異なります。1745年以降、出生登録で男女を区別するようになった時期から1809年末まで、この首都で洗礼を受けた男の子が393,386人、女の子が377,555人いました。これらの数値の比率はほぼ25対24に近いものでした。パリにおいては、男女の洗礼児数がほぼ等しいという結果から、特定の原因が平等な洗礼児数をもたらしていることが推測されます。

確率論の計算を適用すると、この原因の存在は賭け事における238対1の有利性を示し、これによりこの原因の調査が正当化されます。深く考えると、私はこの観察された差異が以下のように解釈できると考えました：田舎や地方で親たちは男の子を家に留める利点を見出し、結果としてパリの孤児院への送り込み数が男女の出生比率に基づいて少ないものでした。これはこの孤児院の登録簿によって証明されています。1745年から1809年末まで、孤児院に登録された男の子が163,499人、女の子が159,405人いました。最初の数値は後者の約-$$\frac{4094}{2755}$$だけ多いもので、少なくとも$$\frac{22}{21}$$の比率を上回るべきでした。これは、男児と女児の出生比率がパリでは22対21であり、孤児院への送り込みには注意が払われなかったという原因が存在することを確認しています。

これらの前段階の結果は、無限数の白球と黒球から各抽選で等しい確率を持つように混ぜられた urn からボールを引き出すような状況と比較できるものとして、出生数を推定することができるという前提に基づいています。しかし、異なる年ごとの同じ季節における変動が年間の比率に影響を与える可能性もあります。

不確定な増加の事象に関する考察 69

出生の男児数に対する女児数の比について、フランスの緯度局は毎年その国の人口の年間動向表を『年報』に発表しています。既に発表された表は1817年から始まっており、その年とその後の5年間に生まれた男児数は296,2361人、女児数は2,781,997人で、これにより約T$p$という比が得られます。各年の比率はこの平均結果から大きく変動することはありません；最小の比率は1822年に見られ、そのときの比率は約ffでした。最大の比率は1817年であり、そのときの比率は約ffでした。これらの比率は上記で見つかった比率である|fと比べて著しく異なります。この偏差を確率論の分析に適用し、男児と女児の出生比を urnから球を引くような比較仮説に適用すると、その結果はほとんど確実ではないことが分かります。したがって、この仮説は近似的には正確ですが、厳密には正確でないと示唆されます。

我々が述べた出生数においては、自然子供の男児数は200,494人、女児数は190,698人であり、性別比は約ffでした。これは見捨てられた子供の出生比と同様の結果で、自然子供のクラスにおける男女の出生が正規子供のクラスよりもより近い等価関係にあることを示唆しています。

フランス北部から南部までの気候差は男児と女児の出生比に著しく影響を与えるとは思えません。フランス全体の比率と同じく、最も南の30地区では約T|という比率が見られました。

パリとロンドンにおける男児と女児の出生比の恒常的な優越性は、それらが記録される以来、その傾向が変わることはありません。

一部の学者には、これは神の存在の証拠であり、無ければ、彼らは不規則な原因による事件の進行への恒常的な干渉が年に何度も娘の出生数が男児の出生数を超えることになると考えた。しかし、この証拠は、最終的因由という概念を検討する際に必要なデータがあるとその謎が解けるときにはしばしば見出されるような誤った使用の一例である。問題における恒常性は、男児の出生数に優越を与える規則的な原因による結果であり、それが大きな年間出生数に対する偶然性の異常をも含めて拡張される。この恒常性が長期的に維持されることの確率の調査は、過去の事象から将来の事象への確率へと推移する分析の一環である。1745年から1784年の間の観察された出生数を基にすると、パリで百年間にわたり男児の年間出生数が girl の出生数を常に上回るというベットはほぼ 4対 1 である。それゆえ、これがあなた半世紀にわたって実現したことに驚く理由はない。

別の例として、事象が増加するほど恒常的な比が示される傾向があることを考えてみよう。円形に並べられた多くの籠を想像し、それぞれの籠には非常に多数の白球と黒球が含まれているとする。これらの籠における白球に対する黒球の比率は最初から非常に異なるものであり、例えば一つの籠には白球のみで、別の籠には黒球のみであるなどとする。最初の籠から一球を引いて二つ目の籠に入れ、その後、

無限に増加する事象. 7*

次のように urn を揺らしてよく混ぜ、新しい球を他の球と混ぜるためには、二番目の urn から球を取り出し、それを三番目の urn に挿入し、その手順を最後の urn まで繰り返し、そこから最初の urn に戻って球を取り出す。この系列が継続的に再開されることにより、確率の分析はこれらの urn 内の白い球と黒い球の比率が最終的には同じになり、すべての urn に含まれる全ての白い球の総数に対する全ての黒い球の総数の比と同じであることを示す。従って、この規則的な変化により、これらの比率の原始的な不均一性は最終的に消失し、最も単純な秩序が生まれることになる。

もし、これらの urn の中に新しい urn を挿入し、その中の白い球と黒い球の比が前のものとは異なる場合、全体の urn に対する我々が先ほど示したような抽選を継続的に実行することにより、古い urn 内で確立された単純な秩序は最初に破壊され、白い球と黒い球の比率が不均一になる。しかし、少しずつこの不均一性は消失し、新しい秩序が生まれ、最終的には各 urn に含まれる白い球と黒い球の比率が等しくなる秩序となる。

このような結果は、自然のすべての組合せにおいて適用でき、その要素が動かされる定常的な力によって確立された規則的な作用の模式により、混沌の中心から優れた法則に支配されるシステムを生み出すことを示すことができる。

seeming に最も依存している現象は、大量に重ね合わされるとき、一定の比に近づく傾向を示し、そのような形で各々のこれらの比のどの程度でも望むほど小さな区間に想定した場合、観測結果の平均がこの区間に含まれる確率は、確定性からしかるべき量を超える量だけ異なるものと差異するまでに近づく。従って、大量の観測に対して確率の計算を適用することで、これらの比の存在を認識することが可能となる。しかし、原因を探求する前に、無駄な推測に導かれないようにするためにも、これらの比が偶然による異常であると見做す確率が許容されるほど高いことを確認する必要がある。判別関数の理論は、この確率を非常に単純な表現で与え、これは大量の観測から得られる結果が真実から一定未満の定数だけ異なる量によって変動する量の微分の積分により得られる。この積分は与えられた限界間で取り、同じ積分を正無限大と負無限大に適用したものを除いて割ることで、真実からの変動がこれらの限界間に収まっている確率を表す。これが大量の観測によって示される結果の確率の一般法則である。

\newpage

\chapter{自然哲学への確率計算の応用}
\label{ch:ix}

\begin{center}
\large\textit{天文学的観測データの誤差論・最小二乗法・太陽系の力学的安定性}
\end{center}
\vspace{4mm}

自然の現象はしばしば多くの奇妙な状況に包まれており、その影響を乱す多くの原因が混ざり合って作用します。そのため、これらの現象を認識することは非常に困難です。しかしながら、観測や経験を増やすことで、奇妙な効果は互いに相殺し、平均的な結果が現象とその要素を明らかにすることができます。観測の数が多いほど、それらが変動しないほど、その結果は真実に近づきます。

この最終条件を満たすために、観測方法を選択し、器具の精度を確保し、観察に細心の注意を払うことで達成されます。その後、確率論の理論を使って最も有利な平均結果や誤差が最小となる結果を決定します。

しかし、これだけでは不十分です。これらの結果の誤差が与えられた範囲内にある確率を評価する必要があります。これがなければ、得られた正確さの程度についての完全な知識を得ることはできません。

このような問題に対する適切な公式は、科学の方法に真の改善をもたらすものであり、それらをこの方法に加えることは重要です。これらの公式が必要とする解析は確率論の理論の中で最も精密かつ困難なものであり、私はこの理論に関する著作でそのような公式に至り、誤差の法則に関係なく観測自身によって与えられた量とその表現のみを含む特徴的な優位性を持つものを見つけ出しました。

それぞれの観測は、我々が決定したい要素の関数として解析的に表現される。そしてこれらの要素がほぼ分かっている場合、その関数はそれらの修正量に対する線形関数となる。観測自身に等しいようにこの関数を設定することで、条件方程式が形成される。類似した多くの方程式がある場合、これらをどのように組み合わせて、決定したい要素の数だけ最終的な方程式を得るか？それらの要素がまだ残っている誤差の確率法則は何か？これらの質問は確率論の理論によって明らかになる。

条件方程式を使用して最終的な方程式を作成する過程は、各方程式を未定因子に掛けてその積を結合することに相当する。最も少ない誤差の機会を与える因子の系を選択することが必要である。しかし、

確率論と自然哲学。75

もし要素の可能な誤差をそれぞれの確率に掛けて合計したとき、最も有利なシステムはこれらの積和が全て正に最小となるシステムである。正の誤差や負の誤差は損失として考慮されるべきである。この積和を形成し、極小条件によって因子の系統が決定され、適切に採用すべきまたは最も有利なシステムが特定される。したがって、各制約方程式中の第一要素の係数をそれぞれの制約方程式に掛けて全てを結合することで最初の最終方程式を作り出すことができる。同様にして第二要素の係数を使用して二番目の最終方程式を作り出し、そのような方法で多くの観測を集めた結果得られる要素と現象の法則は最も明確に開示される。

各要素がまだ残す可能性のある誤差の確率は、ハイパボリック対数が1となる数の誤差の平方を負の値として取り、その平方に等しいべき乗に比例し、このべき乗を誤差の符号を負としたものとすることにより、誤差の確率が減少するという定数係数で比例します。つまり、誤差が同じである場合でも、前の値が増加するとその確率は急速に減少しますから、真実に近い表現で言えば、要素はこの定数係数が大きいほどより大きな重みを持つことになります。そのため、私はこの定数係数を要素の重さまたは結果の重さと呼びます。

この重さはシステム内の因子の中で最大であり、最も有利な重さです。これが他のシステムよりもこのシステムに優越性を与える理由となります。質量が共通の重心で比較される場合の質量との驚くべき類似性から、同じ要素が多数の観測からなる異なるシステムによって与えられた場合、その最適なシステムはそれら全体の平均結果が各部分結果をその重さで乗じた総和であることを示します。さらに、結果の総合的な重さは各部分重さの総和であり、したがって、その総合的な平均結果の誤差の確率は、誤差を負としたものに等しいべき乗に比例し、そのべき乗はハイパボリック対数が1となる数のべき乗と等しく、そしてそのべき乗は各システムの重さの総和で乗じたものです。

実際には、それぞれの重さは各システムの誤差確率の法則に依存しており、ほとんど常にこの法則は未知です。しかし幸いにも、観測値から平均結果への変動の平方の合計を使用して、その因子を除去することができました。

大数の観測によって得られた結果について完全な知識を得るためには、各結果にその重さを書き加えることが望ましいです。解析はこの目的のために一般的かつ簡単な方法を提供します。

このような方法で指数関数を求めた場合、結果の誤差が特定の範囲内にある確率は、その範囲内でこの指数関数と誤差の微分の積の積分を計算し、それをその重さの平方根で乗じることにより求めることができます。

確率と自然哲学。77

同一の確率に対する結果の誤差は、その重さの平方根に反比例し、これによりそれぞれの正確度を比較することができる。すなわち、直径が1である円周に対する結果の比によって求められる。

この方法を成功させるためには、観測や経験の状況を変える必要がある。これは誤差の恒常的な原因を避けるためにである。観測は多数であり、決定すべき要素が多いほど多くなければならず、すなわち観測の平均結果の重さは観測回数が要素数で割られたとき増加する。さらにこれらの観測において要素が異なる経路をたどらなければならない。なぜなら、二つの要素の経路が完全に同じであれば、条件方程式における係数が比例することになり、これらの要素は単一の未知量となり、これらをこれらの観測から区別することは不可能となるためである。最後に、観測は正確でなければならず、これは最初の条件であり、結果の表現が観測値からの偏差の平方和で割られるので、この条件は結果の重さを著しく増加させる。これらの注意事項により、私たちは前記の方法を使用し、大量の観測から導き出された結果に対する信頼度の程度を測定することができる。

我々が JUST 提出した 方程式の条件、最終方程式を結論付ける法則は、観測誤差の平方和の最小化に等しいものである。各条件方程式は観測値とその誤差を代入することで正確になるため、それらから誤差の表現を導き出すことは容易であり、これらの表現の平方和が最小となる条件は要求する法則を与える。この法則は観測値が多いほどより正確であるが、観測値の数が少ない場合でも、そのような法則を使用することが自然であり、すべての場合において我々が求めたい修正を求めるために探る必要なく単純な方法を提供するからである。さらに、この法則は同じ星の様々な天文表の精度を比較するのに役立つ。これらの表は常に同一形式と仮定できるため、それらは時代、平均運動および引数の係数のみで異なるものとなる。もし一つの表に他の表に見当たらない係数が含まれている場合、その係数をゼロとして他の表に置き換えることと同じである。これらの表を良い観測値全体によって正誤修正すると、誤差の平方和が最小になる条件を満足する。多くの観測と比較して最もこの条件に近い表は優先されるべきである。

上記の原著英文を上記の方法は主に天文学において優れた成果をもたらすことができる。天文台の表は、観測と理論の精密さ、そして条件方程式を使用することによって達成された驚くほど正確な精度を持っている。これらの条件方程式により、多くの優れた観測が同一要素の修正に協力する。しかし、この修正後にもなお残る誤差の確率を決定することが必要であり、私が説明した方法はこれらの誤差の確率を認識するのに役立つ。

確率論と自然哲学．79

その方法を使って興味深い応用例をいくつか示すために、ブーヴァール氏が JUST 終えた巨大な著作『ジュピターとサターンの運動に関する精密な表』を利用してみることにする。彼はブラッドリーおよびその後に彼の観測を行った天文学者によるジュピターとサターンの対立点と四分点を最大限慎重に検討し、これらの運動要素と太陽（単位として1とする）に対する質量を修正した。彼の計算により、サターンの質量は太陽の 3512 分の 1 であることが判明した。確率論の私の公式を適用すると、この結果の誤差がその値の 1/10,000 未満である確率は n,ooo対1であり、つまり新しい観測が前回のものに加えられ、同じ方法で検討された場合、新たな結果がブーヴァール氏の結果から 1/10,000 の誤差を超えることはないという確率はほぼ同様である。この賢明な天文学者はまた、ジュピターの質量を太陽の 9 分の 1 と見出しており、私の確率論の方法はその結果が 1/1,000,000 の誤差であるという確率を 1,000,000 対 1 としている。

この方法は、地理測量作業においてもまた成功を収められるでしょう。地球表面の大円弧の長さを三角測量によって決定しますが、これは正確に測定された基線に基づいています。しかし、角度の測定にいかなる精密性があたえられようと、避けられない誤差は累積し、多くの三角形から導き出される弧の値が実際のものから著しく異なってしまうことがあります。

8o        論理学的確率論に関する哲学的小説。

真理は、その誤差が特定の限界内に収まっている確率を定めることが可能である限り、完全には認識できないものとします。測地学の結果の誤差は、各三角形の角度の誤差の関数です。私は引用された著作で、大量の部分的な誤差の確率分布が分かっている場合に、一つまたは複数の線形関数の値の確率を求める一般公式を与えました。これらの公式を使用して、測地学の結果の誤差が指定された限界内にある確率を求めることができます。これは、部分的な誤差の確率分布の性質に関係なく、より重要です。最も単純な確率分布も自然に存在する可能性が無限であることを考えると、常に極めて低い確率であるからです。

しかし、部分的な誤差の未知の確率分布は公式の中に不確定要素を導入し、これを数値化するためにはその要素を除去する必要があります。天文学における問題では、各観測が条件方程式を与え、要素を得る最可能値を代入した残差の平方和によりこの不確定要素を除去します。しかし、測地学の問題は類似の方程式を提供しないため、別の除去方法を求めなければなりません。各観測された三角形の角度合計が2直角に加えて球面余弦より大きい量の総和がその手段となります。したがって、条件方程式の残差の平方和をこれらの量の平方の総和で置き換えることができます。

確率論と自然哲学。81

そして、系列の測地作業の最終結果の誤差が特定された量を超えない確率を数値として定めることが可能です。しかし、各三角形の3つの角を観測した誤差の合計をどのように分配すべきか？確率論の解析は明らかに、各角はこの合計の三分の一を減らすべきであることを示します。これは、測地結果の重みが最大となる場合であり、同じ誤差がより不確実になるためです。すると、それぞれの三角形の3つの角を観測し、先ほど述べたように修正することには大きな利点があります。単なる常識もこの利点を示唆しますが、確率計算だけがそれを評価し、この修正によってそれが最大となることを明らかにします。

大いなる弧の値が、その一端点において測定された基線の一方から他方へ向かって別の基線を測ることによって確認される。そしてこれらの基線を結ぶ三角形鎖は非常に正確であると推測され、それに従って大いなる弧の値もほぼ正確なものとなる。この値は、測定された基線に従って三角形の角度を修正することで正しくなるが、これを行う方法は無数にある。その中で最も好まれるのは、測量結果が最大の重みを持つものであり、同一の誤差がより低い確率となるからである。確率論の解析は、それぞれの基線の測定値と確率の法則に基づいて最適な修正を直接求めることを提供し、多くの基線により生じるこれらの確率の法則は、その数が増えれば増えるほど急速に減少するものである。

多くの観測から得られた結果の誤差は、各観測の部分的誤差の線形関数である一般に、これらの関数の係数は問題の性質と結果を得るための過程によります。最も有利な過程は明らかに、同じ誤差が他のどの過程よりも結果においてより不確実であるようなものであることが明白です。確率の計算を自然哲学に適用するとは、これらの関数の値の確率を解析的に決定し、その不定係数を選択して、この確率の法則が最も速く減少することを選び出すことです。そして、質問のデータから導入される確率の部分的誤差の未知の法則の因子を取り除き、結果の誤差が特定の量を超えない確率を数値的に評価することが可能となります。これにより、多くの観測から得られた結果に関するすべてが必要なものが得られます。

他の考慮事項によっても非常に近似的な結果を得ることができます。例えば、同じ量について千以上の観測があるとします。これらの観測の算術平均は最も有利な方法で与えられる結果となります。しかし、その結果を選択するためには、

確率論と自然哲学。83

各部分値からの変動の総和がすべて正にとられたとき、その合計が最小となる条件は重要である。確かに、この条件を満たす結果を非常に近似的なものと考えるのは自然である。もし観測によって得られる値を大きさの順序に並べ替えれば、中央に位置する値は前述の条件を満たし、微分法により、無限数の観測の場合にはそれが真実と一致することが明らかになるが、最も有利な方法で与えられた結果は依然として優れている。

それによって我々は確率論が観測誤差の分配方法において何ら任意性を持たないことを理解する。それはこの分配に最適な公式を与え、可能な限り予想される誤差を最小限に抑えるものである。

確率の考慮は、観測誤差の中に包まれた天体運動の小さな不規則性を区別し、これらの運動で観察された異常の原因に戻るのに役立つ。

すべての観測を比較した結果、ティコ・ブラヘは太陽や惑星に適用された時間方程式とは異なる月に対しての時間方程式が必要であることを認識しました。同様に、多くの観測からメイアーサーが前進運動の係数が月に対して僅かに減少するべきであると認識しました。しかし、この減少はマッソンによって確認され、さらには増大したにもかかわらず、万有引力から生じるものではないように見えたため、多くの天文学者は計算においてこれを無視していました。

私たちは、目的のために選んだ月の観測の中から相当数を確率計算に提出し、ボウヴァール氏が私の依頼により検討することに同意したところ、その原因を調査することが強く示唆されたように思われました。私はすぐに、それは地球球体の楕円性であると気づきました。この要素はこれまで理論において無視されてきたもので、その影響は微細なものと見做されていましたが、微分方程式の逐次積分によって可視化されることが判明しました。

これらの項を特定するために特別な分析を行い、まず月の緯度における運動の不等関係を見つけ出しました。これは月の経度の正弦に比例し、以前の天文学者が疑念を抱かなかったものでした。この不等関係により、月の経度における別の不等関係が存在することを認識し、メイアーサーによって観測された月に対する前進運動方程式の減少と一致しました。

この減少量と緯度における先述した不等関係の係数は、地球の扁平率を決定するのに非常に適しています。私の研究結果をバールに伝え、彼が当時全ての良好な観測との比較によって月の表記を完成させようとしていたことに言及しました。私は彼にこれらの2つの量を特に注意深く定めるように依頼しました。

驚くべき一致性で、彼が見つけた値は地球の扁平率1/298.3とほぼ同じものを与え、これは測地度と重力 Pendulum から導かれた平均値とわずかな差異があります。

確率論と自然哲学。85

観測の誤差やこれらの測定における干渉要因による影響から、月の不規則性が完全に決定されているとは思われていなかった。

再び確率論を考慮することで、私は月の永久方程式の原因を見つけることができた。この星の現代観測と古代日食を比較すると、天文学者には月の運動における加速が示唆されたが、幾何学者、特にラグランジュは、この運動が経験した干渉からこの加速に依存する項を探すことに失敗し、その存在を否定していた。アラブ人によって観測された中間日食と古い現代の観測結果を慎重に検討すると、その存在が非常に高い確率で示唆されていることがわかった。それから私はこの視点から月の理論を取り上げ、月の永久方程式は太陽によるこの衛星に対する作用と、地球軌道の偏心率の永久変動が組み合わさった結果であることを認知した。これにより、月の軌道の節点と近地点の運動の永久方程式を発見することができた。これらの方程式は天文学者すら疑問視していなかったものだった。この理論が古代から現代までのすべての観測結果と非常に高い一致度を持つことにより、その信頼性は非常に高まっている。

確率の微分法は、木星と土星の大規模な不規則性の原因を同様に導き出した。ハーレーは現代の観測結果と古代の観測結果を比較し、木星の運動が加速し、土星の運動が遅延していることを発見した。これらの観測結果を調和させるために、彼は木星と土星の運動を1700年からの経過時間の平方に比例する符号反対の二つの世紀方程式に還元した。

オイラーとラグランジュは、これらの2つの惑星間の相互引力がこれらの運動に与えるべき変更を分析的に検討した。彼らはこれにより世紀方程式を見つけ出し、しかし結果は異なり、少なくともその中の一つは誤りであることが推測された。

私はこの天体力学における重要な問題を取り上げ直し、ハーレーの木星と土星の表において導入した世紀方程式を無視することに決定した。つまり、惑星の平均運動は不変であり、これはハーレーが導入した世紀方程式を無効化する。

これにより、これらの惑星の大規模な不規則性を説明するために残るものは、彗星による引力や多くの天文学者が実際的に採用した多くの惑星の運動に長期間の不規則性が生じること、または互いの作用によって影響を受け、それぞれに対して符号反対となる2つの惑星間の相互作用によって生じた長い周期の不規則性である。

この種の不規則性に関する定理を発見した私は、その定理によれば、木星が加速すると土星は遅延するという事実に合致し、これはハーレーがすでに指摘していたことと一致している。さらに、同じ定理から導かれる木星の加速は、土星の遅延に対して非常に近い比率で、ハーレーが提案した世紀方程式とほぼ同様である。

木星と土星の平均運動を考慮すると、私は容易に2つの時間の木星の倍数が土星の5倍から非常に小さな量だけ異なることを認識することができた。この不規則性の周期は、

$$
\frac{t_{\text{J}}}{t_{\text{S}}} = \frac{2 t_{\text{J}} - 5 t_{\text{S}}}{3 t_{\text{S}} - 2 t_{\text{J}}} \approx 1.6
$$

であることが示唆される。

確率論と自然哲学。87

この議論に対する差異は約九世紀分となるでしょう。実際、その係数は軌道の偏心率の立方のオーダーであることが知られていましたが、私はその係数が時間の非常に小さな乗数を含む引数において繰り返し積分によって除数として得られるものであることを知っており、これにより大きな値を持つ可能性があると判断しました。この不等関係の存在は当時私にとって非常に確実なことでした。その後の観察から得られた次の観察結果がその確率をさらに高めました。ティコ・ブラヘの観測時点において引数をゼロと仮定すると、ハレーは現代の観測値と古代の観測値との比較によって彼自身が示した変化を見つけ出すべきでした。一方で、現代の観測値同士の比較はラムブートから得られた結果に似た相反する変化を提供すべきでした。私はこの不等関係の確認のために必要となる長い且つ煩雑な計算を躊躇することなく着手しました。その計算結果は完全に予想通りであり、さらに他の多くの不等関係も確認され、これらは全て同じ観測値の精度で木星と土星の表を傾けたものです。

確率の計算によって、私は木星の最初の3つの衛星の平均運動に関する驚くべき法則を認識しました。この法則によれば、第一衛星の平均経度から第二衛星の3倍の平均経度を引いて第三衛星の2倍の平均経度を足したものは厳密に半円周と等しいです。これらの星の平均運動がその法則にどの程度近いかは、それらの発見以来の近似精度から極めて高い確率でその存在を示しています。

その後、その原因を探し求め、互いの作用を探求しました。その作用についての詳細な検討により、これらの衛星の平均運動がこの法則に特定の限界内まで近づいていた場合、それらの相互作用はそれを厳密に確立し維持したと結論付けました。

したがって、これらの3つの天体は、彗星などの突然な原因による木星周りの運動の変化を除き、前記の法則に従って永遠に空間で均衡を保つことになります。

従って、多くの観測結果から得られる自然の示唆を注意深く見ることはいかに重要かが明らかになります。既知の方法で説明できないとしても、その結果は重要なものです。宇宙系に関する問題に対する極めて困難な性質により、幾何学者たちは常に無視される量が著しい影響を持つ可能性があるという懸念を抱きつつ、近似法に recur せざるを得ませんでした。観測によってこの影響が警告されたとき、彼らは分析に戻り、その修正において常に観察された異常の原因を見つけ出し、そして法則を決定し、しばしばまだ示唆していない不等式を発見して予言しました。したがって、私は自然自身が普遍的重力の原理に基づく理論の解析的な完全性に協力していると言えるでしょう。そして、このことは私にとってこの素晴らしい原理の真実性の最も強い証拠の一つです。

私が JUST 考慮した例では、問題に対する分析的な解は原因の確率を確定に変換しましたが、ほとんどの場合、

確率論と自然哲学。89

この解は不可能であり、しかるにその確率をますます増加させるのみである。多くのかつ計算不能な変化が原因の作用から受け取られる中においても、これらの原因は奇妙な状況から与えられたものと同様に、常に観察された結果と適切な比を保ち、それらを識別し、存在を検証するために保持している。これらの比を決定し、多くの観測値と比較する。もしこれらが恒常的に満足しているなら、原因の確率は事実に対する確信と同じレベルまで高まることができる。原因からその結果までの比の調査は、直接的な問題解決よりも自然哲学において少なくとも同じくらい有用である。それはこれらの原因の現実性を検証するか、効果から法則を決定するかに関わらず、直接的解が不可能な多くの質問で使用でき、最も有利な形でその代わりとなる。私はここでは、自然の最も興味深い現象の一つである海流と潮の流れに対するこの方法の適用について議論したい。

プラインはこの現象について、その正確さで知られる説明を与えています。そこでは、古代の人々が満朔と新弦の時期には潮が最も高く、四分点の時期には最も低いことを観察していたことがわかります。また、月がperiastron（近地点）にあるときよりもaphelion（遠地点）にあるときに潮は高いこと、そして春分と秋分の時期よりも夏至と冬至の時期に潮が高いことも確認されています。彼らはこれらの観察結果から、この現象は太陽や月が海に対する作用によるものであると結論付けました。

ケプラーの『火星星に関する著作』の序文では、海水が月に向かう傾向があることを認めていますが、その法則を知らずしては、この問題について確かな理解を与えることができませんでした。ニュートンは、彼の大原則である普遍引力原理にこの考えを結びつけることで、この可能性を確実なものにしました。彼は潮の満ち引きを生み出すtractive力の正確な表現を与え、それぞれの瞬間に海がこれらの力を満たす均衡位置にあると仮定することで、潮の主な現象を説明しました。

しかし、ニュートンの理論から導かれた結果によると、太陽や月が大きな傾斜角を持つ場合、私たちの港では同じ日に2つの潮が非常に不等になることが予想されます。例えばブレストでは、満朔の時期における夕方の潮は朝の潮の約8倍に達するでしょう。これは観測結果と矛盾しています。

ニュートンの理論から得られたこの結果は、それぞれの瞬間に海が均衡位置にあるという仮定に対して成り立つものではありません。しかし、海の真の形状を調査することは大きな困難を伴います。幾何学者たちが流体の運動や偏微分計算に関する最近の発見を活用することで、私はこの問題に取り組み、地球全体を覆う海の運動の微分方程式を与えました。

このような自然に近い方法で調査を行うことで、私の結果が観測と一致していることを確認することができました。特に、同じ日に存在する夏至満朔の2つの潮汐間のわずかな差異について、その観察と一致していました。私はこれらの結果を見つけ出しました。

確率と自然哲学。91

海が常に同じ深さである場合、潮の高さは等しくなるでしょう。さらに、この深さに適した値を与えることで、港における潮の高さを観測に従って増加させることができました。しかしながら、これらの調査は、それらが示す地域的な状況の影響を全く満足させるものではなく、隣接する港でも見られる大きな差異を説明しませんでした。これらの状況を知る不可能性と、海の湾岸の不規則性、そして部分微分方程式の積分の困難性は、私は先に述べた方法で不足を補うことを余儀なくされました。その後、すべての海洋分子が影響する力の最大比を求め、それらの効果を我々の港で観察するために努力しました。この目的のために、私は次の原理を利用しました。これは多くの他の現象にも適用できるものです。

「物体のシステムにおいて、原始的な運動条件が運動に遭遇する抵抗によって消失した状態は、その運動を活性化する力と同様に周期的である。」

この原理と非常に小さな振動の共存性との組み合わせにより、私は港ごとに異なる地域的な状況の効果を含むかつて最小限の項数にまで簡略化された潮の高さの表現を見つけることができました；必要なのは、それを多くの観測結果と比較することだけです。

科学アカデミーからの招待により、前幾何学の初めにブレストで潮汐観測が行われ、その観測は連続して6年間続けられた。この港の位置はこのような観測にとって非常に有利であり、港は海と接するための運河によって接続され、運河の終端には広大な路地があり、港が建設されている。したがって、海面の不規則性は港内に僅かな程度まで及ぶ。また、ブレストでは潮汐が大きいため、風による偶然の変動は微弱である。これらの潮汐観測において、どれだけ観測を増やしても、非常に高い規則性が見られ、これは私に政府に対してこの港で新たな潮汐観測シリーズを、月の軌道の節点運動期間中に継続することを提案する動機を与えました。これが実現しました。観測は1806年6月1日から始まり、それ以来、中断なく毎日行われています。私は天文に関連するすべての部分について、私の分析と観測との比較に要した膨大な計算を要求されたボワール氏の不倦の熱意に感謝しています。約六千回の観測が1807年とその後の十五年間に実施され、これらの観測から、私の式は月の逆行から太陽までの位相変化、これらの星の傾斜、それらが地球からの距離、そしてそれぞれの要素における最大値と最小値での変動法則に関するすべての潮汐の多様性を驚くほど正確に表現していることが分かりました。この一致から得られる確率は、満潮と乾潮

海の潮汐は太陽と月による引力に起因し、その確実性は合理的な疑念を容認する余地を残さないほど高まっている。これは、この引力が宇宙の万有引力の法則から導かれるものであり、すべての天体現象によって証明されているためである。

月による海への影響は太陽よりも約二倍大きい。ニュートンとその successors は、この影響を立方距離で割られた項にのみ注目し、その後の項による効果は無視できるほど小さいと考えていた。しかし、確率論の計算は私に明らかにしたように、規則的な原因による最小限の効果も、適切な順序で多数の観測結果の中で現れる可能性があることを示している。この計算はまた、その確率を決定し、どの程度の観測が必要かを定める方法を提供する。マесьieur ブオワールによって議論された多くの観測結果を適用することで、ブレストでは満月時には月による海への影響が新月時よりも大きいことを認識した。また、南半球の月の位置の方が北半球のときよりも大きいことを確認した。これらの現象は、月からの距離の四乗で割られた項からしか得られないものである。

太陽と月の作用は大気を通過し、海に及ぼす影響を受け、それに類似した動きを経験するべきである。

これらの動きは、気圧計に周期的な振動を生み出します。解析によって私には明らかになりましたが、これらは私たちの地域ではほとんど感じられません。しかし、地元の状況が私たちの港での潮の高さに著しく影響を与えることを考慮に入れると、私は類似した状況が気圧の振動を著しくするかどうか再び調査しました。これを行うために、王立天文台で長年にわたって毎日行われている気象観測を利用しました。そこで、朝9時、正午、夕方3時、夜11時に気温と気圧の高さが観測されています。ボーヴォー卿は、1815年10月1日から1823年10月1日までの8年の間隔の観測結果を検討しようとしました。パリでの太陽潮汐を示す最も適切な方法でこれらの観測結果を配置したところ、対応する気圧の振動の範囲は約十八分の一ミリメートルでした。これが特に、結果の確率性を決定するための方法が必要であることを私たちは感じさせました。この方法が適用された場合、多くの観測結果を使用してもその結果の確実性は高くありませんし、十分な確率を得るには観測数を十倍に増やす必要があります。

私の潮汐理論の基礎となる原理は、規則的な法則に従って変動する原因が結合されるすべての偶然効果にも拡張できます。これらの原因の作用は、平均値において

確率論と自然哲学。         95

多くの効果の多様性は同じ法則に従い、その分析を通じて確率的に認識することができます。これらの効果が増えるにつれて、その多様性は確率でますます現れ、結果の数が無限大になると、それは確実性に近づきます。この定理は、既に開発した恒常原因による作用と類似しています。つまり、規則的な進行を持つ原因が特定の事件種に影響を及ぼすたびに、観測を増やし、その影響を示す最も適切な順序で並べることで、その影響を探求することができます。この影響が現れると、確率論の分析はその存在と強度の確率を決定します。例えば、日中の温度変化が大気圧を変えることから、バロメータの高さの観測値の増加は太陽熱の影響を示していると考えられます。実際には、赤道付近ではこの影響が最も大きい場所で、一日のうち9時頃に最大値と3時頃に最小値がある小さな日変動が長く認識されてきました。また、11時頃に二番目の最大値と4時頃に二番目の最小値があります。夜間の振幅は日中のものよりも小さいです、その範囲は約2ミリメートルです。私たちの気候の不確実性はこの変動を観測者から取り除きませんでしたが、それは熱帯よりはるかに少ないかもしれません。モン・ラモン氏は、プイ・ドゥ・ドーム地方の中心地クレマンタンで数年間にわたる精密な観測により、その変動を認識し、決定しました。彼はさらに冬の月よりも他の月の方が小さいことを発見しました。

パリでの気圧高さに対する太陽と月の引力の影響を評価するために行った多くの観測は、これらの日変動の一日間の変化を決定するのに役立ちました。午前9時に観測した高さと同日の午後3時に観測した高さを比較すると、この変動は非常に明確に現れ、その月ごとの平均値は1817年1月1日から1823年1月1日までの72ヶ月間常に正の値でした。これらの72ヶ月間の平均値はほぼ0.8ミリメートルで、クレマンタンよりも少し小さく、赤道よりもはるかに小さいです。私は午前9時から午後3時の一日間の変動の平均結果を確認し、11月、12月、1月では僅か0.5428ミリメートルであり、その後の3ヶ月では1.0563ミリメートルに上昇したことを認識しました。他の月には類似するものは見られませんでした。

これらの現象に対する確率計算を適用するために、まず一日間変動の異常性の確率法則を決定しました。それからこの現象の観測にそれを適用すると、それは規則的な原因がその結果を生み出すよりも30万対1以上の賭けであることがわかりました。私はその原因を特定することはありませんが、

確率論と自然哲学。97

私自身がその存在を述べることにとどまる。昼夜変動の周期は太陽日によって規律されていることが明白に示すところから、この変動は太陽による作用により生じていることが明らかである。太陽が大気に対して及ぼすtractive作用の極めて小ささは、太陽と月の合算した引力による効果の小ささから証明される。したがって、太陽の熱によって気圧の昼夜変動が生じることは明白であるが、その作用による気圧の高さや風への影響を微分積分学的に解析することは不可能である。磁針の昼夜変動も確実に太陽による作用の結果である。しかし、この星は気圧の昼夜変動における熱による作用と同様に、電気や磁気に対する影響により作用しているのか、それともこれらの影響の合計によるものなのか？異なる国で行われた長期的な観測によってこれを理解することができるだろう。

世界的システムにおいて最も驚くべき現象の一つは、惑星と衛星のすべての回転運動および革命運動が太陽の回転方向に一致し、また太陽の赤道平面と同じ平面で行われることである。このような驚くべき現象は偶然の結果ではない：これは特定の原因が全てのこれらの運動を決定したことを示している。

この原因が示唆される確率を得るためには、次の観察を行うことにより得られるものである。今日我々が知る惑星システムは、 Herschel に従ってウーニュスに六つの衛星を加えた場合、十一個の惑星と少なくとも十八個の衛星から成り立っている。太陽の回転運動、六つの惑星の回転運動、月の回転運動、木星の衛星の回転運動、土星のリングおよびその一つの衛星の回転運動が認識されている。これらの回転運動は革命運動と組み合わせて、四十三個の同じ方向に指向する運動を形成している。

しかし確率解析により、この配置が偶然の結果ではないという賭けは約4兆対1であることが判明し、これは歴史的事実に対する確率よりも高いものである。つまり、我々は少なくとも同様の信頼性で、原始的な原因が惑星運動を導いたと信じるべきである。特にこれらの多くの運動の多くが太陽の赤道平面に対して非常に小さな傾きを持っていることを考慮に入れると、その理由はより明白になるだろう。

太陽系におけるもう一つの興味深い現象は、惑星や衛星の軌道の偏心率が非常に小さいことと、彗星の軌道は非常に長楕円形であることです。この系の軌道には、大きな偏心率と小さな偏心率の中間的な影を見ることはできません。ここでもまた、規則的な原因の効果を認めざるを得ません。運命はすべての惑星や衛星の軌道がほとんど円形であるような形を与えることはありません；これらの天体の運動を決定した原因がそれらの軌道をほぼ円形にしたということです。

彗星の軌道の大偏心率がこの原因の存在から生じた結果であり、その原因が彗星の運動の方向に影響を与えたわけではなかったことを確認しなければなりません。なぜなら、逆行彗星と進行彗星はほぼ同じ数であることが見られ、すべての彗星の軌道に対する平均傾斜は太陽黄道面に対してほとんど直角に近い値を示しており、これらの天体が偶然で投げ出されたものであるべきであるからです。

どのような性質を持つ原因であれ、それが惑星の運動を生み出したかまたは導いたのであれば、それはすべての天体とそれらが離れている全ての距離を考慮に入れたはずです。したがって、この液体はこの恒星を中心に太陽の大気のように包囲していたことになります。

惑星の運動を考慮すると、太陽の大気は元々すべての惑星の軌道を超えて広がっていたことが考えられ、その後徐々に現在の限界まで縮小したという結論に至ります。

原始状態において、私たちは太陽を想像するとき、それは望遠鏡が我々に示すように、よりまたはそれほど brillant な中心核を取り巻く nebula で構成される nebulae のようなものに類似していました。表面で凝縮し、ある日には星へと変化することになるでしょう。この方法で形成されたすべての星を比喩的に考えると、その前段階の nebulosity 自体がさらに前の星々によってなされ、これらの星々では nebulous 物質がますます希薄になり、中心核はますます暗く、密度も低くなります。可能な限り遡ることを考えれば、nebula は希薄さに達し、その存在すらほとんど疑うことができなくなるでしょう。

実際、これこそが Herschel が特に彼の強力な望遠鏡を使用して観察した最初の nebula の状態であり、凝縮の進行を追跡しました。これらの段階は個々の nebula においては認識できず、我々にはその変化が著しく現れるのは数世紀後のことです。

100     確率論に関する哲学的考察

数世紀にわたる研究ではありますが、その全体を見渡すと、広大な森林の中で各年齢の木々が増加する様子を追跡することほど、それほど多くはありません。彼は、宇宙の異なる部分に広がり散らばった濁った物質から始め、これらの物質が一部で僅かに凝縮され、明るい星雲となる過程を見ました。他の星雲では、その中心核が周囲の星雲の量に比例して輝きます。各中心核がさらに凝縮することで、明るく近接する複数の星雲が形成されます。また、星雲の物質が均一に凝縮すると、惑星と呼ばれる星雲が生じます。そして、より大きな凝縮度はこれらの星雲を全て恒星へと変化させます。

この哲学的な見解に基づいて分類された星雲は、現存する恒星の過去の濁った状態とその将来的な恒星への転換を極めて確実に示唆しています。以下に述べる考察は、これらの類推から得られる証拠を補完します。

長い間、裸眼で見える特定の星々の配置が哲学的な観察者たちの注意を引いてきました。ミッチルは、例えば Pleiades 星団の星々が偶然だけでその狭い領域に留まっている確率は非常に低いと指摘し、これはこの星団群と天体が示す類似したグループが原始的な原因による結果であることを結論付けています。

確率論と自然哲学。101

あるいは、自然法則の一般化。これらのグループは、星雲がいくつかの核に凝縮する必然的な結果である；星雲物質が多様な核によって継続的に引き寄せられることが明白であり、時間とともにこれらはピレアスの星団と同等の星群を形成すべきである。

星雲が2つの核に凝縮すると、非常に近接した星が形成され、互いに回転するものも存在し、これらは Herschel が既に考慮したそれぞれの運動と等しい。さらに、天鹅座的第6星和其后的恒星，在 Bessel 刚刚认出它们具有特别显著且几乎相同的运动，这些恒星彼此接近并围绕共同质心运行的事实应无异议。

由此，从星雲物质的凝缩逐渐过渡到考虑太阳被广泛的大气层所包围的状态。正如之前通过太阳系现象的研究所见，这种考虑又回到了这一点。如此一来，一个如此显著的例子使太阳存在这一先前状态的概率接近于确定性。

但是，太阳大气如何决定了行星和卫星的旋转运动及公转运动呢？如果这些天体深入了大气层，其阻力会使它们坠入太阳；因此，很可能行星是在太阳大气的相继边界处形成的。随着冷却而收缩的大气层，在其赤道平面内应遗留下蒸汽带区，这些蒸汽带区由于分子间的相互吸引力而转变成了球状团块。同样地，卫星也是由各自行星的大气层形成的。

これらのグループは、星雲がいくつかの核に凝縮する必然的な結果である；星雲物質が多様な核によって継続的に引き寄せられることが明白であり、時間とともにこれらはピレアスの星団と同等の星群を形成すべきである。星雲が2つの核に凝縮すると、非常に近接した星が形成され、互いに回転するものも存在し、これらは Herschel が既に考慮したそれぞれの運動と等しい。さらに、天鹅座的第6星和其后的恒星，在 Bessel 刚刚认出它们具有特别显著且几乎相同的运动，这些恒星彼此接近并围绕共同质心运行的事实应无异议。由此，从星雲物质的凝缩逐渐过渡到考虑太阳被广泛的大气层所包围的状态。正如之前通过太阳系现象的研究所见，这种考虑又回到了这一点。如此一来，一个如此显著的例子使太阳存在这一先前状态的概率接近于确定性。しかし、太陽の大気はどのようにして惑星と衛星の回転運動および公転運動を決定したのでしょうか？これらの天体が大気層に深く penetrat した場合、その抵抗はそれらを太陽に落下させるでしょう；そして、おそらく惑星は太陽の大気の逐次的な境界で形成されたものと考えられます。冷え込むことで収縮する大気は、赤道平面内に蒸気帯区を残すべきであり、分子間の相互吸引力によりこれらの蒸気帯区は球状团块へと変化したはずです。同様に、衛星もそれぞれの惑星の大気から形成されたものです。

私は『宇宙体系の説明』において、この仮説を詳しく展開し、私にはこの体系が提示するすべての現象を満足させるものだと考えています。ここではその仮説について、太陽と惑星の回転速度が、それぞれの表面で連続的に気体の凝縮によって加速されるという点に着目し、それゆえ近接する天体の回転速度を上回るはずです。観測結果は実際、惑星や衛星に関してこの予想を裏付けており、さらには土星のリングについても同様で、その公転周期が0.438分である一方で、土星自体の自転周期は0.427分であることが確認されています。

この仮説によれば、彗星は惑星系とは無関係と見なされるべきです。彗星の形成を星雲の形成に付属させるとすれば、それは中心核を中心にした小さな星雲であり、太陽系から他の星系へと漂流していると考えられるでしょう。その形成は宇宙全体で広く散らばっている星雲物質の凝縮によって成り立つものです。彗星は、私たちの系において地球に類似した小天体である流星群のように、外人として見なされるべきです。これらの星が我々に見えるときには、その形態が星雲と非常に類似しているためしばしば星雲と混同されますが、運動を確認することや、その出現する部分の天空において閉じ込められているすべての星雲を知ることによって、それらを区別することができます。この仮説は彗星の頭尾が太陽に近づくにつれて著しく拡大し、またその尾が非常に希薄であるにもかかわらず、その存在自体が星の光を著しく減衰させないという現象をうまく説明します。

確率論と自然哲学。103

彗星の頭尾は太陽に近づくにつれてどれだけ広がるか、そしてこれらの尾が非常に希薄であるにもかかわらず、その存在自体が我々が見渡す星の光を著しく減衰させないという現象は、この仮説によってうまく説明されます。

太陽の引力が主導する空間領域、すなわちこの星の活動 sphere に小さな nebule が入ると、これらの nebule を引き寄せ、楕円軌道や双曲線軌道を描かせる。しかし、その速度はすべて方向で等しく可能であるため、これらは楕円のすべての向きと傾きにおいて同等に動くべきであり、これは観察結果と一致している。

彗星の軌道の大きな偏心率も前述の仮説から再び導かれる。実際、これらの軌道が楕円形である場合でも、その長軸は太陽の活動 sphere の半径以上であり、非常に延長されている。しかし、これらの軌道は双曲線形である可能性もある。そして、これらの双曲線の軸が太陽から地球までの平均距離に対して非常に大きいわけではなく、彗星がそれらを描く運動は感覚上双曲線形に見える。しかしながら、既知の百個以上の彗星の要素から見ても、一つも確実に双曲線形で動いているものがない。したがって、双曲線形の運動を生む確率は、その逆の確率に対して極めて稀であることが必要である。

彗星は非常に小さく、視認するためにはその近日点距離が極めて小さい必要があります。これまでのところ、この距離は地球軌道の直径よりも大きく2倍を超えたことは二度だけであり、ほとんどの場合、それは地球軌道の半径未満でした。彗星が太陽の活動領域に近づくためには、その進入時に太陽周囲空間において一定の範囲内で速度と方向を持つべきであると考えられます。

確率解析により、これらの範囲内で十分な楕円を描く機会に対する、パラボラと混同される可能性のある軌道を描く機会との比を決定したところ、その比は6000対1の賭けであることがわかりました。すなわち、太陽の活動領域に penetrates し観測される星雲が非常に長軸を持つ楕円または双曲線を描くことを示唆しています。この双曲線は観測部分においてパラボラと混同されやすいので、これまで双曲運動が認識されていないのは不思議ではありません。

惑星の引力、そしておそらくさらに太陽周囲空間の虚無中心による抵抗も、大軸が太陽活動領域の半径よりも小さい楕円軌道に多くの彗星の軌道を変更する可能性があります。これは太陽近接時に遭遇する蒸発によって見えなくなるまで、1200日周期で再出現し続ける彗星の場合にも当てはまります。

彗星が近日点に戻るたびに遭遇する蒸発により、その存在が無視される可能性があるため、この変化が1759年の彗星や1200日に一度しか現れない彗星においても生じていると考えられます。

さらに、確率の分析によって、その作用が存在すると信じられる一定の原因による影響を有機的な生物に及ぼすか、またはその存在自体を確認することができます。我々が自然の不顕微な因子を認識するために使用できるすべての装置の中で、最も感受性が高いのは神経であり、特にパラプロブ abilities があるときです。

確率論と自然科学. 105

特に感度を増加させる原因が存在します。これらの原因は、二つの異種金属の接触によって発展する微弱な電気を発見するために用いられました；これにより物理学者と化学者の研究領域が広がりました。

極端に敏感な神経の現象がいくつかの個体で起こることから、動物磁場の存在についての多様な意見が生まれました。これは通常の磁場に対する作用や太陽と月の影響による一部の神経疾患への影響、さらに金属や流れている水の近接性による印象についても言及されています。

これらの原因の作用は非常に微弱であると考えるのが自然であり、偶然の状況によって容易に妨げられる可能性があるため、その存在を否定するべきではありません。我々が全ての自然の因子とそれらの多様な作用形態を認識しているわけではないのであり、現時点での知識において説明できない現象を否定するのは非哲学者的であると言えます。

しかし、これらの現象を検討する際には、その承認がより困難になるほど、より慎重に注意を払うべきです。ここでは確率の計算が不可欠となり、それらが示唆する因子の存在に対する確率を他の場面で否定する理由よりも高い確度を得るためには、観察や経験をどれだけ増やす必要があるかを決定するために不可欠となります。

確率の計算は、各方法の利点と不都合を評価するのに重要な役割を果たすことができます。

speculative  sciences において 用いられる ものは ある でしょう。したがって 、病気 の 治療法 の 中で 最も 良好な 方式 を 確認する ためには 、患者を 同等の 数に 分割し 、各治療法を 全ての 条件が 正確に 類似した状況で 対象とする 患者に 施す ことが 必要です。最も有利な治療法 の優越性は 、試験対象者の数が 増えるにつれて ますます 明らかになります；そして 計算によって その優越性の 確率と 他の治療法に対する 相対的な 比較関係が 分かるでしょう。

\newpage

\chapter{道徳科学（人文・社会科学）への応用}
\label{ch:x}

\begin{center}
\large\textit{人間行動の統計法則と意思決定論への展開}
\end{center}
\vspace{4mm}

我々は、原因が不明または複雑でその結果を計算にかけることができない自然現象の法則を調査する際に確率解析の利点を見た。これは、倫理的な科学のほとんど全ての主題についての場合である。多くの予期せぬ要因、隠されたものもしくは認識不能なものも影響し、人間の制度に対する判断を事前に定めるのは不可能である。時間が経つにつれて発生する一連の出来事はこれらの結果を明らかにし、有害な結果を補正する方法を示す。多くの賢明な法律はこの観点から制定されたが、我々がその動機を保存しなかったため、多くの法は無用とされ廃止された。不快な経験が再び必要性を感じさせた場合でも、それらの法を再建立するべきであった。

各分野の公共行政において、様々な手段が生み出した結果について正確な記録を維持することが非常に重要である。政府による大規模な実験によって得られた多くの経験を反映しているからである。政治的・倫理的な問題にも適用すべきである。

107

108      確率論に関する哲学的小説。

観測と微分積分に基づく方法、自然科学において私達にそのような成功をもたらした方法について言及しよう。科学の進歩に対して、有用でしばしば危険な抵抗を提供するのではなく、知識の必然的な影響に対する最少限の抵抗を提供することを望む。しかし、我々は既に長年にわたって従来の制度や習慣を変更する際には極めて慎重であるべきである。

過去の経験から、これらの制度がもたらす困難さについてはよく知っているが、その変更によって生じる悪影響の程度については知らない。この知識不足においては、確率論はすべての変更を避けることを我々に導く。特に、道徳的世界と物理学的世界において突然の変化は常に大きな生命力の損失を伴うため、突然の変化を避けることが必要である。

既に確率論の計算は、道徳的な科学のいくつかの主題に対して成功裏に適用されている。ここではその主要な結果を提示する。

\newpage

\chapter{証言の確率について（Probability of Testimonies）}
\label{ch:xi}

\begin{center}
\large\textit{伝言連鎖による信憑性の減衰・複数証言の総合・歴史的批判}
\end{center}
\vspace{4mm}

私たちの多くの意見は証拠の確からしさに基づいており、そのため確率論に提出することは実際には重要です。確かに、証人の信頼性を評価する困難さと、彼らが証言する行為に関連する多くの状況により、しばしば不可能になりますが、いくつかの場合では、提案された問題と類似した問題の解法は誤った推論による誤りや危険から私たちを導き、守るための適切な近似解として見なすことができます。このような近似解はよく作られれば常に、最も説得力のある推論よりも優れています。それでは、いくつかの一般的なルールを制定してそのような近似を得ることを試みましょう。

千個の数字が入ったurnから一つの数が抽出されました。この抽選を見た証人が79という数字が抽出されたと報告しました。この数字が抽出される確率を求めます。経験により、この証人の信頼性は既に知られています。

$$
\log \text{no}
$$

「哲学的確率論の Essay」

証人が一度に十回中一回だけ嘘をついており、その証言の確率は \( \frac{1}{10} \) です。観察された事象は、証人が79が引かれたと主張することです。この事象は、証人が真実を述べる仮説または嘘をついた仮説から生じることがあります。観察された事象の確率を推定するためには、各仮説における事象の先験的な確率を求めることから始めなければなりません。

まず、証人が真実を述べる仮説では、証人が79を告げる確率は、79が引かれる確率そのものであり、これは \( \frac{1}{1000} \) です。この確率を証人の正直性の確率 \( \frac{9}{10} \) で掛けると、この仮説における観察された事象の確率は \( \frac{9}{10000} \) となります。

次に、証人が嘘をついた場合、79が引かれなかったということであり、その確率は \( \frac{999}{1000} \) です。しかし、証人が引かれた数字を選択する際には、残りの999個の中から79を選ぶ確率は \( \frac{1}{999} \) となります。この確率を先ほど求めた嘘をついた仮説における事象の確率 \( \frac{999}{1000} \) で掛けると、証人が79を告げる確率は \( \frac{1}{1000} \) となります。

この確率を嘘をついた仮説の先験的な確率 \( \frac{1}{10} \) で掛けると、この仮説における観察された事象の確率は \( \frac{1}{10000} \) となります。

最後に、各仮説における事象の確率を分母として、その和を分子とする分数を作ります。第六原理により、この分数が最初の仮説の確率となります。

この確率はT9ffである。つまり、証人の信憑性そのものである。これはまた、79という番号を引く確率と等しい。

もし証人が欺こうとするなら、未引いた番号の中から79を選ぶ利害関係がある場合、例えば、この番号に大きな賭け金をかけた上で、その抽選結果の発表が彼の信用度を高めるという判断があれば、彼がこの番号を選択する確率は最初のT9ffではなく、$p$や$p$など、その発表によって得られる信用度に関連して変化する。仮にそれが$p$である場合、虚偽の仮定下での観察された事象の確率を求めるには、この分数でT9ffを乗じる必要がある。さらに、信憑性の仮定下では、その結果を再びyで乗じることになる。これにより、第二の仮定下的に観察された事象の確率はTihhjTとなる。

したがって、最初の仮定、すなわち79という番号が引かれた確率は、前記の規則によってT9ffから減少する。証人が79という番号が引かれた場合にその抽選結果を発表する利害関係を考慮すると、この確率は著しく低下する。

しかし、同じ利害関係は、79という番号が引かれた場合に証人が真実を語る確率-を増加させる。しかしながら、この確率は1または$p$を超えることはない。したがって、79という番号が引かれた確率はTihhjTを超えない。

普通の感覚では、この利害関係が証人に与える不信感を示唆するが、計算はその影響を評価する。

証人が発表した番号の事前確率は、全可能な番号数で1に等しい。

urn中の数が証明によって真実性自体に変化し、証明により減少する可能性がある。例えば、urnの中に2つの数しか存在しない場合、その数をa prioriの確率 $p$で引く確立として与えられ、そして証人の信憑性がT%であるとすれば、この引く確立は減少する。実際、証人が虚偽に傾きやすいという事実は明らかであり、したがって、その証言は常にa prioriの確立が $p$ 以上の場合、その事実を証明する確率を減らすべきである。しかし、urnの中に3つの数が存在する場合、証人の信憑性が $p$ を超えると、引く確立は増加する。

例えば、urnの中に999個の黒い球と1個の白い球があり、そして1個の球を引いた後、その引いた球が白いという証人が現れた場合を考える。最初の仮定で観察された事象の確率は、前問と同じく -foooir- となる。しかし、証人が虚偽する場合の仮説では、白い球が引かれていないので、このケースの確立は TV である。これは虚偽の確立 $p$ に掛ける必要があり、これにより観察された事象に対する第二の仮説での確率は T||| となる。

前問ではこの確率は Tol7nr であったが、この大きな違いは次のように生じる。引いた球が黒い場合、虚偽を望む証人は未引かれた999個の球から白い球を引くことを宣言する選択肢を持っていないためである。

これらの2つの確率を分数として形成し、その分子がそれぞれの事象に対する確立である場合、

証言の確率についての考察。113

各仮説に対して、そしてそれらの共通分母はこれらの確率の和であるとき、最初の仮説と白い球を引く確率は \(i^{\prime}\)、第二の仮説と黒い球を引く確率は \(\frac{V_{\text{o}} V_{\text{s}}}{f}\) となる。この最後の確率は非常に高い確実性に近づき、百万個のボールが含まれる urn において一つだけが白い場合、白い球を引くことがさらに不思議なことになるほど、その値はより近づき、そして \(\frac{V_{\text{o}} V_{\text{s}}}{f}\) となる。したがって、行為がより不思議なものになるにつれて、虚偽の確率は増大する。

これまで証人が全く誤っていないと仮定してきましたが、しかし彼の誤りの可能性を考慮に入れると、その不思議な出来事はさらに不可能になります。すると、二つの仮説ではなく、以下の四つの仮説を持つことになります：証人が一切欺かずかつ完全に正しく記述した場合；証人が欺かず但し誤った場合；証人が欺いたが完全に正しく記述した場合；そして最後に、証人が欺き且つ誤った場合。それぞれの仮説において観察された事象の確率を事前に決定すると、第六原理により、その事実が虚偽である確率は、分母が分子と加算される、証人が完全に正しく記述し且つ誤る確率または欺くが正しく記述しない確率の和によって与えられる分数の分子が黒い球の数とその積であることがわかります。

すべての事実が真実であるか、または彼は欺き、また誤りをしているか。このことから、 urn 内に黒い球が非常に多い場合、白い球を引くことは異常であり、その証言が偽である確率は最も近づいてくるほど確実性に近づくことがわかります。

この結論をすべての異常な行為に適用すると、証人の誤りまたは虚偽の確率は、証明された事実がより異常であるほど大きくなることが分かります。ある著者は、異常な事実と普通の事実を見分ける視点において完全に類似しているという理由から、両方の事実を主張するとき、同じ動機により証人に等しい信用を与えるべきだと主張しています。単純な感覚はそのような奇妙な主張を拒否しますが、確率論の計算は普通の感覚と一致しつつ、異常な事実に関する証言の最大の不確実性を評価します。

これらの著者は、同等に信頼できる2人の証人について強調し、第1の証人は15日前に死んだ個人を見たと主張し、第2の証人はその個人が昨日元気な状態で見たと主張するという仮定を立てます。これらの事実のいずれも不確実性を提供しません。個々人の保留はそれらの組み合わせの結果ですが、証言自体はこの結果に直接到達させません。しかし、これらの証言に対する信用は、その組み合わせの結果が異常であるという事実によって減少するべきではありません。

ただし、証言の組み合わせから導かれる結論が不可能な場合、

証言に関する確率について
115

不可能な結論は必ずしも偽であるとは限らないが、不可能な結論の限界は驚くべき結論の限界であり、誤りは不確かな結論の限界である。不可能な結論の場合、証言の価値がゼロになるならば、驚くべき結論の場合にはその価値は非常に大きく減少するべきである。これは確率の計算によっても確認される。

これを明瞭にするために、まず二つの壺を考える。AとBという二つで、最初の壺には百万個の白い球が入っており、第二の壺には百万個の黒い球が入っている。一つの球を引き出し、それをもう一つの壺に戻し、その後もう一度別の球を引く。この二つの証人がそれぞれ最初と最後の引きの球が白いという証言をするが、どの壺から引き出されたかは明示しない。それぞれの証言だけでは不確かなものだが、その証言が示す事実の確率は証人の正直さそのものであることが容易に理解できる。しかし、両方の証言を組み合わせると、最初の引きでAの壺から白い球を取り出し、それをBの壺に戻し、その後再び取り出した球が白いという非常に驚くべき事態が生じる。この第二の壺には百万個の中の一粒だけが白いので、その中から白い球を引く確率は$\frac{1}{1000001}$である。

二つの証人の証言によって示された事実の確率がどのように減少するかを決定するために、観察された事象はそれぞれの証人が見た球が白いという主張であることを注意深く見つめる。各証人の正直さを$T$と表すと、これは

$$
\frac{1}{1000001}
$$

となります。

現在の場合において、証人が完全に嘘をつかずかつ誤解をしていないとき、あるいは嘘をついてかつ誤解をしているときには、以下の4つの仮説を形成することができます：

1. 第一の目撃者と第二の目撃者が真実を述べている場合。最初に、白い球が urn A から抽出され、このイベントの確率は $\frac{1}{2}$ です。なぜなら、最初の抽選で球がどちらかの urn から抽出される可能性があるからです。したがって、その白い球が urn B に入れられ、二度目の抽選で再び現れたというイベントの確率は、目撃者が真実を述べる確率 $\frac{1}{2} \times \frac{1}{2}$ の積に相当します。このイベントが観察された確率は、$\left(\frac{1}{2}\right) \times \left(\frac{1}{2}\right) = \frac{1}{4}$ となります。

2. 第一の目撃者が真実を述べている一方で第二の目撃者が嘘をついている場合、彼が誤解をしていないか、または嘘をつくが誤解していないかに関係なく。最初に白い球が urn A から抽出され、このイベントの確率は $\frac{1}{2}$ です。その後、その白い球が urn B に入れられ、黒い球が抽出されたというイベントの確率は $\frac{1}{3}$ です。したがって、複合的なイベントの確率は $\left(\frac{1}{2}\right) \times \left(\frac{1}{3}\right) = \frac{1}{6}$ となります。この複合的なイベントが観察された確率は、第一の目撃者が真実を述べる確率 $\frac{1}{2}$ と第二の目撃者が嘘をついている確率 $\frac{1}{2}$ の積に相当します。したがって、観察されたイベントの確率は $\left(\frac{1}{6}\right) \times \left(\frac{1}{2} \times \frac{1}{2}\right) = \frac{1}{24}$ となります。

3. 第一の目撃者が真実を述べていない一方で第二の目撃者がそれを述べている場合。最初に、urn B から黒い球が抽出され、その後のイベントの確率は $\frac{1}{2} \times \frac{1}{3}$ の積となります。この複合的なイベントが観察された確率は、第一の目撃者が嘘をついている確率 $\frac{1}{2}$ と第二の目撃者が真実を述べる確率 $\frac{1}{2}$ の積に相当します。したがって、観察されたイベントの確率は $\left(\frac{1}{6}\right) \times \left(\frac{1}{2} \times \frac{1}{2}\right) = \frac{1}{24}$ となります。

これらの仮説についての考察は、証言の確率に関する問題を理解する上で重要です。

urn Aに白い球が入れられた状態で、このurnから白い球が引かれた。これらの事件の最初の確率は$p$であり、二番目の確率はT$p$\#$S$Tである。それらの複合事象の確率はi-fとなる。それを最初の証人の真実性を持たない確率13jrと、第二の証人がそのことを告げた確率yVを掛け算すると、観測された事象に対するこの仮説での確率はsiiHfTmHtaとなる。

4番目。最後に、どちらの証人も真実性を持たない場合。最初の抽選でurn Bから黒い球が引かれた後、それをurn Aに入れると、二度目の抽選で再出現する確率はaooooo?-である。それをそれぞれの証人が真実性を持たない確率-faとy1-を掛け算すると、観測された事象に対するこの仮説での確率は200000200となる。

今、二つの証人の告発、つまり各抽選で白い球が引かれたという事実の確率を得るためには、最初の仮説に対応する確率を四つの仮説に対する確率の合計で割らなければならない。そして、その確率はy-g-ooooasとなる非常に小さな分数である。

二つの証人がそれぞれurn AとBから白い球が引かれたことを主張し、第二のurn A'およびB'も同様に類似した urn である場合、二つの証人の告発に対する確率はその証言の確率の積、すなわち$y \cdot v$となる。それにより、

第4章. 統計的確率論の哲学的考察.

少なくとも百八万千分の一万倍以上大きい。これを見れば、最初の場合において、最初に白い玉を引いた後で二度目の引きで同じ白い玉が再び引かれるという証言がいかに稀有な結論であるかがわかる。

我々は、百個のサイコロを投げて全て同一面に落ちたと主張する男の証言には信用しないだろう。もし我々自身がその様子を見ていたら、他の目も含めて慎重に状況を検討し、誤覚や詐欺がないことを確認した上で初めてそれを信じるだろう。しかし、その検討後でも極度の不確実性にもかかわらずそれを認めることなく逡巡しないだろう；そして誰も視覚の法則を否定する説明に訴えることは好ましくないだろう。我々はそれから自然の法則の恒常性の確率がこの事象が全く起こっていないという確率よりも大きいと結論付けるべきである。これは、我々が不動の事実として受け入れている多くの歴史的事実の確率よりも大きい。これにより、自然法則の中断を認めるために必要な証言の重みが判断できるし、そのようなケースでは通常の批判の規則を適用することは適切でないこともわかる。

しかし、学識のある人々の信頼性は減少する一方で、未開化の人々の信頼性はしばしば増加することがある。彼らは常に驚異的なことに飢えているからである。

証言に関する確率 119

したがって、教育された人々の信頼性を減らすものは、未開化の人々の信頼性をしばしば増加させることがある。

あることに驚異的なものがある。その不確実性は、いかなる証拠にも匹敵しないことがある。しかし、主導的意見の影響によって、その驚異性は弱まり、証言の確度に劣るまでに至ることがある。そしてこの意見が変化すると、16世紀に生まれたこの荒唐無稽な主張は、その後の時代には、一般意見によるより明哲なる心の影響を示す新たな証拠として提供されるのみとなる。

ルイ14世の時代における2人の偉大な人物、ラシネとパスカルは、これに典型的な例である。ラシネ、この人間の心を絶妙に描き、かつて生きた最高峰の詩人であるが、パスカルの甥女ペリエール・パスカルの奇跡的な回復を報告し、そのことを神聖なるものと報告するという態度は、痛々しいほどである。そしてパスカル自身が、この奇跡が宗教に必要であり、モンクたちの教義を正当化するために必要な理由を探求する論理も、痛々しく感じられる。

ペリエール・パスカルは3年半にわたって涙管 Fistula という病気で苦しんでいた。彼女は救世主の冠に装飾された棘の一つとされる遺物を自分の患った目につけることで、瞬時に回復する信仰を持っていた。数日後、医師たちと外科医たちは彼女の回復を確認し、自然と治療薬がその回復に関与していないことを宣言した。この出来事は1656年に起こり、パリ全城がポート・ロワールの修道院に駆けつけたという。人々の数は日に日に増加し、神自身もこの教会で行われる奇跡の数々によって、人々の信仰を認めるよう喜んでいるかのように見えた。

当時の時代において、奇跡と呪術がまだ不確実性を伴うものとして認識されていなかったため、そのような自然の異常な現象は、それ以外の説明が得られない限り、奇跡や呪術に帰属されることが了承された。

このような結果を観察する方法は、ルイ14世の時代における最も優れた作品で見ることができます；すなわち、哲学者のロックによる「人間理解論」においても見られます。彼は証言される信ずる程度について述べているところにこう言っています。「一般的な経験と日常的な事象が人々の心に影響を与え、何かを信じるために信用を与えるか否かを判断する上で正当な影響を持つことは認めますが、それでもある一つのケースにおいては、その出来事の不思議さが証人の公正な証言に対する信頼性を損なうことはありません。なぜなら、自然の流れを変える力を持ち、その目的に向かって行動する者が存在する場合、そのような超自然的な出来事が、通常の観察を超えるほど不思議であるほど、そのような状況下ではより信頼性が高まるからです。」

このような証言の確率に関する真実の原理は、主に理由に基づいて進歩を遂げた哲学者たちによって誤解されており、私はこの重要な主題について計算法の結果を詳細に提示する必要があると感じました。

ここで自然に生じるのが、パスカルの有名な議論、それをケイグという英国の数学者が幾何学的形式で展開したものです。証人たちは神からそのようなことを知っていると主張し、ある特定の事柄に従うことで無限の幸せな生活を得ることを約束されていると述べています。証明の確率があまりにも弱くても、それが無限小であるとは限らない限り、それらが従う人々にとっての利得は無限大であり、それはその確率と無限の善の積に等しいからです。それゆえ、この利得を得るために自己を確保すべき迷いはありません。

証言の確率論についての考察については、これ以上詳しく述べることにします。

この議論は、神の名前で証言する無限に続く幸運な人生を約束した者たちに基づいています。彼らが極端にその約束を無制限の限界を超えて広げているので、それを好ましく思わない理性から、これらの証人にはその約束を明確に定めることが必要です。また、微分積分学は、この過度の言及自体が証人の証言の確率を無限小またはゼロまで薄めることを教えてくれます。

実際、これは、大量の番号が詰まった壺から最高の番号が引かれた後に、その番号が引かれたことを告げる最大の利益を得るために証人が嘘をつこうとするケースと似ています。既に、この利益が証人の証言をどれだけ薄めるかについて説明しました。

証人が嘘をついた場合に最高の番号を選ぶ確率を$p$と評価したとき、微分積分学はその告白の確率を、虚偽の事前の考慮から独立して考えた確率の半分の製品に加えて数値の数によって与えられる分数よりも小さいものとして与えます。この証明とパスカルの議論を比較するには、壺の中の数字で全ての可能な幸運な人生を表現し、これらの数字の数が無限大であることを示すだけで十分です。そして、証人が嘘をついた場合にその虚偽を認めるために最も大きな利益があることに注意すると、彼らの証言の確率は無限小になります。

この証言の確率を無限に続く幸運な人生の数で掛けると、積が示すパスカルの議論から生じる利得が消失し、パスカルの議論自体が破綻します。

現在、いくつかの証言が既存の事実についての確率を考慮しましょう。まず、我々の考えを明確にするために、百個の数字が入ったurnから一つだけ数字が引かれたという事実を考えましょう。この引き取りを証する2人の証人が数字2が引かれていると報告し、これらの証言全体の結果的な確率を求めることにします。証人たちは真実を語るか、嘘をついているかの二つの仮説を形成することができます。

最初の仮説では、数字2が引き出され、このイベントの確率は$\frac{1}{100}$です。これは証人の信憑性（信ぴょう性）の積に乘じる必要があります。証人の信憑性をそれぞれ$T_{97}$と$T_1$と仮定すると、この仮説下での観察されたイベントの確率は$T_{97} \cdot T_1 \cdot \frac{1}{100}$となります。

二番目の仮説では、数字2が引き出されない確率は$\frac{99}{100}$です。しかし、証人の一致した報告は、彼らが嘘をつこうとする際に引く数を99つの未抽出の数字から2にすることを意味します。証人が秘密の合意を持っていない場合、この選択の確率は$\left(\frac{1}{99}\right)^2$となります。したがって、これらの二つの確率を乗じて、そして証人が嘘をついている確率である$\frac{99}{100}$と$\frac{99}{100}$の積を乗じることになります。これにより、第二の仮説下での観察されたイベントの確率は$\left(\frac{99}{100}\right)^2 \cdot \frac{99}{100} = \frac{9801}{10000}$となります。

現在の事実が証明されるか、つまり数字2が引かれることの確率は、最初の仮説の確率を二つの仮説の確率の合計で割ったものです。この確率は$\frac{\frac{1}{100} \cdot T_{97} \cdot T_1}{\frac{1}{100} \cdot T_{97} \cdot T_1 + \left(\frac{99}{100}\right)^2 \cdot \frac{99}{100}} = \frac{T_{97} \cdot T_1}{T_{97} \cdot T_1 + 9801/10000}$となります。そして、数字2が引かれなかったことと証人が嘘をついた確率は$\left(\frac{99}{100}\right)^2 = \frac{9801}{10000}$となります。

CONCERNING THE PROBABILITIES OF TESTIMONIES. 123

もしurnの中にしか数字Iと2が含まれていれば、同様の方法で数f $p$を2の引き出しなどの確率として見つけ出し、従って証人の虚偽の確率はとなり、前記の確率よりも少なくとも九十四倍大きいことがわかる。このことから、証人によって証明される事実がより低い確率を持つとき、証人の虚偽の確率がどれほど減少するかが理解できる。確かに、そのような場合、証人が欺くときの一致は困難になり、少なくとも秘密の合意がない限り、我々はここではそれを仮定していない。

前記の場合でurnの中にしか二つの数字だけが含まれていたとき、事実の事前の確率はであり、証言結果の確率は証人の正直性の積をその積に加えたものと虚偽のそれぞれの確率の積との和で割ったものである。

現在、私たちは伝統的な証人の鎖によって伝えられる事実の確率が時間にどのように影響されるかを検討する必要があります。これは明らかに、その鎖が延びるほどこの確率は減少するべきであるということです。もし事実は自身から確率を持たないものであれば、例えば無限個の数字が入っているurnから数を引くような場合、その事実が証人の証言によって獲得する確率は証人の正直性の連続的な積に比例して減少します。もし事実が自身から確率を持っているものであれば、例えばurnから2という数を引くという事実で、 urnには一つだけ数字が引かれておりそれが確定している場合、伝統的な鎖がこの確率に加えるものは、最初の因子がurn内の数字の個数から1を引いた数と同数の比であり、他の各因子はそれぞれの証人の正直性からその証人の虚偽の確率との比を減算したものです。つまり、事実の限界的な確率は、事実が証言以前に先ず考えられた確率、すなわちurn内の数字の個数で割った1と等しい確率となります。

時間の作用は歴史的事実の確率を絶えず弱め、最も耐久性のある記念物も変えるように働きます。証言やその支持する記念物を増やすことによって、この確率を減少させることが可能です。印刷はこれに非常に有用な手段を提供しますが、不幸にも古代の人々には知られていませんでした。物理的・精神的な革命により地球の表面は常に揺れ動くため、時間の必然的な効果と結びつき、千百年後にその確率が疑問視されるでしょう。

歴史的事実として今日最も確かなものとされる諸事実についての確率論。125

数年間、歴史的事実が今日もなお最も確かなものとされている。

Craigは、キリスト教の宗教証明の逐次的な衰弱を微分積分に適用しようと試みました；彼は世界が確からしさを失ったときには終わりるべきであるとする仮定に基づき、その結果としてこれは彼が書いている現在から1454年後になるべきだと結論付けました。しかし、彼の分析は彼が月の持続期間についての仮説ほど奇妙なものではありませんでした。

\newpage

\chapter{議会の選択と決定について（Selections of Assemblies）}
\label{ch:xii}

\begin{center}
\large\textit{多数決原理・投票制度の数理比較（ボルダ方式とコンドルセ基準）}
\end{center}
\vspace{4mm}

議会の決定の確率は、その構成員の投票数、知識と公正さによって左右されます。多くの情熱や個別の利益がしばしば影響を及ぼすため、この確率を計算する事は不可能です。しかし、単純な常識から導かれるいくつかの一般的な結果は、計算によって確認されています。例えば、議会が提出された主題について十分に情報がない場合、その主題が精密な考慮が必要である場合、または既存の偏見と一致しない真実がある場合、各投票者が誤る確率が1対1より高いと予想されるならば、大多数による決定はおそらく間違っているでしょう。そして、議会が大きいほど、その懸念はより正当になります。

公共の事務においては、多くの人々が理解できる範囲にある主題について議会が判断することが重要です。また、情報が広く普及され、理由と経験に基づいた善行によって、決定を下す者たちに誤った考えや無知による偏見から警戒されることが重要です。

学者たちはしばしば、最初の概念はしばしば欺き、真実は必ずしも確実ではないという事実を指摘してきました。

議会の欲求を多様な意見の中から理解し定義することは難しいですが、この問題に関していくつかの規則を与えるために、最も一般的な2つのケースを考えることで試みます。すなわち、複数の候補者の選挙と、同一主題に関する複数の提案の選択です。

ある集団が、同じ種類の職位に対する一人または複数の候補者の中から選ぶとき、最も単純な方法は、各投票者が優先順位を付けて候補者の名前を票に記入することである。彼が誠実に評価した場合、これらの票を検討することで選挙結果を得ることができ、これにより候補者同士を比較し、新たな選挙でより良い結果が出ることはない。

現在の問題は、これらの票によって候補者の間にどのように優先順位を決定するかである。各投票者が無限の球体を含むurnを与えられ、彼がどの程度の優遇度を与えるかを選択できると想像してみよう。また、彼が各候補者に対するその優遇度に比例した数の球体を引き出し、その数を票の側方に候補者の名前に記入するという手順を考える。明らかに、すべての票で各候補者に対し記載された数値の合計を求めることにより、最も高い合計を持つ候補者が集団が好む候補者であり、そして候補者の優先順位はその合計の順序となることがわかる。

しかし、これらの票には各投票者がどの程度の球体を与えたかを明示していない。それらは単に最初の候補者が次の候補者よりも多く、次がさらに多いという順序だけを示している。したがって、特定の票に対して一定数の球体を仮定した場合、その前の条件を満たすすべての下位の組み合わせが等しく可能であり、各組み合わせが与える数値の合計を全て加算し、全体の組み合わせ数で割ることにより、各候補者に対する球体の数を得ることができる。

非常に単純な解析によれば、最後の名前の隣に記入すべき数値、次に最後の名前、そしてその他の名前に対しても記入すべき数値は等差数列1, 2, 3, 等の項に比例していることがわかる。したがって、各票に対してこの数列の項を記入し、各候補者に対する項の合計を加算することで、その合計の大きさによって候補者の優先順位が示される。

このような選挙方法は確率論が提案するものである。もちろん、各投票者が自身の票に彼が与える優遇度の順序で候補者の名前を記入することが望ましいが、個々の利害や多くの奇妙な評価基準によりその順序が影響を受け、時には最後の位に候補者が置かれることがある。

会議における選択と決定。129

最も恐ろしいのは、自分が好む候補者に対して与える過度の優位性であり、これは中程度の才能を持つ候補者の利益を大きくするからである。また、経験は、この選挙方法を採用した社会でその使用を見舍てることをもたらした。

絶対的多数による投票により行われる選挙は、多数が拒否するすべての候補者を受け入れない確実性と、会議の希望を最も頻繁に表現する利点を結びつける。この方法は、候補者が二人だけである場合においては、前記の方法と常に一致する。確かに、これは会議に対して選挙が永远に続く可能性のある不便を与える危険性がある。しかし、経験は、この不便さは全く問題でないことを示し、そして選挙を終結させたいという一般的な希望は多数の投票を集めることが早く行われるということが判明した。

同じ対象に対するいくつかの提案を選択する場合における選択は、候補者のいくつかを選ぶ場合と似たような規則に従うべきであるように思われる。しかし、二つの状況間にはこの違いが存在し、すなわち、候補者の才能はその競争相手の才能を排しないが、提案が相反するものである場合には、一つの真実は他のすべての真実を排除することになる。それでは、この問題についてどのように見ればよいのか見てみよう。

各投票者に無限の球が入ったurnを与えることとして、彼はそれぞれの候補に対する自己 atributoする確率に応じてこれらの球を分配することとする。確実性を表す球の総数は明らかであり、仮定によれば少なくとも一つの候補が真であることが保証されているので、投票者は最終的にその数を各候補に分配するだろう。問題は、最初の候補よりも多くの球が第二の候補に、第二の候補よりも多くの球が第三の候補にと、それぞれの候補に対して球の数が増加するような組み合わせを求めることに帰着される。各候補に対する球の数をすべての組み合わせで合計し、その合計を組み合わせの数で割ることにより、商は特定の投票用紙に対する各候補への球の数となる。

分析によって得られる結果として、最後の候補からこれらの商は以下の量に比例する。まず、候補の数で1を除く量；次に、その前の量に1を加えたもので候補の数から1を引いた量；第三に、その第二の量に1を加えたもので候補の数から2を引いた量、そして他のものは同様である。これらの量はそれぞれの投票用紙の側方に記載され、異なる投票用紙における各候補に対する相対的な量を合計すると、その大きさが議会がこれらの候補に対して与える好意度の順序を示す総和となる。

全体会議の更新方法について言及したい。これらの会議が一定期間ごとに完全に更新されるべきであるとする場合、その更新は一度に行われるべきか、それともこれらの年数間に分散されるべきか。

後者の方法によれば、会議はその更新時期における多様な意見の影響下で形成されることになり、当時の意見がおそらくすべてこれらの意見の平均となるでしょう。したがって、会議は選挙の期間を代表する全土にまで拡張することによって得られる同じ利益を得ることになります。

現在、経験が明確に教えていることは、選挙は常に支配的な意見により最大限に導かれるということです。そのため、一部の更新を通じてこれらの意見を他の意見で和らげることはいかに有用であるかを感じることができます。

\newpage

\chapter{裁判における判決の確率について}
\label{ch:xiii}

\begin{center}
\large\textit{陪審員制度・冤罪確率の最小化と法的証明度}
\end{center}
\vspace{4mm}

解析は、単なる常識が教えているように、判決の正しさは裁判官がより多く、より賢明であるほど、その確率も高まるという判断を確認します。従って、上告 tribunal がこれらの二つの条件を満たすことが重要です。最初の instance tribunal が直接関与している被告に対して、上告 tribunal に第一審判決という既に確度のある判決を提供し、しばしばその判決と一致することがあります。これは妥協または請求を取り下げることで行われることがあります。

しかし、争点の不確定性と重要性が当事者に上告 tribunal の利用を決定させた場合、彼はより高い正義を得る確率により財産上の安全が確保され、新たな手続きによる手間と費用の補償を見つけるべきです。この点は、地域 tribunal 互いの相互上告制度においては考慮されておらず、その結果、市民の利益に非常に有害でした。おそらく、判決を無効にするために上告 tribunal の多数決が少なくとも2票であることを求めることが適切で、かつ確率論と一致しているかもしれません。

上告 tribunal が偶数の裁判官から成る場合、投票が均衡した場合でも、判決が維持されることになります。私は特に刑事事件の判決について検討します。

疑罪から有罪を断定するためには、裁判官が被疑者の犯罪に対する確固たる証拠を持つことが必須であることは間違いない。しかし、道徳的な証明は常にただの確率に過ぎず、経験は刑事判決がいかなるものであれ、誤りを犯す可能性があることをはっきりと示している。これらの誤りを修正する不可能性は、死刑を禁じようとする哲学者たちにとって最強の論拠である。もし数学的な証拠が必要なほど疑罪から有罪を断定する根拠がなければ、我々は判断を避けるべきだという結論に達しなければならない。しかし、犯罪者が未処罰でいる危険性があるため、判決が必要となる。この判決は、被疑者が無罪でありながら有罪とされた場合、市民が誤りの恐れを感じる理由が、彼が有罪でありながら無罪とされた場合に、再犯や他の不幸な人々の例によって激励される新たな犯罪を懼れる理由よりも少ないかどうかという問いに帰着する。この問いに対する解答は複数の要素により左右され、それらは非常に困難であることを認識しなければならない。その一つが、犯罪者が未処罰でいる場合に社会に脅威となる危険性の高さである。時としてこの危険性は極めて高く、裁判官は適切な無罪保護のための形式を正当化して放棄せざるを得ない状況が生じる。しかし、被疑者の犯罪に対する確率を正確に評価し、有罪判決が必要とする確率を定めることがほとんど不可能であるという事実が、この問いをほぼ常に解けないものとしてしまう。各々の裁判官はこの点で自分の判断に依拠せざるを得ず、被疑者の証言や状況と彼の反省と経験を比較し、真実を判明するために長い間有罪者を尋問し、裁く経験が大きな利点となる。

次の質問は、犯罪の調査における注意深さにかかっているため、死刑を科す際には、数ヶ月の拘留よりもはるかに強い証拠が求められます。重大な罪に対して適切な注意を払うことは、無実者を誤って処罰する危険性を減らします。裁判所において、軽微な事案に対する注意を緩和できる権限を与える法律は、犯人への人道的な配慮と社会の利益に合致しています。

被疑者の有罪の確率とその重大さの積が、被疑者が無罪放免された場合に社会に及ぼす危険性を測る指標であるため、注意の程度はこの確率に依存すべきと考えられます。裁判所では、非常に強い証拠があるにもかかわらず、まだ有罪と断定するには不十分な被疑者に対して一定期間拘束し、新たな証拠を待つという方法が採られています。このような措置は、再び市民の間で彼を見ることなく安心して生活できるようにすることを目指しています。

しかし、この措置の恣意性とその濫用により、個人の自由を最も重視する国々では、この制度が拒否されています。

現在の裁判所が特定の多数決により裁く場合、その決定が公正である確率はどれくらいか？つまり、提案された問題に対する真実の解答に合致する確率はどの程度か？この重要な問題を適切に解決すれば、異なる裁判所間の比較手段を得ることができる。多数決の一票差は、争点が非常に不明瞭であることを示す。被告の裁決は人道の原則、無罪を保護するものから逸脱することになる。全員一致の判決は公正な決定に対する強固な確率を与えるが、それを避けると多くの有罪者が無罪となる可能性がある。したがって、全員一致を求めたい場合は裁判所の人数を制限するか、より多くの多数決が必要な場合、裁判所が増加するにつれて裁決の必要な多数決を増やすことが必要である。私はこの問題に微分積分学を適用してみようとする。実感から得られるデータに基づいて最善のガイドであると確信しているからだ。

各判事の意見が公正である確率は、この計算における主な要素となる。千一名前の裁判所において五票差がある場合、それぞれの判事が公正な見解を持っている確率はこの計算に重要な役割を果たす。

百一は一方の意見に一致しており、五百人は反対の意見に立っています。したがって、各裁判官の意見が$p$より僅かに優れていることが明白です。もし各裁判官の意見が明らかに大きいと仮定すると、単純な一人の投票差異は不確実な出来事となります。しかし、裁判官が一致している場合、それは証拠の強度を示し、その結果信頼性が生じます。各裁判官の意見がそれほど近い確率であることを示します（ただし、情熱や一般的な偏見が同時にすべての裁判官に影響を与えない場合）。これらの状況以外では、有罪か無罪に対する投票数の比だけがこの確率を決定すべきです。したがって、私はその確率が$p$から1に変化する可能性があると仮定しますが、それは$p$未満にはならないと考えます。もしそうでないなら、裁判所の決定は運に等しくなり、価値は各裁判官の意見が誤りよりも真実性に向かう傾向が高いという点でのみ存在します。したがって、有罪または無罪に対する支持票と反対票の数の比によってこの意見の確率を決定します。

これらのデータは、判決が既知の多数決により行われる裁判所の公正な決定が正しかった場合の一般的な確率を確定するのに十分です。八人の裁判官で五人の賛成が必要な裁判所では、その決定における誤りの可能性は$p$を超えるでしょう。もし裁判所が六人になり、四票の多数により有罪判決を行うことができれば、その決定における誤りの可能性は$p$以下になります。したがって、被告には裁判所の縮小によって利益があると言えます。両方の場合で必要となる多数は同じであり、

裁判所の判断確率に関する論理はこのように定義されます。

これらのデータにより、既知の多数決により裁かれる裁判所の決定が公正である確率の一般的な表現を特定することができます。八人の裁判官で五票が必要な有罪判決の場合、その決定における懸念すべき誤りの確率は$p$を超えるでしょう。もし裁判所が六人になり、四票の多数により有罪判決を行うことができれば、その決定における懸念すべき誤りの確率は$p$以下になります。したがって、被告には裁判所の縮小によって利益があると言えます。両方の場合で必要となる多数は同じであり、

必要となる多数は同じであり、

PROBABILITY OF THE JUDGMENTS OF TRIBUNALS. 139

二つと等しい。従って、多数決が一定である場合、審判者の数が増えるにつれて誤りの確率は高まります。これは、必要な多数決がどのようなものであれ、その多数決が一定である限り普遍的です。

次に、算術的な比を規則とするとして、被告は tribunal がより多くなるほど有利な立場にはならず、むしろ不利になります。一人でも多くの審判者が必要となると、十二人の多数決が必要な英国での有罪判決の状況とは大きく異なるでしょう。少数派の投票が多数派と同じ数の投票を中和するのであれば、十二人の残りの投票は十二人の陪審員による一致した判決を表していると考えるかもしれませんが、これは大きな誤りです。

普通の感覚では、二百人と十二人の審判者で、それぞれ百十二人が被告を有罪とし、百人が無罪とする場合と、十二人の全員が一致して被告を有罪とする場合との間には違いがあることがわかります。前者の場合、被告に有利な百票は証拠が十分に強力であるとは思われないことを示していますが、後者の場合は審判者たちの一致した判断から、その程度まで到達していると信じられます。

しかし、この二つの場合における誤りの確率の極端な違いを理解するには単なる普通の感覚だけでは不十分です。そこで計算に recur し、まず最初の場合では約五分の一、そして後者の場合は $-\frac{1}{500}$ の誤りの確率を見つけます。これは最初の確率の千分の一未満です。

これは、二つの場合における誤りの確率の極端な違いを確認するものです。

判決に及ぼす数人の陪審員の数による算術的な比率が被告にとって不利である一方で、幾何学的な比率を基準とすると、陪審員の数が増えるにつれて誤った決定の確率は減少します。例えば、三分の二以上の多数決で被告を裁くことができる裁判所では、陪審員が6人の場合、誤りの可能性は約四分の一です；陪審員の数を12人に増やすとその確率はさらに低くなります。したがって、誤りの確率が常に特定の分数以上または以下になることがないよう望むなら、算術的な比率や幾何学的な比率に従うべきではありません。

しかし、どの分数を決定すべきか？ここから恣意性が始まり、裁判所はこの点で最も多様な選択肢を提供します。8人の陪審員中5人が賛成すれば被告を裁くことができる特別な裁判所では、正義の観点からの誤りの確率は約6/13、つまりより半分以上です。この分数の大きさは恐ろしいものですが、少し心を落ち着かせるべきなのは、判決で無罪を宣告した陪審員が被告が実際には無罪であると認識していないことが多いという考慮点です。彼はただ、有罪に至る十分な証拠がないことを述べているだけです。特に自然が人間の心に植え付けた同情心により、被告が被疑者であることに不情愿地目を向けさせることで、無知な陪審員による不便を補うことができます。12人の陪審員からなる陪審団において、多数決が得られなければなりません。

十二人のうち八人の賛成で-condemnation-を宣告することが求められる場合、懸念される誤りの確率は|fffであり、これより少し大きい約八分の一である。この多数が九票の場合、ほぼ同じ確率であると言える。

一致判決の場合、懸念される誤りの確率は-g-T17jであり、つまり、私たちの陪審員による判決と比較して千倍以上も少ない。これは、一致判決が被疑者に対する有利または不利な証拠のみから生じるという仮定に基づいており、しかし、全く異なる動機もしばしば必要条件として陪審員の判断に影響を与える場合がある。

陪審員の決定は、陪審員の性質、性格、習慣、そして彼らが置かれた状況によって左右されるため、陪審員の大多数が証拠のみを考慮して判決を下す場合よりもしばしば異なる結果となることがある。この点については私には大きな欠点があると感じている。

私たちの陪審員による決定の確率は過不足があり、無罪を十分に保証するためには十二人のうち九票以上の多数が必要であると考える。

\newpage

\chapter{生命表・平均余命・婚姻に関する確率論}
\label{ch:xiv}

\begin{center}
\large\textit{人口動態統計・死亡率曲線・天然痘ワクチンの有効性分析}
\end{center}
\vspace{4mm}

次に述べる死亡率表の作成方法は非常に単純です。民事登録簿から、出産と死が記載された多くの個人を取り扱います。これらの個体中で、最初の年齢で亡くなったものが何人いたか、2番目の年齢で亡くなったものが何人いたかを定めます。そして、各年の開始時に生存していた個体数から結論を導き出し、その数を表の側方に記入します。

零の側には出生数を、一年目に達した乳児の数をIの年齢の側に、二年目に達した乳児の数を2の年齢の側にと、それぞれ記入します。しかし、生命の最初の二年間における死亡率は非常に高いので、この最初の年齢においては、各半年にわたる生存者数も示す必要があります。

すべての個体の寿命の合計年数を死亡率表に記載された全ての個人の数で割った値は、平均寿命を計算する上で重要です。

これらの個人について、この表に対応する平均寿命を得るために、最初の年における死者数を二分の一年に、側方に記された0年と1年の個体数の差に等しいもので乗じます。彼らの死亡率が一年にわたって分布されているため、その平均寿命はただの半分の年となります。第二年における死者数を一年半に、第三年における死者数を二年半にそれぞれ乗じます。そして、このような積の総和を出生数で割ったものが平均寿命となります。

この結論を得るには、各年ごとの表側に記された個体数の合計を出生数で割り、その商から半分を引くと、同じ結果が得られます。年は1として取ります。任意の齢から始めても、その齢に到達した個体数に基づいて同様の方法で平均寿命を求めることができます。しかし、平均寿命が最も大きいのは出生直後ではなく、幼少期の危険を脱したときであり、その頃は約43歳となります。ある齢から特定の齢まで生き延びる確率は、表においてこれらの2つの齢に対応する個体数の比に等しいです。

これらの結果の精度は、表を作成する際に非常に多くの出生数を使用することを要求します。解析は、それらの表に示された数字が真実からどの程度異なるかを評価するために非常に単純な公式を与えることになります。これらの公式により、考慮すべき出生数が多いほど、限られた範囲内でしか誤差がないことがわかります。したがって、使用する出生数が無限大である場合、表は正確に死亡率の真の法則を表現します。

次に、死亡率の表は人間の生存確率の表です。各年ごとの表記された個体数と出生数の比は、新しい出生がその年まで生き延びる確率を示します。希望の価値を評価するために、各利益の期待値を取得する確率にそれぞれの製品を足し合わせることで計算すると、同様に各年ごとの平均生存期間を評価することができます。これは各年の半分の総和とその年の開始と終了の生存確率の積を加えることで得られます。これにより先ほど得られた結果が導かれます。

しかし、この方法は平均生存期間を見る際の利点として、恒常的な人口における平均生存期間が年間出生数に対する人口比であることを示すことができます。つまり、出生数と死亡数が等しい恒常的な人口において、各連続年の間に存在する個体数は、年間出生数をその年の開始と終了の生存確率の総和の半分で乗じたものに等しくなります。これらの製品の合計は全体の人口となります。したがって、この合計を年間出生数で割ると、先ほど定義した平均生存期間と一致します。

mortality の表を使用して、人口が一定であると仮定した場合の対応する表を作成することは容易です。これには、死亡率の表における年齢 0 および 1 年、1 および 2 年、2 および 3 年などの対応する数値の算術平均を取ることで行います。これらのすべての平均値の合計が全体の人口であり、その合計は年齢 0 の側に記載されます。この合計から最初の平均値を引くと、残りが 1 年以上の人々の数となり、その数は年 1 の側に記載されます。さらに最初の余剰量から第2の平均値を引くと、これが 2 年以上の人々の数となり、その数は年 2 の側に記載されます。そして同様にして続きます。

死亡率を表すものは、多くの変動要因によって影響を受けますので、これらの表は場所や時間に応じて変更されるべきです。生活状態の多種多様な状況においては、それぞれの状態から不可欠であるとされる労働と危険が異なる点があり、これらは寿命に基づく計算において考慮する必要があります。しかし、これらの差異は十分に観察されていません。いずれかの日にはそれらが観察され、そして各職業が要求する生命の犠牲を知るようになるでしょう。そしてその知識を利用して危険を減らすことができます。

太陽の影響度、その高度、温度、住民の習慣、そして政府の活動は、死亡率に大きな影響を与えます。しかし、これらの原因が観察された差異を調査する前に、その原因が示唆される確率についての調査を行うことが常に必要です。例えば、フランスで年間出産数に対する人口比が二十八分の三まで上昇した一方で、古代ミラノ公国では二十五に留まっています。これらの比率は、多くの出産に基づいて確立されているため、ミラノの人々には特別な死亡率の原因があると疑う余地はありません。この原因を政府が調査し、除去することが我が国の重要な課題です。

人口比が出生数に対して再び増加する可能性があります。もし特定の広く蔓延している危険な病気を減らし、取り除くことができれば、その結果です。小児麻疹については幸いにも、最初はこの病気の接種により、その後はジェナーによって発明された無痛で優れた方法であるワクチン接種により、その問題が解決されました。ジェナーは人類にとって最も偉大な恩恵をもたらした一人の人物となりました。

小児麻痺は特有の点として、同じ個人が二度とその病にかかるわけではなく、または少なくともそのような稀なケースは計算から除外されるべきであるということが挙げられる。この疾患は、 Vacxin の発見以前にはほとんど生存しなかったものであり、しばしば致死的で、攻撃された人々のうちの七分の一が死亡する原因となる。しかし時々は軽い症状を示し、経験から学んだところでは、適切な食事と好適な季節に健康な人に対してその疾患を接種することで、この後者の方へ性格を変えることができる。すると、接種を受けた人々のうち死亡する個人との比は三百分の一である。この接種による大きな利点は、顔貌を変更しないことと自然の小児麻痺がしばしばもたらす重篤な後遺症から保護することにあるため、多くの人々によって採用された。その実践は強く推奨されたが、それは

TABLES OF MORTALITY, ETC. 145

不便を伴うことで強く争われたことには、通常通りです。この論争の真っただ中でダニエル・ベルヌは、予防接種が平均寿命に与える影響を確率の計算に委ねることを提案しました。当時の麻疹による死亡率の正確なデータが不足していたため、彼は各年齢での麻疹感染とその結果死亡する危険性が同じであるという仮定を行いました。これらの仮定に基づき、彼は精密な分析を行い、一般的な死亡率表を麻疹が存在しない場合や、またはその影響が非常に少ない場合に使用されるものへと変換し、これにより平均寿命が少なくとも3年延びることが示されました。彼はこの結果から、予防接種が疑問の余地なく生命の平均寿命を大幅に増加させる利点があると結論付けました。

ダランベールはベルヌの分析を攻撃しました。最初には彼の2つの仮定の不確実性について、次にはその不足について批判しました。それは自然発症の麻疹への即時的な危険性、非常に小さいものの、予防接種による死亡危険と比較することなく、自然発症の麻疹に屈服する非常に大きなが非常に遠い危険を考慮していないからです。この考慮は多くの個人を考えるときには消えますが、政府にとっては重要で、予防接種の利点は依然として存在します。しかし、家族の父親にとっては重大な問題であり、子供たちを予防接種に恐れを感じさせ、最も愛する子が死ぬことを恐れるのです。多くの親はこの恐怖によって制約されましたが、 Vacxin の発見により幸いにもその懸念は解消されました。

自然が我々に頻繁に提供する一つの神秘的な現象として、Vacxin は Variolar 病毒と同様に確実な予防接種であり、危険性はありません。それはどの病気も引き起こさず、非常に少量の注意を要するだけです。したがって、その実践は迅速に広まり、それを全人類へ普及させるには、人々の自然的な惰性を克服する必要があります。これは彼ら自身の最も切迫した利益に関わる問題であっても、その抵抗を常に乗り越える必要があるのです。

疾患の絶滅が生む利益を計算する最も単純な方法は、特定の年齢で毎年その疾患により亡くなる個人数を観察し、同じ年齢での全死亡者数からこれを引くことです。特定の年齢におけるこの差異の比は、疾患が存在しない場合にその年齢で死ぬ確率となります。その後、出生から任意の特定年齢までこれらの確率を合計し、それを1から引きます。残りの数値は、疾患が存在しない場合にその年齢まで生き延びる確率となります。これらの確率の系列は、この仮定に対する死亡表となり、それから先に述べた通り、平均寿命を推定することができます。Duvilard は、 Vacxin 接種により生じる平均寿命の増加が少なくとも3年であることを発見しました。このような大きな増加は、他の理由による人口制約がない限り、人口に非常に大きな増加をもたらすでしょう。

人口の進行的な増加は主に食料不足によって止まることです。

死亡表等に関する考察 14?

すべての動物と植物について、自然はunceasingに個体数を増加させる傾向があります。その数は生活費の平均水準まで達するまで続きます。人類の場合には、道徳的な要因が人口に大きな影響を与えます。森林地帯の容易な開発が新しい世代に豊富な食料を提供できるなら、多くの家族を養う確実性は結婚を奨励し、その結果産まれる子供の数も増加します。同じ土地においては、人口と出産は同時に等比級数で増加するべきです。しかし、森林地帯の開発が難しくなると、人口の増加も鈍ります。それは生活費の変動状態に近づき続け、その周りを周期的に振れることになります。これは、懸垂点を変えた振り子が重力によってその点を中心に振れるように、周期性は遅れて変化します。

人口の最大増加率を評価するのは難しいですが、観察によれば、有利な状況下では人類の人口は15年ごとに倍になることがわかります。北アメリカにおいては、この倍増期は22年と推定されます。このような状況で、人口、出産数、結婚、死亡率はすべて連続する項の比が一定である等比級数に従って増加します。

人間の寿命を表す死亡率のテーブルを使用して、婚姻期間を決定することができます。両性の死亡率が同じと仮定することで、結婚の確率を求めることができます。

夫婦の婚姻が一年、二年、三年など持続する確率は、表に記載された二人の年齢に対応する数値の積を公分母とする分数列によって形成され、その分子はそれぞれの年齢に加えて一、二、三等年数を増やした数値の連続的な積である。これらの分数の和に半分を加えたものが婚姻の平均持続期間であり、年は単位として取り扱われる。同じ規則を三人以上の個体で構成される協会の平均持続期間にも拡張することが容易である。

\newpage

\chapter{事象の確率に依存する制度の利益}
\label{ch:xv}

\begin{center}
\large\textit{海上保険・生命保険・年金基金とリスク分散の数理}
\end{center}
\vspace{4mm}

事象の確率に基づくもの。

ここでは、希望について述べた際に見たことを思い出す。複数の単純な事象があり、そのうちのいくつかは利益をもたらし、他のいくつかは損失をもたらす場合、各有利な事象の確率とそれがもたらす利益の積を加え、それらの合計から各不利な事象の確率とそれが付随する損失の積の合計を引くことで、得られる利点を得ることができるということが見てきた。しかし、これらの総和の差が示すいかなる利点も、これら単純な事象から成る複合事象が発生しても、損失を経験する懸念から解放されるわけではない。この懸念は、複合事象を増やすことで減少すると想像される。確率の分析は、以下の一般定理に到達する。

有利な単純な事象や複合事象を繰り返し行うことで、実際の利益はますます確実になり続け、無限回の繰り返し仮定下では確定的になる。その数で割ることにより、各事象の平均的な利益が得られる。これは数学的希望自身であり、事象に対する相対的な利点である。同様に、損失も長期的には確定的になり、事象の不利性がいかに小さくても同じである。

この利益と損失に関する定理は、既に述べた単純なまたは複合な事象の無限回繰り返しによって示される比と類似しており、それらと同じように、規則性は最も危険と名付けられるものでも最終的に確立されることを証明する。

多くの事象が存在する場合、解析はその利益が特定の範囲内に収まる確率を非常に単純な表現で与えます。これは、先に述べた確率に関する一般法則に再び登場する表現です。

確率に基づく制度の安定性は、前記の定理の真実性に依存しますが、それを適用するためには、これらの制度が多くの利益をもたらすために有利な事象を増やすことが必要です。

人間の寿命に関する確率に基づいて様々な制度が構築されてきました。例えば、生命年金やトントーニンなどがあります。これらの制度の利益と費用を計算する最も一般的で単純な方法は、実際の額に還元することです。1年の利息率は、

確率に基づく制度. 151

そのように呼ばれます。毎年末、ある金額は利息率加1の因子により増加します；そしてそれはこの因子が比となる幾何級数の項として増加します。したがって時間の経過とともにそれは巨額になります。例えば、利息率が5\%である場合、資本は約14年でほぼ2倍になり、29年で4倍になり、3世紀未満で200万倍に増大します。

莫大な増加が、公債の返済に利用するための考えを生み出し、その実現に向けて定期基金を設けることを考案した。この基金は年間で一定額を用いて公共債券の返還を行い、毎回返還される債券の利子によって継続的に増加させるものである。長期的には、この基金が国家債務の大きな部分を吸収することになることが明らかだ。

国の需要により借入が必要となったとき、その一部を定期基金の増額に充てることで公共債券の変動を抑えることができる。これにより貸与者の信頼と、借入した資金を失わずに返還できる確率が高まり、借入条件がより有利になる。実践的な経験はこれらの利点を完全に確認している。

しかし、このような機関の成功には、立法権がいくつかの独立した権力によって分割されている政府による信頼性と安定性が必要である。これらの権力間の必要な協力を生み出す信頼は、国家の力を倍増させ、その結果、君主自身も恣意的な権限を失うよりも法的権限を得ることができる。

前段落から導き出される結論として、ある金額が一定年数後に支払われる確率に相当する実質的な資本は、その支払い時に支払われる可能性の確率を支払時期までの期間数で累乗し、1加算した利子率で割り、そしてその結果を得ることになる。

すなわち、ある金額が特定の年数後に支払われるための実質的な資本は、その支払い時に支払われる可能性の確率をその支払時期までの期間数で累乗し、1加算した利子率で割ったものと等しい。

この原理を一人または複数の人に対する生命年金や貯蓄銀行、どのような性質を持つ保険会社にも適用することは容易である。例えば、ある人が与えられた死亡率表に基づいて生命年金の表を作成することを考える場合、五年後に支払われる生命年金が実際の額に換算されるのは、その原則により、以下の二つの量の積に等しいことになる。すなわち、年金を単位の五乗から増加した利子率で割ったものと、支払い確率との積である。

この支払い確率は、表において年金が確定される人の現在の年齢に対応する生存者数と、その年齢に五年後の生存者数との逆比である。つまり、

$$
\text{支払い確率} = \frac{\text{生存者数（現在の年齢）}}{\text{生存者数（現在の年齢 + 5年）}}
$$

と表すことができる。

次に、生存率表において指定された年齢で生きている人数を分母とする一連の分数を考える。その分子は、年金を同じ年齢で生きている人数とその人数がそれぞれ1年、2年、等々に増加した数との積である。これらの分数の和が求められる生命年金の額となる。

例えば、ある人が生命年金制度を設立することを考える場合、

$$
\text{分子} = \text{年金} \times (\text{同じ年齢で生きている人数}) \times (1 + r)^n
$$

$$
\text{分母} = \text{生存者数（現在の年齢）}
$$

ここで、$r$ は利子率であり、$n$ は年数である。これらの分数を合計することで、指定された年齢での生命年金の額を求めることができる。

以上が原著の

彼は自身の死後の年を跨ぐ金額を受け継がれる生命保険を設定する。その価値を決定するために、彼が生前銀行から借入し、永続的な利子で同じ銀行に預け入れるという想定を行うことができる。明らかに、この資本は彼の死後の年終了時に遺族へと返還されるが、彼は毎年その生命保険金よりも永続的利子の方が大きい部分だけを支払うだろう。その後、生命保険表は彼が死亡後にもその資本を受け継がれるために銀行に年々支払うべき額を示す。

火災や嵐などの海難事故に対する保険、そしてその他の種類の同様な機関も同じ原理に基づいて計算される。海上で商船を持っている商人は、これらの船舶と貨物が遭遇する危険からその価値を保護するために保険を設定したいと考える。これを行うために、彼は一定の金額を保証会社に支払い、その会社は彼に対してその貨物と船舶の評価された価値に対する責任を持つことになる。この価値が支払われる保険料に対する割合は、船舶が遭遇する危険性によって決まり、港から同じ目的地に向かって出航した船舶の運命に関する多数の観察結果によってのみ評価されることが可能である。

もし被保険者が確率論の計算によって示された金額を保険会社に支払う場合、その保険会社は機関運営費を賄うことができず、それらはそのような保険のコストよりも遥かに大きな金額を支払わなければなりません。では、被保険者はどのような利益を得るのでしょうか？ここでは、不確定性に付随する道徳的な不利な点を考慮することが必要です。最も公正なゲームであっても、既に見たようにプレイヤーが確実な賭け金と不確かな利益の交換を行うことで不利になることがわかります。保険は不確かなものを確実なものと交換することで有利であるべきであり、これは我々が上記で与えた道徳的な希望を決定する規則からも明らかです。さらに、この規則により、常に道徳的な利益が残されることで、被保険者が支払うべき犠牲の範囲がどのように広がるかを見ることができます。その結果、この会社はこの利益を得ることによって非常に大きな利益を上げることができ、被保険者の人数が非常に多い場合に限って、その継続的な存在が可能となります。するとその利益は確実なものとなり、数学的希望と道徳的希望は一致します。解析は一般の定理へと導き、それが極めて多くの期待を持つ場合、二つの希望は絶えず近づき合い、無限の数の場合には必ず一致することを示します。

数学と道徳的な希望について述べたとき、期待される利益のリスクをその一部に分散することには道徳的な優位性があるということが言及されました。したがって、遠くの場所に金銭を送る際は、一つの船で送るよりも複数の船で送ることの方が遥かに良いです。これは相互保険によって達成されます。

例えば、同じ港から同じ目的地に向かう2隻の異なる船に乗っているそれぞれが同じ額を持つ2人の者が、到着したお金全てを半分ずつ分け合うという約束を結ぶ場合、明らかにこの約束により各者は二つの船間で期待する金額を均等に分配することになります。

確率に基づく制度 155

このような保険は常に損失の懸念について不確定性を残します。しかし、被保険者の数が増えるほどこの不確定性は減少し、道徳的な優位性もますます高まり、最終的には数学的優位性と一致する自然な限界に達します。これは、相互保険の協会が非常に多くの参加者を持つときには、被保険者がその協会よりもより有利であることを意味します。これらの協会は、与える利益に対して道徳的な優位性を常に数学的優位性未満とし、その管理の監視によって相互保険の優位性を補うことができます。

これらすべての結果は、既に述べたように、道徳的な優位性を表現する法則によらず独立しています。

自由な国民を一つの大きな団体と見做すことができる。そのメンバーは、この保証の費用を比例して負担することで、互いの財産を保護する。数多くの人々が結成した連合国においても、個々の個人が社会の中で享受しているような利点を得ることができるだろう。彼らの代表者による議会が、共通の利益に関わる対象について討議し、おそらくフランスの科学者の提案する重量、測定法、および貨幣制度は、この議会において最も商業関係に有用なものの一つとして採用されるであろう。

人間の寿命に関する確率に基づく諸制度の中で、より優れたものは、その収入を軽微な犠牲にして、将来の生活と家族の存続を確保するものである。ゲームが不道徳であるほど、これらの制度は習慣に有益であり、最も強固な人間の性向を促進するからである。政府は公共の運勢の起伏に関わらず、それらを奨励し尊敬すべきである。なぜなら、その希望は遠い将来を見据えているため、存在中に全ての不安から守られている場合のみに繁栄できるからである。代表制政体の制度がそれを確保するという点で優位性がある。

次に借入について述べます。永続的に借入するためには、毎年利息率による資本の積み立てが必要です。しかし、その本金を一定期間を通じて等しい支払いにより一括で返済したい場合があります。このような支払いは「年金」と呼ばれ、その価値はこのように求められます。各年金は実際の瞬間に減少させるためには、その支払期日から数えた年数分の利息率を加えた単位のべき乗で割るべきです。

したがって、最初の項が単位のべき乗で割った年金である幾何級数を形成し、最後の項は同じ量のべき乗で割った年金であるとき、この級数の和は借入された資本と等価であり、これにより年金の価値が決定されます。基金は底から見ればただの手段であり、永続的な賃料を年金に変換するもので、唯一の違いは借入による年金の場合には利息が一定であるのに対し、基金によって獲得された利息は変動していることです。両者の場合で同じであれば、基金によって獲得された年金はその資金から形成され、この年金に対して国は毎年基金に寄付します。

生命保証借入を設ける場合、生命年金表は任意の年齢での生命年金のために必要な資本を提供します。これらの表を使用して、借入した資本金から個々人に支払うべき賃料は単純な比例関係で求められます。これらの原理に基づいて、すべての種類の借入が計算できます。

我々がこれまで説明した、制度による利益と損失に関する原理は、既に実施された何らかの観測結果に対する偏差を考慮する際に、それらの観測結果の数々に対する平均的な結果を決定するために利用できる。まず、最小の結果に対する修正を x とし、次々と g, q', q", 等によって増加した修正を次の結果に付与することを定義しよう。また、観測誤差 e, e', e", 等を観測の法則が既知であるとする。各観測は結果の関数であり、この結果に対する修正 x が非常に小さいと仮定すると、最初の観測の誤差 e は x に特定の係数を乗じたものと等しいことが容易に理解できる。同様に、次の観測の誤差 e' は q に x を加えた和に特定の係数を乗じたものと等しくなる。そして、誤差 e の確率が既知の関数で与えられた場合、それは最初の前の積分の同じ関数で表現される。誤差 e' の確率は、これらの積分の二番目の同じ関数で表現され、他のものは同様に表現される。誤差 e, f', c", 等の同時存在の確率は、これらの異なる関数の積に比例し、この積は x に依存する関数となる。これに基づいて、x の値が x 軸上にある曲線の縦座標であると仮定すると、この曲線は x の様々な値に対する確率を表す。これらの x の限界は誤差 e, e', e", 等の限界によって決定される。

必要とする x の値を X とし、X から x を引くと、x が真の修正である場合に生じる誤差が得られる。この誤差を確率 x または曲線の対応する縦座標で乗じると、その積は損失の確率として考慮される。これを x の微分で乗じて、曲線の最初の端点から X までの積分は X に対する x 値が X よりも小さい場合の不利を表す。x 値が X よりも大きい場合、X から x を引くと、x が真の修正である場合に X の誤差となる。この積分は、曲線の対応する縦座標と x の微分との積を X から曲線の最後の端点まで積分することで、X に対する x 値が X よりも大きい場合の不利を表す。これらの不利を加えると、合計は X の選択に付随する不利となる。この選択は、この不利が最小になるという条件によって決定されるべきであり、これにより X は曲線を二等分する縦座標を持つ x 軸であることが示されるので、そのような選択は確実なものとなる。

確率に基づく制度。159

真の x の値が X の一方または他方には落ちないという事実を考慮に入れます。

有名な幾何学者たちは、X を x の最も確かな値として選び、従って曲線の縦軸の最大値に対応する値を選んだのですが、しかし私には、確率論の理論が示す明らかな値はその前のものであることが明らかです。

\newpage

\chapter{確率評価における人間の錯覚と誤謬}
\label{ch:xvi}

\begin{center}
\large\textit{ギャンブラーの誤謬・認知バイアスと迷信の批判}
\end{center}
\vspace{4mm}

確率論について。

心には視覚の錯覚のようなものがあり、触覚が後者を訂正するように、反省と計算が前者を訂正します。日常的な経験に基づく確率や、恐怖と希望によって過大評価された確率は私達に強く印象を与えますが、それはただ単なる計算の結果です。したがって、小さな利益のために危険な状況に身を置くことを恐れず、フランスのlotteryでquinを引くよりも遥かに低い確率であるにもかかわらず、それでも誰もそのような利点を得るために自分の命を失う可能性がある certaintyで自己を危険から守る手段を望む人はいません。

私たちの情熱、偏見、支配的な意見は、有利な確率を過大評価し、反対の確率を薄めることによって、危険な錯覚の豊富な源となります。

現在の悪事とその原因が私達に与える影響は、逆の原因によって引き起こされた悪事が記憶されるよりも遥かに大きいです。それらは私達が一方の不利益と他方の不利益の不適切な評価をし、適切な手段による防衛の確率を理解するのを防ぎます。これが人々を長期間そして残酷な揺れ動きから休憩状態に戻るまで、断続的に独裁主義と無政府状態に導く原因となります。

事件の存在によって与えられる鮮明な印象は、他の人が観察している反対の事象をほとんど認識することができず、これは誤りに対する主な原因であり、充分に警戒することもできません。

確率的なゲーム、特に宝くじのようなものにおいては、多くの幻想が希望を支え、不利な状況にもかかわらずそれを維持するために働きます。 lottery をプレイする大多数の人々は、自分たちにとって有利な機会は何であるのか、不利な機会は何であるのかを知りません。彼らは小さな賭け金で大きな金額を得る可能性だけを見ていますし、想像力が生み出す計画は、その得られる可能性を彼らの目には過大に見せます；特に貧しい人は、運命を改善したいという欲求に駆られて、最も不利な組み合わせに賭けて、それらが彼に大きな利益をもたらすことを期待して、自分の必要品を失う危険性を無視します。もし彼らが失われた賭けの数々を知っていたら、疑いなく驚くでしょう；しかし逆に、勝利は大きな宣伝を受け、これによりこの葬送的な遊びに対する興奮がさらに高まります。

フランスのlotteryで長い間引かれていない番号が抽選に出てこない時間が長くなると、人々はその番号を賭け金で覆そうとする。彼らはこの番号が長い間引かれていなかったことから、次の抽選では他の番号よりも優先して引かれるだろうと考える。このような誤りは私には、過去の出来事への無意識的な思い返しによる幻覚に基づいているように思える。

例えば、コインを投げて表が出る確率が非常に低い場合でも、九回連続で表が出たとき、第十回では裏が出ると信じてしまう。しかし、コインが表の出やすさよりも裏の出やすさの方が大きいことを過去の経験から示すと、最初の出来事が次の出来事より確率的に高い可能性を持つことになる。そして、次に表が出る確率はさらに高くなる。

同様の幻覚は多くの人々を説得し、同じ番号に賭け金を続けることでlotteryで必ず勝つことができるという誤った信念を持たせる。一度の抽選でもそのような予想が持続できなければならなくなったとしても、それらは投資家たちの数学的な不利を減らすことはなく、むしろ彼らの道徳的な不利を増大させる可能性がある。なぜなら、各回の抽選で非常に大きな財産の一部を賭けることになるからである。

男たちは子供が欲しいと強く願っており、その中には、彼ら自身が父親になる月の生まれる男児の数に心配を感じながら学ぶ人もいました。彼らはこれらの男児の出産数が女児の出産数に対して同じ比率であるべきだと想像し、既に出産された男児が次の出産が女児である確率を高めると判断しました。したがって、 urn（玉入れ）の中に一定数の白い玉と黒い玉があるとき、白い玉を取り出すことで次の取り出しでは黒い玉が出る確率が高まります。しかし、 urnの中にある玉の数が無限である場合、つまり生まれつきの男児と女児との比較を行うために考えなければならない場合、この状況は変わります。

163  illusions in the estimation of probabilities

もし一ヶ月間に生まれた男児の数が女児の数を遥かに上回るようなことがあれば、そのときには彼らの受精時期に一般的な原因が男性受精を促進し、次に生まれる子供も男児である確率が高まる可能性があると疑うかもしれません。自然の不規則な出来事は、各抽選ごとにすべての番号が完全に等しく抽出されるように混ぜられるlottery（宝くじ）の数字の抽選とは正確には比較できないでしょう。これらの一つの出来事が頻繁に起こること自体は、その出来事がやや促進されている原因を示し、次回の再発現の確率を高めます。そして、長い期間にわたる連続した雨日のような不規則な出来事の繰り返しは、その変化の未知の原因を開示する可能性があり、各予想される出来事において、宝くじの抽選ごとに同じ状態の迷いから解放されることはありません。しかし、これらの出来事が観察される回数が増えるにつれて、それらの結果を宝くじの結果と比較することにより、その正確さはより高まります。

前回の幻覚とは異なるものとして、フランスのlotteryの過去の引当票から最も多く引かれた数字を見つけて、その上に賭けを組み立てるという方法があります。しかし、このlotteryで数字が混ざられる方法を考えると、過去は将来には影響を与えるべきではありません。頻繁に引かれた数字はただの運の異常であり、私はいくつかの例を計算にかけ、常にその結果がすべての数字が等しい確率で引かれることを前提とする可能性の範囲内にあることを確認しました。

同じ種類の事件の長い系列においては、偶然の単独の機会は時として運の良い流れや悪い流れを提供し、多くのプレイヤーはこれを運命とみなす傾向があります。同時に運とプレイヤーの能力に依存するゲームでは、失敗した人が自分の損失に悩むと、それを補うために危険な投げ方をすることがあり、別の状況では避けるべきものでもあります。これにより彼自身の不運を悪化させ、その期間を延ばすことがあります。そのときには慎重さが求められ、不利な機会に関連する道徳的な欠点が不運自体によって増大することを確認することが重要です。

人間は昔から自己中心的な立場に置かれており、自然が自分を特別な対象として特別な配慮をしていると信じている。この信念に基づいて、個々人は自分の周りに広がるより狭いまたは広い範囲の領域を自らの中心とするし、運命が自分に対して優遇していると考えることが多い。その信念を支えられると、賭博家たちは不利な確率でも大きな金額を賭けてしまうことがある。

生活における行動においても、同様の信念は時として利点を持つこともあるが、多くの場合、それは災害的な事業につながる。ここでもどこでも幻想は危険であり、真実だけが一般的に有用である。

確率論の計算法の大きな利点の一つは、最初の印象を疑うことを教えることにある。我々が計算によってその誤り性を認識するとき、他の事柄についても極めて慎重に考慮した上で信頼を示すべきであると結論付けるべきである。

このことを例によって証明しよう。

四つの玉で満たされた壷があり、これらの玉はすべて同じ色ではない。その中から白い玉が一つ引き出され、それを元の位置に戻して同様の引き出しを行うことを繰り返す。引き出しの四回目までに黒い玉だけを引き出す確率を求めよ。

白球と黒球の数が等しい場合、この確率は各回引き出される黒球の確率 $p$ の四乗に相当します。そのときそれは T4 となります。しかし、最初の一回で白球を引くことは、urn 内の白球の多さが優位であることを示唆します。例えば、urn 内に3つの白球と1つの黒球がある場合、各回白球を引き出す確率は | です；2つの白球と2つの黒球がある場合、その確率は $p$ になります；さらに、3つの黒球と1つの白球がある場合、その確率は J となります。事件から原因の確率の原理に基づいて、これらの三つの仮定の確率は $p$, f, $p$ の量比で互いに等しいと推測されます。したがって、黒球の数が白球の数よりも少ないか、または少なくとも同じであるというベットは 5 対 1 です。

最初の一回で白球を引いた後、次々と4つの黒球を引き出す確率は、色の等しい場合や十六分の一未満であるべきだと考えられます。しかし、それはそうではなく、非常に単純な計算により、この確率が四分の一以上であることがわかります。実際には、その確率は十六分の一よりも大きいことが分かります。これは $p$ の四乗

$p$の冪乗、|の冪乗、および|の冪乗は、前段落で述べた玉の色に関する仮定の最初のもの、第二のもの、第三のものにおいてそれぞれ発生します。各冪乗を対応する仮定の確率、すなわちf, f, と$p$にそれぞれ乗じて、その積の合計が連続して4個の黒い玉を引き出す確率となります。したがって、この確率は2f9fとなります。これはfより大きい分数です。このパラドックスは、最初の引く際に白い球の優越性が黒い球の優越性を完全に排除しないことを考慮することで説明できます。つまり、白い球の優越性が存在する可能性があり、これが色の等価性の仮定を否定します。しかし、この優越性は非常に低い確率であるとしても、与えられた黒い玉の数が相当大きい場合に引き連続して黒い玉を引く確率が等価性の仮定よりも大きくなるべきです。そして、先ほど見たように、これは与えられた数が4個になるときに始まります。

次に、いくつかの白い球と黒い球で構成された別の urn を考えてみましょう。最初に、その中には1つの白い球と1つの黒い球のみがあると仮定します。この場合、1回の引くことで白い球を引き出す確率は等しい賭けとなります。しかし、平等な賭けとしては、 urn が2個の黒い球と1個の白い球を持っている場合に白い球を引き出すためには2回の引くが必要であり、3個の黒い球と1個の白い球を持っている場合に3回の引くが必要であるなどと考えられます。各回の引く後に引き出した玉は urn に戻されます。

私たちは容易にこの最初の考えが誤りであることを確信します。実際、2個の黒い球と1個の白い球がある場合、2回連続して黒い球を引く確率はfの二乗、すなわち9fとなります。しかし、

確率評価における幻影。167

確率を引くと、二回の引きで白い球を引く確率に加えると、確定性または一に等しい確率があります。なぜなら、二つの黒い球または少なくとも一つの白い球が引かれることが確定しているからです。その場合の確率は -| となり、これは f よりも大きい分数となります。一方で、五回の引きで一つの白い球を引く賭けでは有利性があります。この urn 内に黒い球が五つ、白い球が一つあるときです。さらに、四回の引きでも有利性があります。これは六面のサイコロを四回投げた場合と同じ確率に戻ります。

メーテルリンク伯は、彼の友人であるパスカルという偉大な幾何学者に確率論の発明を促し、その結果確率論が生まれました。彼はパスカルに対して次のように言いました：「私たちは六を出す確率について計算したところ、一回のサイコロで六を出す場合と四回投げて六を出す場合との間で誤差があることに気づいた。一回のサイコロで六を出すには 671/625 の確率があり、一方で二つのサイコロを使って二つの六を出す場合は 24 回投げたときが不利である。少なくとも 24/36 は二つのサイコロの面の数であり、4/6 は一つのサイコロの面の数です」と。

パスカルはフェルマーに対して次のように述べました：「彼はこの比から誤差を見つけていたのです。一回のサイコロで六を出す場合と四回投げて六を出す場合との間で、671/625 の確率がありますが、二つのサイコロを使って二つの六を出す場合は 24 回投げたときが不利です。少なくとも 24/36 は二つのサイコロの面の数であり、4/6 は一つのサイコロの面の数です」と。彼は「このことは彼にとって大きな混乱を引き起こし、彼は大胆に言いました：「これらの命題は不変ではなく、算術が狂っていると言ったのです。. . . 彼は非常に良い頭脳を持っていますが、幾何学者ではないという点では大きな欠点があります」と。

メーテルリンク伯は誤った類推により、賭けの等価性において投げ回数が可能なすべての確率の数に比例して増加するべきであると考えていました。これは正確ではありませんが、この数値が大きくなるにつれて正確さに近づきます。

これらの議論を説明しようとする試みがあります。

男児の出産数が女児よりも多いという一般的な父親たちの願望により、名前を継続させるために息子を得ることを目指す。したがって、等しい数だけの白い球と黒い球で満たされた urn を想像し、多くの人がこの urn から球を取り出し、白い球を引き出したときに止むという意図を持って何度も繰り返すと、白い球が取り出される回数が黒い球よりも多いと考えられた。実際には、すべての引き取り後、取り出した白い球の数は参加者の数に等しくなり、黒い球を一度も引かない可能性がある。しかし、この意図は全体の引き取りにおいて白い球を取り出す回数に父親たちの意図が影響を与えるとは限らないことが容易に理解できる。最初の引き取りですべての人々が一度に urn から球を引くことを想像すると、その意図は球の色に影響を与えることは明らかではない。唯一の効果は、初めて白い球を引いた人を次回の引き取りから除外することである。新たに参加する人々の意図も、その後の引き取りにおいて球の色に影響を与えるとは限らないことが明白であり、これは続くすべての引き取りでも同じである。この意図は全体の引き取りにおいて球の色に影響を与えないが、各回の参加者数を増減させるだけである。白い球と黒い球の比率は従って1から大きく異ならない。そして、参加者の数が非常に大きい場合、抽出された色の比率が観察によって1から著しく異なるものであればあるほど、これは一種の錯覚であることが推測できる。

一様性において、白球と黒球の比率が玉入れの中にあるとき、その差異が単位と同様であることは非常に可能性が高い。

また、ライプニッツとダニエル・ベルヌーイが確率計算を級数の和に適用した illusions として再分類する。分数の分子が単位で、分母が単位に加えて変数を持つとき、この分数を特定の比率に基づくこの変数の冪の系列によって定義された系列に分割すると、変数を単位と仮定したときに分数が $p$ となり、系列は +1, -1, +1, -1, 等という形式となることが容易に理解できる。最初の二項を加え、次に次の二項を加えていくと、系列は各項がゼロとなる別のものに変換される。グランディ、イタリアのジェシトは、この事実から創造の可能性を結論付けた；なぜなら、系列が常に $p$ であることを観察し、この分数が無限数のゼロからもしくは何も生まれるのだと見たからだ。ライプニッツは彼の二進法アラビア数字において、単位と零という2つの文字のみを使用したときに、神を単位と表現し、何も零と表現することから、最高なる存在が何もないところからすべての存在を引き出したと考えた。彼は、全ての数がこの算術システムで表されるように、単位と零が全ての数を表現すると考えた。このアイデアはライプニッツに非常に好ましかったため、 jesuit グリマーディに、中国の数学裁判所の院長に、その創造の象徴が科学を特に愛するそこで在る皇帝をキリスト教に導くことを望んでそれを伝えた。この出来事を報告するのは、幼少期の偏見がいかに偉大な人物を誤らせるかを示すためである。

17番目の 論理学的確率論の小説。

ライプニッツは、常に独特で非常に緩い形而上学に導かれるまま、項が奇数個のとき零であり、偶数個のとき1であるというこれらの項からなる級数「+1, -1, +1, ...」が無限大において1または0になると考えていた。そして、無限大において偶数と奇数を好む理由はなく、確率論の規則に従ってこれら二種類の結果の半分を取り、これは級数の値として$p$を与える。

ダニエル・ベルヌーイは以来、周期項からなる級数の総和についてこの議論を拡張した。しかし、これらのすべての級数は厳密な意味での値を持たない；その項が変数の連続するべき乗に乘じられる場合のみに値を持つ。そして、この変数が1からどれだけ小さいものでも、これらの級数は常に収束し続ける。ベルヌーイによって確率論の規則に基づいて与えられた値は、これらの項を1と仮定した生成分数の値と同じであることが容易に証明できる。これらの値はまた、変数が1に近づくほど、級数がより近づく限界である。しかし、変数が正確に1に等しいとき、これらの級数は収束しなくなる；それらは項を止めるまでしか値を持たない。

この確率論の計算における驚異的な比は、周期的系列の値の極限と一致するものであり、これらの系列の項がすべての連続するべき乗の変数に乘じられる場合にのみ成り立つ。しかし、このような系列は異なる分数の展開からも得られることもある。したがって、「+1, -1, +1, ...」という級数は、

確率論における誤解：確率評価における幻影。

分数が、分子が変数に加えた単位と等しいもので、分母がその分子に変数の平方を加えたものから発展する可能性があります。変数を1と仮定すると、この発展は提案された級数において変化し、生成分数はfに等しくなります。確率の法則はそれによって誤った結果を導き出し、これは類似の推論を使用することの危険性を示しています。特に数学的な科学では、その操作がより厳格であるべきであるからです。

自然に導かれて、私たちは地上で見られるものの順序が常に存在し続けていることを信じるでしょう。実際、現在の宇宙の状態が過去の生成した状態と完全に類似している場合、それはまた類似の状態を生み出すでしょう。これらの状態の連続は永遠に続くことになります。私は解析を万有引力の法則に適用し、惑星や衛星の回転運動と公転運動、そして軌道および赤道の位置が周期的な不等式によってのみ支配されていることを発見しました。古代の日食と月の secular 方程式の理論を比較すると、ヒッパルコス以来、100分の1秒にも達しない時間変化と、地球の平均温度が100分の1度にも達しない減少が確認されました。したがって、実際の秩序の安定性は理論と観測によって同時に確立されています。しかし、その秩序は多様な原因によって形成され、それらは注意深く調査しても計算に提出することは不可能です。

海洋の動き、気象の変動、そして火山の噴火などによる地表への影響は、長期間にわたり大きな変化をもたらすでしょう。気候の温度、大気の体積、そしてその構成するガス類の割合がわずかながらも変化することがあります。これらの変化を定量化するための装置や方法がまだ新しいことから、観測はこれまでにそれらの変化について何も教えてくれていません。しかし、大気中のガス類が吸収されたり再生成されたりする原因が完全には一定しているとは思えません。数世紀という長い期間が経過することで、生命維持に不可欠なこれらの要素が経験した変化が明らかになるでしょう。

歴史的な記録は古くまでさかのぼりませんが、自然の作用によるゆっくりとした継続的な影響によって生じた大きな変化を示しています。地殻内部から見れば、過去の自然の痕跡が多数見つかります。これらは現在のものと全く異なるものです。

また、地球全体が当初液体状態だったとする説があり、その状態から現在の固体状態へと移行した際には、表面に驚くべき変化があったと考えられます。天体自体も、運動の秩序があるとはいえ、絶対的な不変性はないと考えられます。光や他の虚無気体に対する抵抗力、そして星への引力は、長い年月を経て惑星の動きに著しく影響を与えるでしょう。

illusions in the estimation of probabilities. 173

星や星雲の既に観察された変化は、これらの巨大な天体の系において時間が発展させるであろう変化を予感させてくれる。宇宙の連続的な状態を表す曲線を描くことができ、その曲線における縦軸は様々な状態であり、横軸は時間である。この曲線の一部しか知らない我々は、その起源にさかのぼる能力がまだ十分ではない。そして、想像力を満足させるために、私たちの知識からくる原因を理解することができない現象に対する無限の好奇心から慎重さを持って予想を提示することも賢明である。

確率の評価においては、特に知的組織の法則に依存する一種の幻覚が存在し、これらの法則を理解し、それによって生じる幻覚から自分自身を守るために深い検討が必要となる。未来への洞察と、ある重要な出来事の比率、占星術師や占い師による予言、予感や夢、運勢や不運とされる日付や数字に対する欲求は、多数の偏見を生み出し、これらは今でも広く存在している。多くの非一致が印象に残らなかったり、知られていないことに注意を向けないでいる。しかし、これらの一致が属性づけられる原因の確率を評価するために、そのような非一致についての知識を得ることは必要である。これにより、理由に基づいてこれらの偏見について語っている理性が我々に教えていることを確認することができるだろう。

古代の哲学者は、ある神殿で示され、そこで崇拝される神の力強さを高めるために、船難から救われた人々の祈りの効果を示す例を見せる。しかし、その中でも祈りを行っていながら船難に見舞われた人々の名前は記されていないという出来事を提示し、これは確率計算と一致していることを指摘する。

キケロは『占星術論』において、これらの偏見を理知と雄弁さで根拠なく否定しており、その結論として次のように述べている。古代の先哲たちの中に、理由の雷がすべての迷信を消し去った後、人間社会の制度の唯一の基礎となるべきものとなった場面を見つけるのが好きである。

ローマの演説家は言います：「占いや夢の兆候、その他の類似した偏見を排除する必要があります。広く根付いた迷信が多くの心を支配し、人間の弱さに侵食しています。我々は『神々の性質』に関する著作において、特にこの作品の中で、迷信を破壊することで他者と自分自身にサービスを提供できると信じてそのことを説明してきました。しかし（そして私は特にこの点について自分の考えがよく理解されるように願っています）、迷信を破壊することによって宗教を混乱させるつもりはありません。賢明は我々に、祖先の制度や儀式を維持し、神々の崇拝に関わる慣習を守ることを命じています。また、宇宙の美しさと天体の秩序が人間の心に優越した存在の認識を迫ります。そのような存在は人類にとって注目すべきものであり、称賛されるべきものです。しかし、宗教の普及は自然に対する知識と共に適切に行われれば、迷信の根絶に向けて努力することが必要です。迷信は人々を苦しめ、困らし、常にどこでも追いかけてきます。

確率の誤解：幻影

占い師や巫者に相談したり、犠牲を捧げたり、鳥の飛翔を見たり、カルデア人や占星術師に遭遇したり、雷が鳴ったり、 thunderbolt が落ちたり、最後には何か異常な出来事が起こったりした場合、迷信は支配し、安らぎを与えないままにしてきます。眠りさえも、人々の苦悩と労働からの避難所としての地位を失い、新たな不安と恐怖の源となります。」

これらの偏見と、それらが引き起こす恐怖は、理由がそれらから私達を解き放つ後も、時折生理学的な原因によって強く働き続けます。しかし、これらの偏見に反する行為の繰り返しは常に、それらの偏見を破壊することができます。

\newpage

\chapter{確実性に接近するための諸手段}
\label{ch:xvii}

\begin{center}
\large\textit{帰納法・類推・仮説検証・科学的方法論の哲学的基礎}
\end{center}
\vspace{4mm}

帰納法、類推、事実に基づいた仮説で、新たな観察によって継続的に修正されるもの、そして自然から与えられた幸運の直感と、その指針が経験との比較を通じて強化されることにより得られるもの、これらは主に真実に到達するための重要な手段です。

同じ性質を持つ一連の物体を考慮すると、それらの中にそして変化の中で、一定の比率を見ることができます。これらの比率は系列が延長されるほどより明確に現れ、継続的に拡大・一般化され、最終的にはその比率から導かれた原理へと至ります。しかし、これらの比率は多くの奇妙な状況で包まれており、それらを解きほぐし、元の原理に戻るためには大きな洞察力が必要です：これが科学の真の才能の所在です。解析学と自然哲学はこの豊かな手段、すなわち帰納法によって最も重要な発見を行いました。ニュートンはその二項定理の定理と普遍的重力の原理を得るためにそれを用いました。帰納法による結果の確率を評価することは困難ですが、これは最も一般的である最単純な比率が最も頻繁に現れるという事実に基づいています。これは解析学の公式において確認され、自然現象、結晶化、化学結合においても見ることができます。これらの比率の単純さは、自然界のすべての効果が少数の不変の法則の数学的結果であると考慮すると驚くべきものではありません。

しかし、帰納法は科学の一般原理の発見に導くことはできますが、絶対的な証明には至りません。それらを確立するためには、証明または決定的な実験によって確認する必要があります。科学の歴史は、帰納法がしばしば不正確な結果を導き出すことを示しています。例えば、フェルマーの定理について述べます。この偉大な数学者は、この定理について深く考え、単なる素数で構成される式を見つけ出そうとしました。その式から直接、任意に与えられた他のどの数よりも大きい素数が得られるものを求めたのです。

帰納法は彼を導き、二のべき乗が自身もまた二のべきであるとき、それと一を組み合わせると素数になるという結論に至らせました。したがって、二の平方加一から五という素数が得られ、二の二乗の二乗（十六）から十七という素数が得られます。彼はさらに、二の八乗と十六乗のそれぞれを一に加えたものが依然として素数であることを発見しました。これらの帰納法はいくつかの算術的考察に基づいており、彼はその結果を一般化したものと考えました。しかし、彼はそれを証明していなかったことを認めました。

実際、オイラーは二の三十二乗加一が六百四十一で割り切れるという事実を認めていました。

私たちは帰納的に判断し、様々な出来事や動きが長年にわたって単純な比率で結びついてきた場合、その比率は継続されることを推測します。そして、確率論の理論から、この比率は偶然によるものではなく、規則的な原因によるものであると結論付けます。月の回転と公転の運動の等しさ、軌道の節点と月の赤道との運動の等しさおよびその一致、木星の最初の三つの衛星の運動の比率、つまり第一衛星の平均経度が第二衛星の三倍から第三衛星の二倍を引いたものに二直角に等しいもの、潮汐間隔と月が子午線を通るまでの期間の等しさ、最大潮の再現と朔望との一致、最小潮の再現と弦月との一致。これらは最初観測されたときから維持されてきたことから、幾何学者が普遍的な重力法則に適応させた恒常的な原因の存在を極めて確実な可能性で示唆しています。

Bacon卿は真の哲学的方法の推進者であるが故に有名な演説家だが、地球の静止性を証明するために帰納法を非常に奇妙な方法で使用した。彼は「Novum Organum」の中で次のように論じる：「東から西への星の運動はその距離が地球から増えるほど速さが増す。これは最も遠い星で最も速く、土星では少し遅くなり、木星ではさらに遅くなる」というような推論を展開している。

月と最高の彗星。それはまだ大気中に存在し、特に熱帯では、空気の分子がそこで描く大きな円によって感じ取れる。最後に、それは海洋においてほとんど感じられず、それゆえ地球上ではほぼない。

しかし、この推論は、土星とその以下の星々は実際的なまたは観測される運動から独立した動きを持っていることを証明するだけである。これは、全体の天体球が東から西へと回転する実際的なまたは観測される運動とは対照的に、これらの動きはより遠くにある星ほど遅くなるという現象を説明し、光学法則に従っている。

Bacon は、地球が動かない場合、恒星がその日間の回転のために必要とする驚異的な速さについて打撃を受けたはずです。そして、その回転が遠く離れた天体、すなわち恒星、太陽、惑星、月すべてがこの回転に従うように見せかける極めて単純な方法であるという事実についても驚いたはずです。

海洋と大気については、それらの運動を地上から分離された恒星の運動と比較する必要はありません。しかし、空気と海は地球の大気に含まれているため、その運動や静止に参加すべきです。

Bacon の天才によって大きな展望が広がったにもかかわらず、宇宙のコペルニクス体系の壮麗な概念には説得されなかったことは奇妙です。しかし、ガリレオの発見を継承したコペルニクス体系に強い類似性を見出すことができたのです。

彼は真実探求のために教えるべき命令を与えましたが、例は示さなかった。

しかし、理屈と雄弁の全力を駆使して、無价值な伝統からの放棄の必要性を強調したことで、その体系を支持することができた。

学校の奥義を理解し、観察と経験に取り組むために、この偉大な哲学者は、現象の一般的な原因へ上昇する真の方法を示すことで、人類の心が巨大な躍進を遂げた18世紀において、その職業終焉までに莫大なる貢献を行った。

類推は、似ているものは同じ種類の原因によって同じ効果を生むという確率に基づくものである。この確率は、類似性がより完璧になるほど増加する。したがって、同じ器官を持ち、同じことをしている生物は、同じ感覚を経験し、同じ欲望に動かされることが疑いなく判明する。私たちは、私たちの種に属する個体と類似性のある動物が、私たちと同じような類推的な感覚を持っている確率は、私たちの種に比べて少し劣るとしても非常に大きいと判断する。宗教的偏見の影響により、一部の哲学者たちが動物を単なる自動機器であると考えるように説得されたことは、そのような信仰を必要とした。器官が私たちのものとより類似性が減少するほど、存在する感覚の確率は減少するが、それでも昆虫すら含めても非常に大きい。同じ種の個体が世代を超えて正確に同じ方法で非常に複雑なことを実行し、学習することなくそれを繰り返す光景を見ると、彼らは分子が結晶の構造を形成する際にも見られるような類似性のある親和性によって行動していると考えるようになる。この親和性には、すべての動物組織に付随する感覚が加わり、化学反応の規則性と同様の規則性で、より独特な組み合わせを生む。おそらく、このような選択的親和性と感覚の混ざり合いを名づけることができる。

確証に近づく様々な方法。181

動物と植物の類似性。植物と動物の組織には大きな類似性があるものの、それから植物に感情を拡張する理由は十分ではないと考えます。しかし、植物に感情があるとする証拠もありません。

太陽がその光と熱によって地球上を覆う動植物を生み出すようにして、我々は類推により他の惑星でも同様の効果があると判断します。太陽の活動が地球のような大きな惑星であるジュピターにおいても無意味であるとは考えにくいからです。ジュピターもまた日、夜、年を持ち、観測結果からは非常に活発な力が存在することを示唆しています。しかし、惑星の住人と地球の住人が類似しているという結論は、類推を過度に拡張したものと言えます。

人間は、その気候や環境に適応して生きており、他の惑星では呼吸できるような要素も温度も異なるため、おそらく生きられないでしょう。しかし、この宇宙の様々な球体に対して無限の組織が存在するべきではないでしょうか？地球の元素と気候の単純な違いでも地上の生物は多様性を示していますから、それぞれの惑星や衛星にはどれほど大きな違いがあるか想像もつきません。最も活発な想像力でもその存在を形にすることはできませんが、それは非常に可能性が高いことです。

太陽と同様の引力を持つ多くの太陽として、我々は星々について類推します。その引力は質量に比例し、距離の二乗に反比例するものであると考えられます。この力は実証されていますからです。

すべての太陽系の天体、そしてその最小限の分子についても、それはすべての物質に固有のものであるように思われる。既に、結びつきによって「双星」と呼ばれる小さな恒星の動きは、互いの引力を示唆しているように見える。正確な観測がせいぜい一世紀程度で、それらの回転運動や相互間の運動を検証すれば、その相互作用は疑う余地なく明らかになるだろう。

各恒星を太陽系の中心とする類推は、前記のものよりもはるかに弱いが、恒星と太陽の形成に関する仮説によって確実性を得る。この仮説によれば、各恒星はもともと広大な気体包囲層を有していたように考えられ、その気体包囲層には太陽の気体包囲層と同じ効果が期待され、それが凝縮して惑星や衛星を生み出したと考えられる。

科学における多くの発見は類推によるものである。最も著しい例として、雷電気現象と電気現象の類推から導かれた大気電気の発見がある。

真実探求において最確かな方法は、帰納法によって現象から法則へ、そして法則から力を引き出すことである。法則とは特定の現象が互いに結びつく比率であり、力の一般的原理を導き出す一般原則が示されたときには、直接的な経験でそれを検証するか、または既知の現象と一致しているかどうかを調べることによって検証する。そして厳格な分析によってそれがその原理から派生したものであることを理解すればよい。

多様な方法による確信への近道 183

原則においては、その小さな詳細まで忠実であり、さらに多種多様かつ非常に多数であるならば、科学は達することができる最高の確証と完璧さを獲得します。天文学者が普遍的な重力の発見によってなったようにです。科学の歴史は、導入の緩慢で労働集約的な道程が発明家たちにとって常にその通りではなかったことを示しています。想像力は原因に到達するのを待つことができず、仮説を作り出す喜びを感じ、しばしば事実を変更してそれを作業に合わせようとします；すると仮説は危険になります。しかし、現象を結びつけて法則を見つけるための手段としてそれらを考慮するとき、そして、現実に属性を与えないことで恒常的に新しい観察によって訂正することでそれらが真の原因へと導くことができるか、または少なくとも観測された現象から与えられた状況が生むべきものについて結論を下す立場を得ることができるとき、仮説は有用です。

我々が現象の原因に関するすべての仮説を試みるならば、排除の過程によって真実なものを到達します。この方法は成功を収めています；時として我々は、全ての既知の事実を同等に説明するいくつかの仮説に至り、それらについて学術者は分かれ、決定的な観察が真実なものを明らかにするまで争います。そのような仮説に対する歴史は人類の心に興味深いものであり、これらの仮説を再考し、多くの事実をどのように説明するかを見るだけでなく、それらが自然の歴史と一致するためにどのような変化を経るべきか調査することです。このようにしてピトゲーリの体系は、ただ天体の現実化であるという点で宇宙の体系から離れてきます。

184  確率論に関する哲学的考察。

観測結果は、太陽の軌道を年間で描かれる円やエピキュクルスと等しく且つ並行にするという仮定によって、惑星の運動に関する仮説に変換されます。彼が年間に描かせる円やエピキュクルスの大きさは未確定ですが、太陽の観測された運動を地球に対して反対方向に移動させることで、この仮説を真実の宇宙体系へと変えることができます。

これらの様々な方法によって得られた結果の確率を計算するの几乎常に不可能です。これは歴史的事実についても同じことが言えます。しかし、説明された全体的な現象や証拠が十分に大きい場合、確率を評価できないとしても、それらに対する疑問を持つことは理不尽であると言えることがあります。他の場合では、慎重に受け入れるべきです。

\newpage

\chapter{1816年までの確率論の歴史的概観}
\label{ch:xviii}

\begin{center}
\large\textit{パスカル、フェルマー、ホイヘンスからベルヌーイ、ベイズ、ラプラスまでの展開史}
\end{center}
\vspace{4mm}

昔々、最も単純なゲームにおいては、有利または不利となる確率の比が決定され、賭け金や賭け事はこれらの比に基づいて調整されていました。しかし、パスカルとフェルマト以前には、この主題を計算に提出する原理と方法が誰一人提示されていませんでしたし、このような種類の複雑な問題も解決されていませんでした。したがって、最初の確率の科学の要素は、これらの二人の偉大な幾何学者に帰属すべきであり、その発見は17世紀に人間の知性を最も栄光に輝かせたことの一つとして並べることができます — それは最大の功績を成し遂げた17世紀です。彼らが異なる方法で解決した主要な問題は、我々が見たように、ゲームを途中で打ち切った場合にプレイヤー間の賭け金を公正に分配する問題であり、プレイヤーは等しく技術があり、ゲームが終了する前に停止することに合意していると仮定されます。勝利には異なるプレイヤーごとに特定のポイント数が必要です。

185

186      確率に関する哲学的考察。

確実に、配分はそれぞれのプレイヤーがこのゲームで勝つ確率に比例して行われるべきであり、その確率はまだ欠けている点数によって依存します。パスカルの方法は非常に巧妙で、最終的にはこの問題の部分差分方程式を用いて連続する各プレイヤーの確率を決定することにより、最小の数字から次の数字へと進むものに過ぎません。この方法は2人のプレイヤーのみに限定されていますが、フェルマーの組合せに基づいた方法は任意のプレイヤー数に対応します。最初パスカルは自分の方法と同じく2人のプレイヤーに限定されていると考えていましたが、これにより彼とフェルマー間で議論が生じ、結局パスカルはフェルマーの方法の普遍性を認識しました。

ヒユーゲンは既に解決された様々な問題と新しい問題をまとめ、この主題に関する最初の短い論文を執筆しました。そのタイトルは「ラティオニシス・イン・ラドゥロ・アレッセ」です。その後も幾何学者たちはこの主題に関心を持ち続けました：大議長ホドー、オランダのウィット、そしてイギリスのハーレイは人間の寿命に関する確率に微分法を適用しました。ハーレイはこの分野で最初の死亡率表を発表しました。同じ頃、ジャック・ベルヌーイは幾何学者たちに対して様々な確率に関する問題を提示し、その後その解を示しました。最後には彼は美しい著作「アール・コンジクランディ」を完成させました。これは1706年に彼の死後7年後に出版されました。確率論の科学はこの作品においてより深く研究されています。著者は組合せと級数に関する一般理論を与えています。

187

この方法は、リスクに関する複数の難問への適用に使用されました。この作品は、その公正さと巧妙さ、二項式の公式の採用によるこの種の問題への適用、そして無限に観測と経験を重ねることによって、異なる性質の事象の比がそれぞれの確率に対する比に近づき、その差がより狭くなるとともに任意に小さな量よりも小さくなるという定理の証明により、依然として驚くべきものである。この定理は、観測を通じて現象の法則と原因を得るのに非常に有用です。ベルヌーイは、彼の証明に対する大きな重要性を正当に付与し、彼がそれを二十数年間思索したと言っています。

その間に、ジャック・ベルヌーイの死からこの作品の出版までの期間には、モンモルとモーブルが確率論の計算に関する二つの論文を発表しました。モンモルの論文は「賭博ゲームについての試み」というタイトルで、その中ではこの計算を様々なゲームに広く適用しています。著者は、ニコラス・ベルヌーイが複数の難問の巧妙な解決策を含む第二版に追加したいくつかの手紙を付け加えました。

モーブルの論文はモンモルの後に出され、1711年の哲学的トランスケーションズ誌で最初に発表されました。その後、著者はそれを単独で出版し、三版にわたって改善しました。この作品は主に二項式の公式に基づいており、その中には、その問題と解も同様に大きな普遍性がありますが、その特徴は再帰列の理論とそれの使用です。この理論は、定数係数を持つ有限差分方程式の積分であり、モーブルは非常に幸運な方法でこれを達成しました。

モーブルは、ジャック・ベルヌーイの大量観測によって決定される結果の確率に関する理論を取り上げ、ベルヌーイが示したように、事象が起こるべき比がそのそれぞれの確率に近づくという点を示すだけでなく、これらの二つの比の差が与えられた限界内にある確率についても優雅で簡潔な表現を与えています。彼は、非常に高い冪の二項式展開の最大項の比とすべての項の和との比、そしてこの項が隣接する項よりも多い部分の双曲線対数を求めます。

最大項がそれ自体多くの因子の積であるため、その数値計算は実行不可能になります。収束する近似を使ってこれを得るため、モーブルはステリングの定理を用います。これは非常に高い冪の二項式の平均項に関する定理で、特にこの定理が超越的な円周率と半径の比の平方根を、そのように思われるべきではなかった表現に導入することに注意深いものです。

さらに、モーブルはステリングから得た結果に強く打たれました。これは無限積から円周率の表式を導き出したものでした。ウォリスはこの表式を得るために驚くべき分析を行い、その中には非常に興味深く有用な定積分理論の芽が含まれていました。

多くの学者が、Deparcieux、Kersseboom、Wargentin、Dupre de SaintMaure、Simpson、Sussmilch、Messene、Moheau、Price、Bailey、Duvillardを含む中で、人口、出生、結婚、死亡に関する正確なデータを集めました。彼らは生命年金、トントネ、保険などに関連する公式と表を提供しました。しかし、この短い概説では、これらの有用な著作を指摘し、元々のアイデアに従って論述することだけが可能です。

この中で、Daniel Bernoulliが与えた数学的および倫理的な希望に対する革新的な原理は特に特筆すべきです。また、彼が1763年の哲学雑誌の出版物において、過去の経験から示唆された可能性が特定の限界内に含まれる確率を求めることで、予防接種への微分積分学の巧妙な応用も特筆すべきです。

これらの原始的なアイデアの中で特に重要なのは、観察された事象から導き出される事象の可能性を直接考慮することです。Jacques BernoulliとMoivreはこれらの可能性を既知として扱い、将来的な経験の結果がそれらにますます近づく確率を求めました。一方で、1763年の哲学雑誌の出版物においてBayesは直接、過去の経験から示唆された可能性が特定の限界内にある確率を求めました。彼はこの問題を非常に巧妙な方法で解決しましたが、少し複雑でした。

この主題は観察された事象から導き出される原因と将来の事象に関する確率理論に関連しています。数年後、私はこれらの原理について説明し、等しい確率が仮定されている場合に存在する可能性のある不均一性の影響について注意を払いました。

これらの不均一性がどの単純な事象を有利にするかは不明ですが、この知識そのものがしばしば複合的な事象の確率を増加させることがあります。

一般化された解析と確率に関する問題を広げた際、私は部分的な有限差分の計算法に至りました。この方法はラグランジュによって非常に単純な手法で扱われており、その優雅な応用例は彼がこのような問題において使用しています。生成関数の理論については、私がほぼ同じ時期に発表したものはこれらの主題を含んでおり、最も難しい確率に関する質問にも最大限の一般性を持って適用できます。この理論は、多くの項と因子からなる関数の値を非常に収束する近似で再定義し、円周率に対する半径の平方根がこれらの値に頻繁に含まれることを示すことで、無数の他の超越関数も導入されることを示しています。

選挙や議決集会における証拠、投票、および判断も確率計算の対象となりました。多くの情熱、多様な利益、そして状況がこれらの主題に関連する質問を複雑にし、しばしばそれらはほとんど解けないものとなります。しかし、そのような問題と非常に類似した非常に単純な問題の解決策は、しばしば重要な難問に対する大きな光を投げかけ、計算法の確実性が最も説得力のある推論よりも常に優れていることを示します。

確率計算における最も興味深い応用の一つは...

確率論は観測結果の中から選択すべき平均値を扱います。多くの幾何学者がこの主題に取り組み、ラグランジュはトゥリーニのメモリアルに、観測誤差の法則が分かっている場合にこれらの平均値を決定する美しい方法を発表しました。私は同じ目的のために、一種の特殊な装置に基づく方法を与えました。この方法は他の分析上の問題でも有利に使用できるもので、これは長い計算全体の過程において無限に拡張でき、最終的な結果となる関数が性質上制限されるべきであるという自然の性質によって各項が受けるべき修正を示唆します。

既に述べたように、それぞれの観測は条件方程式として最初の次数の式を与えます。これらの式は常に第一項にすべての項を配置し、第二項をゼロにする方法で扱うことができます。これらの式の使用は、我々の天文表の大きな精度の一つの主な原因であり、多くの優れた観測がこれによって一致して要素を決定するために利用されました。

要素がただ一つしかない場合、コッツは条件方程式が未知の要素の係数が各々で正であるように準備されるべきであると指示しました。そしてこれらのすべての方程式を加えることで最終的な方程式を作り出し、その結果からこの要素の値を求めることです。コッツの規則は全ての計算家によって従われましたが、彼がいくつかの要素を決定することができなかったため、条件方程式をどのように組み合わせるべきかには固定されたルールはありませんでした。

各要素について、その決定に最も適した観測を選んだことで必要な最終方程式を得る方法は省略しますが、それぞれの要素に対して最適な観測を選択しました。 Legendre と Gauss は、これらの摸索を避けるために、条件方程式の最初の項の平方和を加え、未知の各要素についてこれを最小にするようにしました。これにより直接、要素数分の最終方程式が得られます。しかし、これらの方程式によって決定された値が他の方法で得られるすべての値よりも優れていると評価される資格があるでしょうか？この質問は、確率の計算だけが答えられました。私はその主題に適用し、精密な分析により前記の方法を含む規則的な過程で望む要素を得る利点を加えたルールを得ました。

しかしながら、それらの結果に対する我々の知識は不完全であり、それらが受ける誤差の法則が未知である限りです。これらの誤差が特定の限界内に含まれる確率を指定する必要があります。これは私が「結果の重さ」と呼んだものです。解析はこの目的のために一般的で簡便な公式を導き出します。私はこの分析を地図測量観測の結果に適用しました。一般的な問題は、一つまたは複数の線形関数の値や、非常に多くの観測の誤差が任意の限界内にある確率を決定することです。

193

観測の誤差の可能性の法則は、これらの確率を表す式に定数を導入し、その値はこの法則の知識が必要とされるが、これはほとんど常に未知である。幸いにも、この定数は観測から決定できる。

天文学的要素の調査においては、各観測値と計算された値との差の平方の和によって与えられる。同様に可能な誤差は、この和の平方根に比例するため、これらの平方を比較することで、同じ星の異なる表の相対的な正確さを評価できる。地図作成作業においては、三角形ごとの三つの角度の観測された合計の誤差の平方が代わりに使用される。これらの誤差の平方を比較することにより、角度を測定した際に使用された機器の相対的な精度を判断できる。この比較から、地図作成において繰り返し円分が置き換えられた機器よりも優れていることが明らかになる。

観測においては、多くの誤差の源が存在することが多く、その一つとして、星の位置を経線望遠鏡と円盤で定める際に生じる誤差がある。これらの誤差はそれぞれ異なる法則に従うものであり、それらの法則を同じものとする仮定は適切ではない。したがって、これらの観測値から導き出される要素は、その誤差によって影響を受けている。

条件方程式は、これらの要素を求めることを目的として作られるが、各機器の誤差と様々な係数を含んでいる。これらの条件方程式に適切な因子を乗じて、最終的な方程式を得るためには、どの程度の数の要素が必要であるかに関わらず、最も有利な因子システムを選ぶことが重要である。私は用いた解析は、どの程度の誤差源がある場合でも、同じ誤差が他の任意のシステムよりもより低い確率で生じるような因子システムを容易に導き出すものである。

同様の解析は、これらの結果における誤差の確率法則も決定する。これらの公式には、誤差の源の数だけ未知の定数が含まれており、それらはこれらの誤差の確率法則に依存している。

一つの誤差源の場合を考慮すると、その定数は、各条件方程式の残差の平方和を計算し、要素の値が求められたときにそれを形成することで決定できることが示されている。同様の過程は、どの程度の数の定数であっても、観測結果に対する確率解析の適用を完成させるものである。

私はここで重要な注釈を行うべきである。観測が少ない場合、それらの不確実性は、私がさきほど述べた定数の値に関して、少々の不確定性を残す。そして、これが分析によって決定された確率にわずかな影響を与える。しかし、それはしばしば十分であり、結果の誤差が狭い範囲内にある確率がほぼ1に近づいているかどうかを知るだけでよい。もしもそれがそうではない場合でも、観測を何倍にするべきかを知るだけでよい。そしてそのような確率を得ることで、合理的な疑念が残らないようにするためである。

確率論における計算。195

結果の正確性について考慮する必要がある。確率の解析公式はこの要件を完全に満たしており、その点では観測全体に基づく科学の必要な補完と見なすことができる。また、自然科学や道徳科学における多くの問題の解決に不可欠である。

現象の規則的な原因はしばしば未知であり、あるいは計算に適さないほど複雑である。さらに、その作用は偶然的で不規則な原因によってしばしば妨げられるが、これらのすべての原因により生じた出来事には常にその影響が残され、長期的な観測を通じて特定される修正をもたらす。確率の解析はこれらの修正を開示し、それらの原因の確率を定め、この確率を継続的に増加させる方法を示唆する。

したがって、気象の不規則な原因の中でも、太陽熱の日と夜、そして冬と夏との周期的な変化は大きな液体の圧力とそれに応答するバロメータの高さに一日と一年の振動を与え、多くの観測によってその前者が少なくとも同じ確率で我々が絶対的な事実として考慮しているものと同じ確率を示すことが明らかになった。

また歴史的出来事の系列は、社会がいかなる形でも混乱する感情や多様な利益の中においても、倫理学の大原則の恒常的な作用を示唆している。確率論という科学は、偶然性のあるゲームから始まったにもかかわらず、人間知の最も重要な主題にまで昇華されたことに驚きを感じる価値がある。

私はこれらの方法を全て、『解析的概率論』という著作に収録しました。そこで私は、確率計算の原理と分析を最も一般的な形で説明し、またその計算が提示する最も興味深い難問の解法も提案しました。

この論文から明らかになるように、確率理論は結局はただの常識を数値的に表現したものです。それは正確な思考によって本能的に感じられるものを正確に評価させます。そして、人間の判断における恣意性を排除し、常に最も有利な選択が決定できるようになります。これにより、人間の知識と弱さを補完する最良の方法となります。

この理論による分析的手法、その基礎となる原理の真実性、問題解決に必要な細密で精巧な論理、それに基づく公共利益の確立、そして自然哲学や倫理学における最も重要な問題への適用によって得られた拡張性を考慮すると、また、計算不能な事項においても最良の指針を与えることから、我々はこの科学がより深く且つ有用であることを理解するでしょう。

短題目カタログ

出版物の

ジョン・ウィリーズ\& Sons 社による

ニューヨーク版

ロージッド: シャープマン\&ホール リミテッド

分野別に並べ替えられています。

ご要望により説明書を送らせていただきます。

星印が付いている本は卸価格のみで販売されます。

すべての書籍は布製のカバーで綴じられており、特にその旨が記されていない限り。

農業。

アームズビーの家畜飼育マニュアル（12モ）

ダウニングのアメリカの果樹と果樹樹種（8オクターブ） 5.00ドル

グロテンフェルの現代乳業実践論理（ウォール著）（12モ） 2.00ドル

ケムのランドスケープガーデニング（12モ） 2.50ドル

メイナードの家宅装飾に適用されるランドスケープガーデニング（12モ） 1.50ドル

ストックビルズの岩と土壌（8オクターブ） 2.50ドル

ウォールの農民と乳牛飼育者向けハンドブック（16モ） 1.50ドル

建築。

バロウズの建物における蒸気暖房（12モ） 2.50ドル

ベルグのアメリカ鉄道の建築と構造（4オクターブ） 5.00ドル

ビルキムのアメリカ劇場の計画と建設（大開本） 3.00ドル

"         建築用鋼材と鉄骨（8オクターブ） 3.50ドル

"         鋸接組合せ梁の建築における適用（8オクターブ） 2.00ドル

"         高層事務所ビルの計画と建設（8オクターブ） 3.50ドル

"         建築における骨組み構造（8オクターブ） 3.00ドル

ブラギーズの現代アメリカ学校建築（8オクターブ） 4.00ドル

カーペンターの建物の暖房と換気（8オクターブ） 3.00ドル

フレイタージュの建築工学（8オクターブ） 3.50ドル

"         鋼骨建物の防火処理（8オクターブ） 2.50ドル

ゲラハードの衛生的な住宅検査ガイド（16モ） 1.00ドル

劇場火災とパニック（12モ） 1.50ドル

ハティフィールドのアメリカの家 Carpenter（8オクターブ） 5.00ドル

ホリリーの木工と組み立て手帳（18モ） 75セント

キッドラーの建築家と建設者の Pocket-book（16モ、モラコ） 4.00ドル

メリルの建築と装飾用石材（8オクターブ） 5.00ドル

モンクトンの階段建築（4オクターブ） 4.00ドル

パットンの基礎実用論（8オクターブ） 5.00ドル

シビルとビッジンの現代石材切断と石工（8オクターブ） 1.50ドル

ウェイトのエンジニアリングと建築法的権学（8オクターブ） 6.00ドル

羊、6.60ドル

建設前の操作に関する法律 建築およびエンジニアリングにおける構造計画に先立ち（8オクターブ） 6.00ドル

羊、5.60ドル

"         雇用契約の法理（8オクターブ） 3.00ドル

ウッドバリーの工場における火災保護（8オクターブ） 2.50ドル

ウォースターとアtkinsonの小規模病院、設立と維持、および病院建築に関する提案、附図を含む 小規模病院（12モ） 1.25ドル

世界のコロンビア展1893年 大4度、100

軍隊と海軍。

バーナドゥの無煙薬、硝酸セルロースおよびセルロース分子理論 12度、250

ブリフの砲兵学と射撃術教科書 8度、600

チェイスの螺旋推進器と航海推進 8度、300

クレホアとスカイアールの極性写真時計 8度、300

クロンヒットの非委任官向け射撃術 24度、摩砂皮装丁、200

デイビスの法の原理 8度、250

「      米国軍事法の論理」 8度、750 羊皮紙装丁、750

デブロックの騎兵哨戒勤務。 （カール） 24度、摩砂皮装丁、200

ダイツの兵士初級救急手引き 16度、摩砂皮装丁、125

「      現代フランス砲兵」 半摩砂皮装丁、1500

ドランの船舶抵抗と推進 8度、500

「      明るい砲兵ハンドブック」 12度、300

ツィスラーの現代高爆薬 8度、400

「      領地要塞工事教科書」 小8度、200

ホフの基礎海軍戦術 8度、150

インガルズの直接射撃問題ハンドブック 8度、400

「      精算表」 8度、150

ライオンズの電磁現象論文 8度、600

マーハンの永久要塞 （メルクール） 。 8度、半摩砂皮装丁、750

軍事裁判所用手引き 16度、摩砂皮装丁、150

「      防御地攻撃」 12度、200

「      武器の原理」 8度、400

メタカルフの製造費 — 公私工作場の管理 8度、500

「      砲兵学と射撃術」 12度、500

ムーリーの歩兵訓練規程 18度、紙装丁、10

「      実用航海測量」 8度、250

ポールズの軍事官僚試験 12度、400

シャープの軍隊戦時食糧管理術 18度、摩砂皮装丁、150

ウォルケの爆薬講義 8度、400

「      砲撃作戦と軍事採掘」 8度、200

ウィンスロップの軍事法要綱 12度、250

ウッドヒールの軍事衛生ノート 16度、摩砂皮装丁、150

ヤングの単純航海原理 16度、摩砂皮装丁、100

第二版、拡大改訂版 16モ, モルジョア, 200

試験分析の実用的手引.

フレッチャーの定量的試験分析におけるブロウパイプの使用法 12モ, モルジョア, 150

フルマンの実用的な試験分析マニュアル 8オ, 300

ミラーの試験分析マニュアル 12モ, 100

オドリスコールの金鉱石処理に関するノート 8オ, 200

リッケッツとミラーの試験分析に関するノート 8オ, 300

ウィルソンの cyanide プロセス 12モ, 150

"       氯化プロセス    12モ, 150

天文.

クレッグの方位角 4オ, 350

ドゥリットルの実用的な天文学講義 8オ, 400

ゴアの測地学入門 8オ, 250

ヘイフォードの測地天文学教科書 8オ, 300

メリマンの精密測量と測地学の要素 .... 8オ, 250

* ミッチーとハーロールの実用的な天文 8オ, 300

* ウォートスの理論的および説明的な天文学の要素. 12モ, 200

植物学.

バドウィンの新英格兰 Orchids 小 8オ, 150

ダベンポートの統計的手法、生物学的変異への特別な参照 16モ, モルジョア, 125

スローンとベネットの構造的および生理学的植物学. 16モ, 225

ウェスタマイヤーの一般植物学概論. (シュナイダー) 8オ, 200

化学.

アドリアンの実験室計算と特定重率表,

12モ, 125

アレンの鉄分析用テーブル 8オ, 300

アールノルドの化学概論. (マンドル) (準備中)

アウステンの化学学生用ノート 12モ, 150

ベルナドゥの無煙粉. — ニトロセルロース、およびセルロース分子の理論 12モ, 250

ボルトンの定量的分析 8オ, 150

ブラッシュとペンフィールドの決定的鉱物学マニュアル..8オ, 400

クラーセンの電解による定量的化学分析. (ハerrick—ボルトウッド) 8オ, 800

コーンの指示剤と試験紙 12オ, 200

クラフトの質的化学分析の短いコース. (シュアーファー) 12モ, 200

ドレッセルの化学反応. (メリル) 12モ, 125

アイスラーの現代高爆薬 8オ, 400

エフロンの酵素とその応用. (プリーストレー) (準備中)

エルマンの化学製剤への入門. (ダンラップ)

12mo, 1 25 フレッチャーの量的分析入門：ブロウパイプによる。12mo、モラコ、1 50

フレゼニウスの定量化学分析ハンドブック。（ウェルズ）

8vo, 5 00 定量化学分析の指導システム。（アレン） 8vo, 6 00

フルテスの水と公衆衛生。12mo, 1 50

フルマンの実用的金属分析ハンドブック。8vo, 3 00

ギルのエンジニアリング向けガスと燃料分析。12mo, 1 25

グローテンフェルトの現代乳業実践原理。（ウォール） . 12mo, 2 00 ハムマルストの生理化学教科書。（マンデル）

8vo, 4 00

ヘルムの数学的化学原理。（モーガン）. 12mo, 1 50 ホーラムネンの無機化学教科書。 （コッパー）

（準備中）

ホプキンスの油化学会ハンドブック。8vo, 3 00

キープの鉄板。 （準備中）

ラッドの量的化学分析ハンドブック。12mo, 1 00

ランダウエルのスペクトル分析。（ティンゲー） 8vo, 3 06

ラッサール・コーンの実用尿分析。（ロレンツ） （準備中）

ルーベスの有機化合物の電解と電合成。 （ロレンツ）. 12mo, 1 00

マンデルの生物化学実験室ハンドブック。12mo, 1 50

メイソンの水供給。（主に衛生的観点から検討）。8vo, 5 00

" 水質検査。（化学的および細菌学的）. 12mo, 1 25

マイヤーの炭化物中の基の決定。 （ティンゲー）. 12mo, 1 00

ミラーの分析ハンドブック。12mo, 1 00

ミッテルの化学入門教科書。12mo, 1 50

モーガンの溶液理論の概要とその結果。. . 12mo, 1 00

" 物理化学の要素. 12mo, 2 00

ニコルズの水供給。（主に化学的および衛生的観点から、1883年）。8vo, 2 50

オブライネの化学分析実験ガイド。8vo, 2 00

オドリスコールの金鉱石処理ノート。8vo, 2 00

オストとコルベックの化学技術教科書。（ロレンツ—ボザート）. （準備中）

* ペンフィールドの「決定的鉱物学ノートと鉱物試験記録」8vo、紙装、0 50

ピナーの「有機化学入門」（オーステン）. . .12mo, 1 50

プールの「燃料の calorific 力」8vo, 3 00

* リーゼイグの「点染色ガイド」8vo, 25 00

リチャーズとウォルマンの「衛生学的視点からの空気、水、食料」8vo, 2 00

リチャーズの「 sanitary 科学によって変更された生活費」12mo, 1 00

食品のコスト、栄養に関する研究12mo, 1 00

リッケッツとルーサーの「無機化学ノート（第1部。非金属元素）」. .8vo、モロコシ装、75

リッケッツとミラーの「採鉱ガイド」8vo, 3 00

ライデアルの「 Sewage および Sewage の細菌的浄化」.8vo, 3 50

ルディマンの「処方箋における不適合性」8vo, 2 00

シュンプフの「体積分析教科書」12mo, 2 50

スペンサーの「 beet-sugar 製糖所の化学家用ハンドブック」.

16mo、モロコシ装、3 00

製糖業者およびその化学家のためのハンドブック. .16mo、モロコシ装、2 00

ストックブラッドの「岩石と土壌」8vo, 2 50

* ティルマンの「熱の初等教科書」8vo, 1 50

"          説明的総合化学8vo, 3 00

ターンエアールとルーサーの「公共給水」8vo, 5 00

バン・デ・ベンターの「初学者用物理化学」（ボルトウッド）.

12mo, 1 50

ウォルケの「爆薬に関する講義」8vo, 4 00

ウェルズの「定量的化学分析実験ガイド」

..8vo, 1 50

"      工程学生用初学者の定量的無機化学分析短コース12mo, 1 50

ホプライプの「飲料水の顕微鏡」8vo, 3 50

ウィーチマンの「糖類分析」小8vo, 2 50

理論化学の講義ノート. .  .12mo, 3 00

ウィルソンの「 cyanide 工法」12mo, 1 50

"        氯化処理12mo, 1 50

ウッリングの「無機薬学および医学化学入門講座」12mo, 2 00

土木工学。

橋と屋根。水力学。工程材料。鉄道工学。

ベイカーの「エンジニア用測量器具」12mo, 3 00

ビックスビーの「グラフィカル計算表」紙装、19$p$x24$p$インチ. 25

ダVISの標高とスタディアテーブル 8vo、100円

フォウェルズの下水道。 （設計と維持管理） 8vo、300円

フライタッグの建築工学 8vo、350円

グッドヒューの都市改良 12mo、175円

グッドリッチの都市廃棄物の経済的処分 8vo、350円

ゴアの測地学入門 8vo、250円

ホイフードの測地天文学教科書 8vo、300円

HOWEの土留壁 12mo、125円

ジョンソンの測量理論と実践 小8vo、400円

"        スタディアと土工計算表 8vo、125円

キエルステッドの下水道処理 12mo、125円

マハンの土木工学論文。 （1873年） （木版印刷） . 8vo、500円

* マハンの描画幾何学 8vo、150円

メリマンの精密測量と測地学入門 8vo、250円

メリマンおよびブロックスの測量士ハンドブック 16mo、綴じ革、200円-

メリマンの衛生工学入門 8vo、200円

ヌジェントの平面測量。 （準備中）

オーガデンの下水道設計 12mo、200円

パットンの土木工学論文 8vo、半革装、750円

リードの地図とスケッチ法 4to、500円-

ライルドの下水道と下水道の細菌処理 8vo、350円

シービートとビッジンの現代石切りと石材工学 8vo、150円

スミスの地図測量手引 （マクリーン版） . 8vo、250円

トゥー・ツインの土木技師ポケットブック 16mo、モロコシ革装、500円

ウェイトの土木と建築法的権学 8vo、600円

羊皮紙、650円

"      工程および建築における作業法に関する法律 8vo、500円

羊皮紙、550円

"      契約法 8vo、300円

ウォーレンの立体幾何学 — 石切りに関する問題 8vo、250円

ウェブの工程機器の使用と調整に関する問題 16mo、モロコシ革装、125円

* ワーラーの土木工学入門講座 8vo、400円

ウィルソンの地図測量 8vo、350円

橋と屋根.

ボイルの鉄道道路橋の実用論文 8vo、200円

* ボイルのセント・メアズ川橋 4to、紙装、500円

バールの橋と屋根桁の荷重に関する課程 16mo、モロコシ革装、未定価

弧状の骨組みと懸索橋 8voie, 3 50

ド・ボイズの枠構造における応力 4to, 10 00

フォスターの木製桁架橋論文 4to,

フオラーのピール用仮堤工法 8voie,

グリーンの屋根梁 8voie,

"          橋梁梁 8voie,

"          木材、鉄、石における拱門 8voie,

ホーの拱門論文 8voie,

6

ジョンソン、ブライアンおよびトゥメアールの現代の枠構造設計に関する理論と実践 4to, 10 00

メリマンおよびジャコビの屋根と橋梁教科書：

第I部.—単純梁架における応力 8voie, 2 50

第II部.—図解静力学 8voie, 2 00

第III部.—橋梁設計。第四版（準備中） 8voie, 2 60

第IV部.—高度な構造物 8voie, 2 50

モリスムのメミノローグ橋 4to, 10 00*

ワデルのDe Pontibus、橋梁技術者のポケットブック 16mo, モレット, 3 00

鋼構造橋仕様 12mo, 1 25

ウッドの橋と屋根の理論に関する論文 8voie, 2 00

ライトの引き込み梁設計：

第I部.—板桁引き込み梁 8voie, 2 50

第II部.—接合栓結梁および長距離引き込み梁。8voie, 2 50

両部を一冊にまとめたもの 8voie, 3 50

水力学。

バジンの孔から流出する液体管の収縮に関する実験（トランスライティング） 8voie,

ボヴィーの水力学論文 8voie,

チャーチのエンジニアリング機械学 8voie,

コフィンの水力学問題の図解的解決法 16mo, モレット,

フォルウェルの給水工学 8voie,

フレズエルの水力 8voie,

フルテスの水と公衆衛生 12mo,

"        水処理施設 12mo,

ガングリエールおよびクッターの河川および他の渠道における水の均一な流れの一般公式（ヘリングおよびトラウトウィン） 8voie, 4 00

ハゼンの公衆給水用の水処理 8voie, 3 00

ハズルースの水道塔と貯水槽 8voie, 2 50

ヘルシュールの大口径、接合栓結金属管の荷重能力に関する実験 8voie, 2 00

マッソンの水道。衛生的な観点から主に検討されたもの。8voie, 5.00'

メリマンの hydraulic 工学。8voie, 4.00

ミッチェの解析力学の要素。8voie, 4.00

シュライヤーの灌漑、水力発電、および家庭用の水道用貯水池。大型 8voie, 5.00

ターンエアールとルーサーの公共用水道。8voie, 5.00

ヴェガンのダムの設計と建設。4to, 5.00

"ニューヨーク市の水道、1658年から1895年まで"。4to, 10.00

ウェイスバッハの hydraulic 工学および hydraulic モーター。（ドゥボイ）。8voie, 5.00

ウィルソンの灌漑工学マニュアル。小型 8voie, 4.00

ウォルフの風車。主な動力源として。8voie, 3.00

ウッドのタービン。8voie, 2.50

"解析力学の要素"。8voie, 3.00

エンジニアリング材料。

ベイカーの石造り構築に関する論説。8voie, 5.00

ブラックのアメリカ合衆国公務員用工作所。オブラング 4to, 5.00

ボーヴィの材料の強度と構造理論。8voie, 7.50

バールの材料の弾性と抵抗。8voie, 5.00

ブライアーンの道路建設。8voie, 5.00

"建設に使用された材料および作業の検査"。16mo, 3.00

チャーチのエンジニアリング機械学。8voie, 6.00

ドゥボイのエンジニアリング機械学。第1巻。小型 4to, 10.00

ジョンソンの建設材料。大型 8voie, 6.00

キープの鉄製品。（準備中）

ランザの応用力学。8voie, 7.50

マーテンスの材料試験ハンドブック。（ヘニング）。2巻、8voie, 7.50

メリルの建築および装飾用石。8voie, 5.00

メリマンの材料力学テキストブック。8voie, 4.00

メリマンの材料強度。12ino, 1.00

メタイフールの鋼。鋼を使用するためのマニュアル。12mo, 2.00

パットンの基礎に関する実用論説。8voie, 5.00

ロックウェルのフランスの道路と舗装。12mo, 1.25

スミスのワイヤー：その使用と製造。小型 4to, 3.00

スパルディングのハイドロセメント。12mo, 2.00

"道路と舗装に関するテキストブック"。12mo, 2.00

サーストンのエンジニアリング材料。3部、8voie, 8.00

第1部。—エンジニアリングと金属工学における非金属材料 8vo, 200円

第2部。—鉄鋼 8vo, 350円

第3部。—黄銅、青铜及其他合金及其成分的专论 8vo, 250円

スチュアートの建設材料教科書 8vo, 500円

ティルソンの街道舗装と舗装材料 8vo, 400円

ワデルの橋梁工程指南 (袖珍本) 16mo, モロコシ製, 300円

鋼桥规范 12mo, 125円

ウッドの材料的抵抗论及其附录木材防腐 8vo, 200円

"        分析力学元素 8vo, 300円

鉄路工程。

ベルグのアメリカ铁路建筑与结构. .4to, 500円

ブロックスの街道电车线路手册. .16mo, モロコシ製, 150円

バッツの土木工程师手册 16mo, モロコシ製, 250円

クランダルの过渡曲线 16mo, モロコシ製, 150円

铁路及其他土方工程表 8vo, 150円

ドーソンの电车与有轨电车. 小4to, 半モロコシ製, 1250円 " 工程"及电车牵引袖珍手册.

16mo, モロコシ製, 400円

ドリジーズ的宾夕法尼亚铁路史: (1879年). 纸本, 500円 * ドリンカー的隧道工程、炸药化合物及凿岩机.

4to, 半モロ科西製, 2500円

ホワイト的过渡曲线手册 16mo, モロコシ製, 150円

ハドソン的计算开挖与填土立方体积表格 8vo, 100円

ナガル的铁路工程师手册 16mo, モロ科西製, 300円

フィルブリック的工程师袖珍手册 16mo, モロコ西製, 300円

プラット与奥登的街道电车路基 8vo, 200円

シールズ的现场工程 16mo, モロ科西製, 300円

铁路螺旋线 16mo, モロコ西製, 150円

テイラーの棱柱体公式及土方计算 8vo, 150円

* イェートウィン的利用图表计算开挖与填土立方体积的方法 8vo, 200円

* " 现场铺设圆曲线的实践方法"

鉄道に関する図書：

鉄道12番目、モロコシ装丁、2ドル50セント

断面図　シート用紙、25セント

ウェブの鉄道建設8番目、4ドル00セント

ウィンダムの鉄道路線経済理論。小8番目、5ドル00セント

描画：

バレーズの機械運動学8番目、2ドル50セント

バートレットの機械図面作成8番目、3ドル00セント

ダーリーの機械の運動の基礎教科書。（準備中）

ヒルの影と視覚、および視点に関するテキストブック。8番目、2ドル00セント

ジョーンズの機械設計：

第1部 — 機械の運動学8番目、1ドル50セント

第2部 — 部品の形状、強度と比例8番目、3ドル00セント

マッコードの説明幾何学要素8番目、3ドル00セント

「          力学；または実用的な機械装置8番目、5ドル00セント

「          機械図面作成4番目、4ドル00セント

速度図8番目、1ドル50セント

マッハーンの説明幾何学と石切り8番目、1ドル50セント

マッハーンの工業図面作成。（トーマス）8番目、3ドル50セント

リードの地図上の描画と素描4番目、5ドル00セント

レidの機械図面作成コース8番目、2ドル00セント

「      機械図面作成と基礎的な機械設計テキストブック8番目、3ドル00セント

ロビンソンの機構原理8番目、3ドル00セント

スミスの地図上の描画マニュアル。（マクリーン）8番目、2ドル50セント

ウォーレンの平面と立体の自由手描き幾何学要素12番目、1ドル00セント

ドラッフィング器具と操作12番目、1ドル25セント

基本的な投影図面作成マニュアル。. 12番目、1ドル50セント

「        形状と影の線形視点に関する基礎問題12番目、1ドル00セント

「        平面幾何学の基礎問題12番目、1ドル25セント

「        基礎幾何学12番目、75セント

「        説明幾何学、影と視点の要素8番目、3ドル50セント

「        影と視点に関する一般問題8番目、3ドル00セント

「        機械構築と描画の要素8番目、7ドル50セント

「        説明幾何学の問題、定理と例題8番目、2ドル50セント

ウェイスバッハの運動学と伝達力。（ヘルマンとクライナー）8番目、5ドル00セント

ウェルプリーの彫刻芸術実用的な指導12番目、2ドル00セント

「ウィルソンの地図測量術」8voie, 3ドル50セント

「ウィルソンの自由手描き視覚法」8voie, 2ドル50セント

「ウルフの立体幾何学入門」大型8voie, 3ドル

電気と物理学。

アンドリューとブラケットの物理学教科書。（マージ）
小8voie, 3ドル

アンソニーの電測理論講義ノート。12モー, 1ドル

ベニオンの電気史。8voie, 3ドル

ベニオンのボルタ電池。8voie, 3ドル

クラーセンの電解による定量化学分析。（ヘリックとボルトウッド）
8voie, 3ドル

クリホアとスキュアの極光写真時計。8voie, 3ドル

ダウソンの電鉄および路面電車。小4度, 半綿, 12ドル50セント
ダウソンの「エンジニアリング」および電気運転ポケットブック。
16モー, 綿, 4ドル

スラザーズの動力計と功率測定。12モー, 3ドル

ギルバートの『マグネタ』。（モットレイ）
8voie, 2ドル50セント

ホーマンの精密測量。8voie, 2ドル

「望遠鏡鏡スケール法による調整と検査」大型8voie, 75セント

ランダウエルのスペクトル分析。（ティンゲ）
8voie, 3ドル

レ・シャトリエの高温測定。（ボドワード—バージュス）
12モー, 3ドル

ロブの有機化合物の電解と電合成。（ローレンツ）
12モー, 1ドル

ライオンズの電磁現象論文。8voie, 6ドル

ミッキーの波動運動の要素。（音と光）8voie, 4ドル

ニアドエトの電池入門。（フィシャバック）
12モー, 250セント

パーシャルとホバートの電力発生機。小4度, 半綿, 10ドル
スチュワートの固定蒸気機関車。8voie, 2ドル50セント

ティールマンの熱入門。8voie, 1ドル50セント

トリーとピッチャーの実験物理学ハンドブック。小8voie, 2ドル

法。

デイビスの法律要素。8voie, 2ドル50セント

「軍事法」（合衆国）8voie, 7ドル
シープ, 7ドル50セント

軍事裁判所用マニュアル。16モー, 綿, 1ドル50セント
シープ, 6ドル50セント

「工学と建築における作業法」8voie, 5ドル
シープ, 5ドル50セント

製造業。

ベアモントの羊毛とウールストド生地の製造12判、1.50ドル

ベルナドゥの煙突なしの爆薬—硝酸セルロースとセルロース分子の理論12判、2.50ドル

ボーランドの鉄鋳工12判、布製、2.50ドル

「鉄鋳工」補遺12判、2.50ドル

「鋳造百科事典と鋳造用語辞典」12判、3.00ドル

ブラスターの現代高爆薬8判、4.00ドル

エフォントの酵素とその応用。（プレストン）（準備中）

フィッツジェラルドのボストン機械工18判、1.00ドル

フォードのボイラー製造—ボイラー作業者向け18判、1.00ドル

ホプキンスのオイル化学ハンドブック8判、3.00ドル

キープの鍛鉄。準備中。

メットカフの鋼。鋼使用者用マニュアル12判、2.00ドル

メットカフの製造コスト—そして工場の管理、公的と私的8判、5.00ドル

マイヤーズの現代蒸気機関車構築4判、10.00ドル

リージングの個体染色ガイド8判、25.00ドル

スミスの金属圧延8判、3.00ドル

「線：その使用と製造」小4判、3.00ドル

スパルディングの Hydraulic セメント12判、2.00ドル

スペンサーの Beet-糖分野の化学者向けハンドブック16判、レザーカバー、3.00ドル

「砂糖製造業およびその化学者のためのハンドブック」16判、レザーカバー、2.00ドル

スチュアートの蒸気 Boiler の手引8判、5.00ドル

ウォルケの爆薬講義8判、4.00ドル

ウェストのアメリカ鋳造実践12判、2.50ドル

「型枠工の教科書」12判、2.50ドル

ヴィークマンの砂糖分析小8判、2.50ドル

ウォルフの風車は主な動力源8判、3.00ドル

ウッドバリーの工場の火災保護8判、2.50ドル

数学。

ベイカーの楕円関数8判、1.50ドル

* ベスの微分積分学入門12判、4.00ドル

ブリッジズの平面解析幾何学入門12判、0.00ドル

チャプマンの方程式論基礎講座…12判、50セント

コムポットの対数計算マニュアル12判、1.50ドル

ダヴィスの代数学論理入門8判、50セント

ハルステッドの幾何学入門8判、75セント

「初等的 合成幾何学」8vo版、50円

ジョンソンの三桁対数表：ポケットサイズ、紙装、15円

100部で5円に

厚紙に貼り付けられた、8×10インチ、25円

10部で200円に

「積分微分法の基礎論」

小8vo版、150円

「カートシアン座標における曲線描画」

12mo版、100円

「常微分方程式と偏微分方程式論」

小8vo版、350円

「誤差理論と最小二乗法」

12mo版、150円

* 「理論力学」

12mo版、300円

ルドローとバス。三角法と対数及其他表の要素 8vo版、300円

「三角法」

表は別売。各々、200円に

メリマンとウッドワード。高等数学 8vo版、500円

メリマンの最小二乗法 8vo版、200円

ライスとジョンソンの微分積分学の初等論文 小8vo版、300円

微分積分学。2巻を1冊にまとめた小8vo版、250円

ウッドの座標幾何学要素 8vo版、200円

「三角法：解析的、平面的、球面的」

12mo版、100円

機械工学。

工程材料、蒸気エンジンとボイラー。

バーリントンの建物向け蒸気暖房 12mo版、250円

バーの機械運動論 8vo版、250円

バートレットの機械描画 8vo版、300円

ベネットの小技とレシピ 12mo版、200円

キャプターの実験工学 8vo版、600円

建物向け暖房と換気 8vo版、300円

クリークのガスエンジン論文 小8vo版、400円

クロムウェルの歯車論文 12mo版、150円

ベルトとギアに関する論文 12mo版、150円

ダーリリーの機械運動の初等教科書（準備中）

平坦性ダイナモメータと力の測定 12mo版、300円

ロープドライブ 12mo版、200円

フィルズのエンジニア用ガスと燃料分析 12mo版、125円

ホールの車両潤滑 12mo版、100円

ジョーンズの機械設計：

第1部 — 機械運動 8vo版、150円

第2部 — 部品の形状、強度と比例 8vo版、300円

Kentの機械工のポケットブック。…16モ、モロコシ製カバー、5.00ドル

Xerrの力と力伝送。 （準備中）

MacCordの運動学；または実用的な機構。 8オ、5.00ドル

機械図面。4オ、4.00ドル

"          速度図。8オ、1.50ドル

Mahanの産業図面。（Thompson）8オ、3.50ドル

Pooleの燃料の熱量的力。8オ、3.00ドル

Reidの機械図面入門。8オ、2.00ドル

"      機械図面と初等的な機械設計の教科書。     8オ、3.00ドル

Richardsの圧縮空気。12モ、1.50ドル

Robinsonの機構原理。8オ、3.00ドル

Smithの金属のプレス加工。8オ、3.00ドル

Thirston摩擦と機械および工場作業における無駄な仕事の論理。8オ、3.00ドル

"          生物がマシンとして、原動機としての動物とエネルギーの法則。12モ、1.00ドル

Warrenの機械構築と図面入門。…8オ、7.50ドル

Weisbachの運動学および伝送力。（Herrmann—Klein）8オ、5.00ドル

"          伝送機械と制御装置。（Herrmann—Klein）8オ、5.00ドル

"          水力と水力モーター。（Du Bois）8オ、5.00ドル

Wolffの風車が原動機として。8オ、3.00ドル

Woodのタービン。8オ、2.50ドル

エンジニアリング用材料。

Boveyの材料力学と構造理論。8オ、7.50ドル

Burrの材料の弾性と抵抗。8オ、5.00ドル

Churchのエンジニアリング機械学。8オ、6.00ドル

Johnsonの建設用材料。大型8オ、6.00ドル

Keepの鍛造鉄。（準備中）

Lanzaの応用力学。8オ、7.50ドル

Martensの材料試験ハンドブック。（Henning）8オ、7.50ドル

Merrimanの材料力学テキストブック。8オ、4.00ドル

材料力学。12モ、1.00ドル

Metcalfの鋼。鋼使用者用マニュアル。12モ、2.00ドル

Smithのワイヤー：その使用と製造。小4オ、3.00ドル

Thurstonのエンジニアリング材料。3巻、8オ、8.00ドル

第II部—鉄および鋼。8オ、3.50ドル

第III部—銅、青铜及其他合金及其成分的论著。8オ、2.50ドル

トーマスの建設材料の教科書。…8voie, 5ドル

ワーズの材料の抵抗論と木材の保存附録。…8voie, 2ドル

解析力学の要素。…8voie, 3ドル—

蒸気機関およびボイラー。

カノンの熱の動力に関する思考。…（トーマス）

12モ、1ドル50セント ダウソンの「エンジニアリング」および電車運転ポケットブック。

16モ、モロコシ装丁、4ドル

フォードのボイラー製造者向けボイラー作成。…18モ、1ドル

ヘメンウェイのインディケーター実践と蒸気機関の経済性。…12モ、2ドル

ホットンの力発生プラントの機械工学。…8voie、5ドル

「熱と熱エンジン」…8voie、5ドル

ケントの蒸気ボイラーの経済性。…8voie、4ドル

クネアスのインジェクターの実践と理論。…8voie、1ドル60セント

マッコードのスライドバルブ。…8voie、2ドル

マイヤーの現代ロコモティブ建設。…4to、10ドル

ピーボディの蒸気機関インディケーターの手引き。…12モ、1ドル50セント

「飽和水蒸気および他の蒸気の性質表」…8voie、1ドル

「蒸気機関およびその他の熱エンジンのthermodynamics」…8voie、5ドル

「蒸気機関のバルブギア」…8voie、2ドル50セント

ピーボディとミラー。蒸気ボイラー。…8voie、4ドル

プライヤーズのインディケーターとの二十年間。…大型8voie、2ドル50セント

ポプリンの可逆サイクルにおけるガスおよび飽和水蒸気のthermodynamics。（オスターベルグ）12モ、1ドル25セント

レーガンのロコモティブ機械とエンジニアリング。…12モ、2ドル

ロンテンゲのthermodynamics原理。…（ドゥボイ）8voie、5ドル

サイクルインのロコモティブエンジン運転および管理。…12モ、2ドル

スマートの技術工場実験所手引書。…12モ、2ドル50セント

スノウの蒸気ボイラー実践。…8voie、3ドル

スパングラーのバルブギア。…8voie、2ドル50セント

「thermodynamicsに関する補注」…12モ、1ドル

トーマスの便利な表。…8voie、1ドル50セント

蒸気機関の手引き2巻、8voie、10ドル

第1部—歴史、構造および理論。…8voie、6ドル

第2部—設計、建設および運用。…8voie、6ドル

トーマスのエンジンとボイラー試験ハンドブック、インディケーターおよびプロニーブレーキの使用。…8voie、5ドル

「静止する蒸気機関 8vo、2.50ドル

蒸気 boiler 爆炸：理論と実践 12mo、1.50ドル

蒸気 boiler 手冊：設計、構築、操作 8vo、5.00ドル

Weisbachの Heat, Steam, and 蒸汽机 8vo、5.00ドル（Du Bois）

Whithamの 蒸汽机设计 8vo、5.00ドル

Wilsonの 蒸气锅炉论：Flather 16mo、2.50ドル

Woodの 热力学，热动机和制冷装置 8vo、4.00ドル

機械学与机械工程。

Barrの 机器的运动学 8vo、2.50ドル

Boveyの 材料强度与结构理论 8vo、7.50ドル

Chordal — 函件摘录 12mo、2.00ドル

Churchの 工程机械学 8vo、6.00ドル

力学笔记与例题 8vo、2.00ドル

Comptonの 铸造术初学者 12mo、1.50ドル

ComptonとDe Groodtの 速度车床 12mo、1.50ドル

Cromwellの 齿轮理论 12mo、1.50ドル

带与滑轮论 12mo、1.50ドル

Danaの 初级力学教科书：用于学院和学校的机械工程 12mo、1.50ドル

Dingeyの 机器模型制作 12mo、2.00ドル

Dredgeの 世界哥伦比亚博览会运输展品的建造记录 1893年 4to，半皮装，5.00美元

Du Boisの 基础力学原理：

第I卷— 运动学 8vo、3.50美元

第II卷— 静力学 8vo、4.00美元

第III卷— 动力学 8vo、3.50美元

Du Boisの 工程机械学。 第I卷 小4to、10.00美元

Durleyの 机器运动学基础（准备中）

Fitzgeraldの 波士顿机工 16mo、1.00美元

Flatherの 动力测量：动力计 12mo、3.00美元

「绳索传动 12mo、2.00美元

Hallの 车辆润滑 12mo、1.00美元

Hollyの 锯片研磨艺术 18mo、76セント

* Johnsonの 理论力学 12mo、3.00ドル

Jonesの 机械设计：

第I部分— 机器的运动学 8vo、1.50ドル

第II部分— 部件的形式，强度和比例 8vo、3.00ドル

Kerrの 力与力传递（准备中）

Lanzaの 应用力学 8vo、7.50ドル

MacCordの 运动学；或实用机械学 8vo、5.00ドル

「速度図」八寸本、150円

メリマンの物質力学入門 八寸本、400円

* ミッチーの解析力学入門 八寸本、400円

レーガンの機関車装置と工学 12寸本、200円

レイドの機械描画入門 八寸本、200円

リチャーズの圧縮空気 12寸本、150円

ロビンソンの機構原理 八寸本、300円

サイレンズの機関車運転と管理 12寸本、200円

スミスの金属のプレス加工 八寸本、300円

サーストンの摩擦と機械・工場での失われた作業 八寸本、300円

「動物はマシンであり原動機であり、エネルギー論法の法則」12寸本、100円

ウォーレンの装置構築と描画入門 八寸本、760円

ワイスバッハの運動学と伝達力（ヘルマン—クレイン） 八寸本、600円

「伝送機械と制御器」（ヘンリ—クレイン） 八寸本、600円

ウッドの解析力学入門 八寸本、300円

「初等力学原理」12寸本、125円

* テルビンス 8寸本、250円

1893年の世界コロンビア博覧会 四寸本、100円

金属冶治学。

エグレストンの銀・金・水銀の冶金学 第一巻—銀 八寸本、750円

第二巻—金と水銀 八寸本、750円

キープの鍛造鉄（準備中）

エンハルトのヨーロッパでの採鉱実践 8寸本、150円

シャトリエの高温測定（ボドウアード—バージス） 12寸本、300円

メットカフの鋼。鋼使用者用マニュアル 12寸本、200円

サーストンの工学材料。三部作 八寸本、800円

第二部—鉄と鋼 八寸本、360円

第三部—銅合金その他の合金及其成分の論文 八寸本、250円

Minerology.

バリンジャーの商業価値のある鉱物の説明 横型、モロコシ装、250円

ボイドの南 virginia の資源 8寸本、400円

「南 virginia 地図」ポケットサイズ、200円

ブラッシュの決定的 mineralogical 手引（ペンフィールド） 八寸本、400円

チェスターの鉱物目録 8寸本、紙装、100円

布製、1 25

「鉱物名辞典」8vo、3 50

ダナの鉱物学体系 大型版、半革装、12 50

「ダナの新鉱物学体系」第一付録 大型版、1 00

「鉱物学教科書」8vo、4 00

「鉱物とその研究法」12mo、1 50

「アメリカ産鉱物の目録」大型版、1 00

「鉱物学およびペトログラフィーの手引」12mo、2 00

エグレストンの鉱物と同義語目録 8vo、2 50

フサックの岩形成鉱物の決定法 (スミス著) 小型版、2 00

ペンフィールドの決定的鉱物学のノートおよび鉱物試験記録 8vo、紙装、50

ローゼンブッシュの岩形成鉱物の顕微鏡的生理学 (イッディング著) 8vo、5 00

ティルマンの重要な鉱物と岩石の教科書 8vo、2 00

ウィリアムズの地質学的手引 8vo、3 00

採掘。

ベアードの坑内の換気 12mo、2 60

ボイドの南ウェストバージニアの資源 8vo、3 00

「南ウェストバージニアの地図」ポケット形式、2 00

ドリンカーのトンネリング、爆薬化合物、および岩石打眼機 (4to装)、半モラコシ装、25 00

ブラスターの現代高爆薬 8vo、4 00

オードリーヨークのアメリカ西海岸の石炭坑 12mo、2 50

イッelsingの採掘手引 8vo、4 00

コンハートのヨーロッパでの鉱物精錬実践 8vo、1 60

オドリスコールの金鉱石の処置に関するノート 8vo、2 00

サウェイアの坑内事故 8vo、7 00

ウォルケの爆薬講義 8vo、4 00

ウィルソンの cyanide 工法 12mo、1 50

ウィルソンの塩素化工法 12mo、1 60

ウィルソンの水力および placer 掘削 12mo、2 00

ウィルソンの実用的かつ理論的な坑内換気講義 12mo、1 25

衛生科学。

フォールウェルの下水道 (設計、建設及び維持) 8vo、3 00

「水供給工学」8vo、4 00

フルテスの水と公衆衛生 12mo、1 50

「水濾過施設」12mo、2 50

ゲルハルトの衛生的住宅検査ガイド 16mo、1 00

ゴードリッジによる都市廃棄物の経済的な処分. . .デミ8voie, 3ドル50セント

ヘズンによる公用水供給の濾過. . .Svoie, 3ドル

キエルステッドによる Sewage 处置. . .12モ, 1ドル25セント

メイソンによる水供給. (主に衛生的観点から考慮)8voie, 5ドル

"        水の検査. (化学的および細菌学的.) 12モ, 1ドル25セント

メリマンによる sanitary 工学の要素. . .Svoie, 2ドル

ニコルズによる水供給. (主に化学的および衛生的観点から考慮) (1883年). 8voie, 2ドル50セント

オーガンによる下水管設計. . .12モ, 2ドル

リチャーズによる食料のコスト. 食事研究. 12モ, 1ドル-

リチャーズとウォルマンによる空気、水、および衛生的観点から考慮した食料8voie,

リチャーズによる衛生科学によって修正された生活費のコスト. 12モ,

ライダップによる Sewage および細菌による Sewage の処理. . .8voie,

ターンエアとルーサーによる公用水供給. . .Svoie,

ホプリによる飲料水の顕微鏡検査. . .8voie,

ウッドヒールによる軍事衛生に関するメモ. 16モ,

雑誌.

バッカーによる深海測量. . .Svoie, 3ドル

エミンスンによるロッキー山脈冒険の国際地質学者会議の地質ガイドブック. 大きな Svoie,

フェールによる気象の通俗論文. . .Svoie,

ヘインズによるアメリカ鉄道管理. 12モ,

モットによる食料の構成、消化性、および栄養価. 指図されたチャート,

"      現在の音理論の誤謬. 16モ,

リケットスによるレンセラーポリテクニックインスティテュートの歴史, 1824-1894年. 小さな Svoie,

ロザンハムによる強調新約聖書. 大きな Svoie,

批判的強調新約聖書. 12モ,

ステールによる犬の病気に関する論文. . .Svoie,

トーテンによるメートロロジーにおける重要な問題. . .Svoie,

1893年の世界コロンビア博覧会.

ワースワースとアtkinsonによる小規模病院、設置と維持、および小規模病院建築に関する提案. 12モ, 1ドル25セント

ヘブライ語とカルデー語の教科書.

グリーンによるヘブライ語文法. . .Svoie, 3ドル

"       元則的なヘブライ語文法. 12モ, 1ドル25セント

"       ヘブライ語のクリストマティ. . .Svoie, 2ドル

ゲゼニウスの旧約聖書用ヘブライ語・アッシリヤ語大辞典

ゲゼニウスの旧約聖書用ヘブライ語・アッシリヤ語大辞典。トリジェレス訳。小4开本、半ynchronously 皮装、5.00ドル。

レターイスのヘブライ語旧約聖書。Svo（大8开）。2.25ドル。

カリフォルニア大学
南地域図書館施設

305 デ・ネーヴェ駅前通り - 17号駐車場 • 郵便番号951388

ロサンゼルス、カリフォルニア 90095-1388

この資料は借りた図書館から返却してください。


\backmatter

\chapter{訳者・数理校閲後記}

本書は、ピエール＝シモン・ラプラスによる不朽の名著『確率の哲学的試論』（1814年）の全18章を、Project Gutenberg（\#58881）の底本に基づき完全翻訳・数学校閲したものである。

原著者が序論で力説するように、「人間の知識のほとんどすべては確率的である」。ラプラスが200年前に構想した決定論的宇宙観と、不完全な情報のもとで推論を行うための「原因の確率（ベイズの逆確率）」、「大数の法則」、「母関数法」、「最小二乗法」の体系は、現代のデータサイエンス、機械学習、ベイズ統計学、金融工学、認知科学の礎石そのものである。

本全訳版が、確率論の哲学的源流と数理的精緻さを同時に味わうための貴重な道標となることを願ってやまない。

\begin{flushright}
2026年10月\\
教育情報分析研究会
\end{flushright}

\end{document}
